\section{Modular holography on the proved core: reconstruction, transposition, quantization, and metaplectic anomaly}
\label{sec:frontier-modular-holography-platonic}
\index{holography!modular proved core|textbf}
\index{modular homotopy theory!holography|textbf}
\index{metaplectic correction!holographic kernel|textbf}
\index{determinant anomaly!holographic sewing|textbf}

The proved modular Koszul core supplies the algebraic protected
package on the boundary curve. Throughout this section, modularity
is the image of the modular trace $\Tr_\cC$ and clutching on the open
factorisation category $\cC^{\mathrm{op}}$ on $(X,D,\tau)$, not a
closed-algebra property of $\cA$ itself
\textup{(}Convention~\textup{\ref{conv:master-primitive-datum}};
the level-$5$ scalar shadow is reconstructed from level~$4$ via
modular reconstruction with the trace + clutching hypothesis package
of Convention~\textup{\ref{conv:master-reconstruction-tower}}\textup{)}.
\begin{enumerate}[label=\textup{(\roman*)}]
\item the bar--cobar transform and its Verdier intertwining
 \textup{(}Theorem~\textup{\ref{thm:bar-cobar-isomorphism-main}},
 Theorem~\textup{\ref{thm:higher-genus-inversion}}\textup{)},
\item the genus-$g$ ambient complex
 $\mathbf C_g(\cA)=R\Gamma(\overline{\mathcal M}_g,\mathcal Z(\cA))$
 together with the Verdier involution $\sigma$
 \textup{(}Definition~\textup{\ref{def:complementarity-complexes}}\textup{)},
\item the complementary-Lagrangian decomposition
 $\mathbf C_g(\cA)\simeq \mathbf Q_g(\cA)\oplus \mathbf Q_g(\cA^!)$
 \textup{(}Lemma~\textup{\ref{lem:involution-splitting}},
 Theorem~\textup{\ref{thm:quantum-complementarity-main}},
 Proposition~\textup{\ref{prop:ptvv-lagrangian}},
 Theorem~\textup{\ref{thm:ambient-complementarity-fmp}}\textup{)},
\item and the scalar genus package
 $\mathrm{obs}_g(\cA)=\kappa(\cA)\lambda_g$
 on the proved uniform-weight lane
 \textup{(}Definition~\textup{\ref{def:scalar-lane}}\textup{)},
 or more generally whenever the full
 scalar package is explicitly assumed
 \textup{(}Theorem~\textup{\ref{thm:genus-universality}},
 Corollary~\textup{\ref{cor:scalar-saturation}}\textup{)}.
\end{enumerate}
These ingredients yield an algebraically complete
\emph{protected transform calculus} on the proved core. This is a
boundary theorem. A holomorphic-topological bulk, a categorical trace
model, or a physical path integral enters only through a specified
comparison functor to this protected calculus. The
model-independent level lives in the homotopy category of
factorization algebras or perfect complexes. The strict model level
uses finite-rank perfect dg representatives of
$\mathbf Q_g(\cA)$ and $\mathbf Q_g(\cA^!)$ on which the Verdier
pairing is strict. The scalar level records the cohomological
shadows: Poincar\'e polynomials, determinant lines, Euler
characteristics, scalar anomaly towers, and reciprocity laws. It does
not record the cross-channel tower unless the all-weight hypotheses
are separately imposed. The
metaplectic Gaussian formulas below use an even finite-rank dg model;
the graded version is obtained from the same calculation by the
Koszul sign rule.

The seven-slot holographic modular Koszul datum
$\cH(\cA) = (\cA, \cA^i, \cA^!, \cC, r(z), \Theta_\cA,
\nabla^{\mathrm{hol}})$
(Definition~\ref{def:holographic-modular-koszul-datum}) is now
subsumed by the nonuple \emph{completed modular Koszul datum}
$\Pi^{\mathrm{oc}}_X(\cA)$
(Definition~\ref{def:thqg-completed-platonic-datum}), which adds
the chiral derived centre
$\cZ^{\mathrm{der}}_{\mathrm{ch}}(\cA)$ (the chiral Hochschild
closed sector of the boundary algebra), the categorical annulus trace
$\mathrm{Tr}_\cA\simeq
\mathrm{HH}_{\bullet,\mathrm{cat}}(\mathcal M)$ (a categorical open
trace after a chosen module-category comparison), and the nonlinear
resonance tower
$\mathfrak{R}^{\mathrm{oc}}_\bullet(\cA)$ (quartic and higher
obstructions in $\cZ^2$).
The displayed symbols are typed slots, not object equalities:
chiral Hochschild, categorical Hochschild/THH, and physical bulk
factorization algebras are compared by maps, and only become
equivalent after the stated finite-type, dualizability, and completion
hypotheses are verified.
See \S\ref{sec:thqg-open-closed-realization} for the full
development.

\subsection{Protected dual transform and the holographic transform}
\label{subsec:frontier-protected-holographic-transform}

\begin{definition}[Protected holographic transform]
\label{def:frontier-protected-holographic-transform}
\index{holographic transform|textbf}
Let $X$ be a smooth projective curve and let $\cA$ be a modular Koszul chiral
algebra on~$X$.
Let
\[
\iota_X\colon \mathrm{ChAlg}^{\mathrm{str}}_X
\longrightarrow D^{\mathrm{co}}\!\bigl(\operatorname{Fact}(X)\bigr)
\]
denote the inclusion of a strict finite-type chiral algebra as its
coderived factorization object whenever such a strict representative
exists.
\begin{enumerate}[label=\textup{(\roman*)}]
\item \textbf{H-level.}
 The \emph{protected holographic transform} is
 \begin{equation}\label{eq:frontier-protected-holographic-transform}
 \mathbb H_X(\cA)
 :=
 \Omega_X\,\mathbb D_{\operatorname{Ran}}\,\barB_X(\cA)
 \in D^{\mathrm{co}}\!\bigl(\operatorname{Fact}(X)\bigr).
 \end{equation}
\item \textbf{H/M-level.}
 For every $g\ge 0$, with ambient complex
 $\mathbf C_g(\cA)=R\Gamma(\overline{\mathcal M}_g,\mathcal Z(\cA))$, define
 the \emph{holographic cofiber}
 \begin{equation}\label{eq:frontier-holographic-cofiber}
 \mathbf H^{\mathrm{hol}}_g(\cA)
 :=
 \cofib\!\bigl(\mathbf Q_g(\cA)\longrightarrow \mathbf C_g(\cA)\bigr).
 \end{equation}
\item \textbf{S-level.}
 The cohomology
 \begin{equation}\label{eq:frontier-holographic-shadow}
 H^*_{\mathrm{hol},g}(\cA):=H^*\!\bigl(\mathbf H^{\mathrm{hol}}_g(\cA)\bigr)
 \end{equation}
 is called the \emph{genus-$g$ protected holographic shadow}.
\end{enumerate}
\end{definition}

\begin{theorem}[Protected dual transform and anti-involution;
\ClaimStatusProvedHere]
\label{thm:frontier-protected-bulk-antiinvolution}
\index{dual transform!protected theorem|textbf}
\index{anti-involution!holographic transform|textbf}
\textup{[}Type signature:
$(\mathrm{quadrant}=\mathrm{Open}\cap\mathrm{Cat},
\mathrm{presentation}=\mathrm{factorisation},
\mathrm{level}=2\leftrightarrow 3,
\mathrm{hypothesis\ package}=\mathrm{modular\ Koszul\ on\ KSDual\ sublocus};
\mathrm{strictification\ of\ Verdier/Koszul\ dual};\mathrm{Mittag\text{--}Leffler}
+\mathrm{separatedness};\mathrm{all\ genera})$
\textup{(}Convention~\textup{\ref{conv:regime-tags}},
Convention~\textup{\ref{conv:hms-levels}};
KSDual: $\mathbb{Z}/2$-fixed sublocus under
$\cA\mapsto\cA^!$, where the Verdier involution
$(\cA^!)^!\simeq\cA$ holds strictly\textup{)}.\textup{]}

Let $\cA$ be a modular Koszul chiral algebra on~$X$ in the KSDual
sublocus.
Assume that the Verdier/Koszul dual branch admits the strict
finite-type representative $\cA^!$, or equivalently that the completed
continuous dual has been strictified under the Mittag--Leffler and
separatedness hypotheses.
Then:
\begin{enumerate}[label=\textup{(\roman*)}]
\item \emph{H-level dual transform.}
 There is a canonical equivalence in
 $D^{\mathrm{co}}(\operatorname{Fact}(X))$
 \begin{equation}\label{eq:frontier-holographic-koszul-dual}
 \eta_\cA\colon
 \mathbb H_X(\cA)\xrightarrow{\;\sim\;} \iota_X(\cA^!).
 \end{equation}
\item \emph{H-level anti-involution.}
 The second transform is defined after applying the same strictification
 to~$\cA^!$, and there is a canonical equivalence
 \begin{equation}\label{eq:frontier-holographic-double-transform}
 \eta_\cA^{(2)}\colon
 \mathbb H_X^2(\cA)\xrightarrow{\;\sim\;} \iota_X(\cA).
 \end{equation}
\item \emph{H/M-level cofiber identification.}
 For every $g\ge 0$,
 \begin{equation}\label{eq:frontier-holographic-cofiber-identification}
 \mathbf H^{\mathrm{hol}}_g(\cA)
 \simeq
 \mathbf Q_g(\cA^!).
 \end{equation}
\item \emph{S-level shadow equality.}
 For every $g\ge 0$,
 \begin{equation}\label{eq:frontier-holographic-shadow-equality}
 H^*_{\mathrm{hol},g}(\cA)\cong Q_g(\cA^!).
 \end{equation}
\end{enumerate}
In particular, on the proved modular Koszul core the protected transform
$\mathbb H_X$ lands in the typed image of the Verdier/Koszul-dual branch
and is a contravariant involution on the strictified locus.
\end{theorem}

\begin{proof}
By Verdier intertwining and higher-genus inversion, in the coderived
factorization category,
\[
\eta_\cA\colon
\Omega_X\mathbb D_{\operatorname{Ran}}\barB_X(\cA)
\xrightarrow{\;\sim\;}
\Omega_X\barB_X(\cA^!)
\xrightarrow{\;\sim\;}
\iota_X(\cA^!),
\]
which proves~\eqref{eq:frontier-holographic-koszul-dual}. Applying the same
argument again and using the strict Verdier involution
$(\cA^!)^!\simeq \cA$ on the same locus gives
\eqref{eq:frontier-holographic-double-transform}.

For the cofiber statement, the involution-splitting lemma gives a canonical
quasi-isomorphism
\[
\mathbf C_g(\cA)
\simeq
\mathbf Q_g(\cA)\oplus \mathbf Q_g(\cA^!).
\]
Under this direct-sum decomposition the canonical map
$\mathbf Q_g(\cA)\to \mathbf C_g(\cA)$ is the inclusion of the first summand,
so its cofiber is canonically the second summand. This proves
\eqref{eq:frontier-holographic-cofiber-identification}. Taking cohomology
proves~\eqref{eq:frontier-holographic-shadow-equality}.
\end{proof}

\begin{remark}[Dependency note for
Theorem~\textup{\ref{thm:frontier-protected-bulk-antiinvolution}}]
\label{rem:frontier-protected-bulk-antiinvolution-deps}
Theorem~\ref{thm:frontier-protected-bulk-antiinvolution} uses only the proved
bar--cobar/Verdier core and the involution-splitting lemma. It does \emph{not}
require the universal Maurer--Cartan class $\Theta_{\cA}$
(Theorem~\ref{thm:mc2-bar-intrinsic}) and it
does \emph{not} claim the existence of an honest three-dimensional bulk
factorization algebra. What it reconstructs is the boundary-controlled
\emph{protected} transform object on the curve~$X$, namely a typed
factorization representative of the Koszul-dual branch packaged by
$\mathbb H_X$.
\end{remark}

\subsection{Transposition, cotangent realization, and fixed-point rigidity}
\label{subsec:frontier-transposition-cotangent-fixed}

For the rest of this section fix a genus $g\ge 1$ and write
\begin{equation}\label{eq:frontier-dg-shift}
 d_g:=3g-3,
 \qquad
 V_g(\cA):=\mathbf Q_g(\cA),
 \qquad
 W_g(\cA):=\mathbf Q_g(\cA^!)[d_g].
\end{equation}
By Theorem~\ref{thm:quantum-complementarity-main} there is a perfect duality
\begin{equation}\label{eq:frontier-vg-wg-duality}
W_g(\cA)\simeq V_g(\cA)^\vee.
\end{equation}

\begin{theorem}[Protected transposition and cotangent realization;
\ClaimStatusProvedHere]
\label{thm:frontier-transposition-cotangent}
\index{transposition theorem!protected|textbf}
\index{cotangent realization!genus-$g$ ambient complex|textbf}
\textup{[}Type signature:
$(\mathrm{quadrant}=\mathrm{Open}\cap\mathrm{Cat},
\mathrm{presentation}=\mathrm{factorisation},
\mathrm{level}=3\to 4,
\mathrm{hypothesis\ package}=\mathrm{curved\text{-}central\ modular\ Koszul};
\mathrm{genus\ }g\ge 1;\mathrm{KSDual\ sublocus\ for\ part\ (iii)})$
\textup{(}Convention~\textup{\ref{conv:regime-tags}},
Convention~\textup{\ref{conv:hms-levels}}\textup{)}.\textup{]}

Let $\cA$ be a modular Koszul chiral algebra on~$X$.
Then:
\begin{enumerate}[label=\textup{(\roman*)}]
\item \emph{H-level shifted cotangent realization.}
 There is a canonical equivalence
 \begin{equation}\label{eq:frontier-shifted-cotangent-realization}
 \mathbf C_g(\cA)
 \simeq
 V_g(\cA)\oplus V_g(\cA)^\vee[-d_g].
 \end{equation}
\item \emph{H-level transposition.}
 The protected holographic cofiber identifies canonically with the shifted dual
 of the deformation branch:
 \begin{equation}\label{eq:frontier-transposition-shifted-dual}
 \mathbf H^{\mathrm{hol}}_g(\cA)
 \simeq
 V_g(\cA)^\vee[-d_g].
 \end{equation}
\item \emph{H-level fixed-point duality on the KSDual sublocus.}
 Assume $\cA$ lies in the KSDual sublocus, the
 $\mathbb{Z}/2$-fixed locus of $\cA\mapsto\cA^!$.
 If $\mathbb H_X(\cA)\simeq \cA$, then
 \begin{equation}\label{eq:frontier-self-holographic-shifted-self-duality}
 V_g(\cA)\simeq V_g(\cA)^\vee[-d_g].
 \end{equation}
\end{enumerate}
\end{theorem}

\begin{proof}
The decomposition theorem gives
\[
\mathbf C_g(\cA)
\simeq
\mathbf Q_g(\cA)\oplus \mathbf Q_g(\cA^!).
\]
By definition,
$W_g(\cA)=\mathbf Q_g(\cA^!)[d_g]$, hence
$\mathbf Q_g(\cA^!)\simeq W_g(\cA)[-d_g]\simeq V_g(\cA)^\vee[-d_g]$.
Substituting this into the direct-sum decomposition proves
\eqref{eq:frontier-shifted-cotangent-realization}. The cofiber statement then
follows from Theorem~\ref{thm:frontier-protected-bulk-antiinvolution}(iii).
Finally, if $\mathbb H_X(\cA)\simeq \cA$, then
$\mathbf Q_g(\mathbb H_X(\cA))\simeq \mathbf Q_g(\cA)$ and the previous part
identifies this with $V_g(\cA)^\vee[-d_g]$.
\end{proof}

\begin{lemma}[Determinant parity for shifted cotangent pairs;
\ClaimStatusProvedHere]
\label{lem:frontier-determinant-parity}
\index{determinant line!parity lemma|textbf}
Let $V$ be a finite-rank perfect complex and let $d\in\Z$.
Then:
\begin{enumerate}[label=\textup{(\alph*)}]
\item
 \begin{equation}\label{eq:frontier-det-dual-shift}
 \Det\!\bigl(V^\vee[-d]\bigr)
 \cong
 \begin{cases}
 \Det(V)^{-1}, & d\ \text{even},\\[4pt]
 \Det(V), & d\ \text{odd}.
 \end{cases}
 \end{equation}
\item Consequently,
 \begin{equation}\label{eq:frontier-det-cotangent-pair}
 \Det\!\bigl(V\oplus V^\vee[-d]\bigr)
 \cong
 \begin{cases}
 \C, & d\ \text{even},\\[4pt]
 \Det(V)^{\otimes 2}, & d\ \text{odd}.
 \end{cases}
 \end{equation}
\end{enumerate}
\end{lemma}

\begin{proof}
Duality gives $\Det(V^\vee)\cong \Det(V)^{-1}$.
A shift by $[1]$ inverts the determinant line, hence a shift by $[-d]$ inverts
it iff $d$ is odd. This proves~\eqref{eq:frontier-det-dual-shift}. The direct
sum formula for determinant lines then gives
\[
\Det\!\bigl(V\oplus V^\vee[-d]\bigr)
\cong
\Det(V)\otimes \Det\!\bigl(V^\vee[-d]\bigr),
\]
from which~\eqref{eq:frontier-det-cotangent-pair} is immediate.
\end{proof}

\begin{corollary}[Spectral reciprocity, palindromicity, and parity rigidity;
\ClaimStatusProvedHere]
\label{cor:frontier-spectral-reciprocity-palindromicity}
\index{spectral reciprocity!holographic|textbf}
\index{palindromicity!self-holographic object|textbf}
\index{Euler characteristic!even-genus vanishing|textbf}
\textup{[}Type signature:
$(\mathrm{Open}\cap\mathrm{Cat},\mathrm{scalar},\mathrm{level}=5,
\mathrm{hypothesis\ package}=\mathrm{finite\text{-}dimensional\ }H^*V_g(\cA);
\mathrm{KSDual\ sublocus\ for\ (ii)\text{--}(iv)})$\textup{]}
Assume $V_g(\cA)$ has finite-dimensional cohomology and set
\[
P_g(\cA;u):=\sum_i \dim H^i\!\bigl(V_g(\cA)\bigr)u^i.
\]
Then:
\begin{enumerate}[label=\textup{(\roman*)}]
\item \emph{S-level spectral reciprocity.}
 \begin{equation}\label{eq:frontier-spectral-reciprocity}
 P_g(\cA^!;u)=u^{d_g}P_g(\cA;u^{-1}).
 \end{equation}
\item \emph{S-level self-holographic palindromicity \textup{(}KSDual sublocus\textup{)}.}
 Assume $\cA$ lies in the KSDual sublocus.
 If $\mathbb H_X(\cA)\simeq \cA$, then
 \begin{equation}\label{eq:frontier-self-holographic-palindromicity}
 P_g(\cA;u)=u^{d_g}P_g(\cA;u^{-1}).
 \end{equation}
\item \emph{S-level even-genus vanishing \textup{(}KSDual sublocus\textup{)}.}
 Assume $\cA$ lies in the KSDual sublocus.
 If $\mathbb H_X(\cA)\simeq \cA$ and $g$ is even, then
 \begin{equation}\label{eq:frontier-even-genus-euler-vanishing}
 \chi\!\bigl(V_g(\cA)\bigr)=0.
 \end{equation}
\item \emph{S-level determinant parity \textup{(}KSDual sublocus\textup{)}.}
 Assume $\cA$ lies in the KSDual sublocus.
 \begin{equation}\label{eq:frontier-ambient-determinant-parity}
 \Det\!\bigl(\mathbf C_g(\cA)\bigr)
 \cong
 \begin{cases}
 \C, & g\ \text{odd},\\[4pt]
 \Det\!\bigl(V_g(\cA)\bigr)^{\otimes 2}, & g\ \text{even}.
 \end{cases}
 \end{equation}
\end{enumerate}
\end{corollary}

\begin{proof}
By Theorem~\ref{thm:frontier-transposition-cotangent},
$\mathbf Q_g(\cA^!)\simeq V_g(\cA)^\vee[-d_g]$.
Taking cohomology gives the reciprocity formula
\eqref{eq:frontier-spectral-reciprocity}. If $\mathbb H_X(\cA)\simeq \cA$,
this becomes palindromicity. Evaluating at $u=-1$ gives
\[
\chi\!\bigl(V_g(\cA)\bigr)=(-1)^{d_g}\chi\!\bigl(V_g(\cA)\bigr).
\]
Since $d_g=3g-3$ is odd exactly when $g$ is even,
\eqref{eq:frontier-even-genus-euler-vanishing} follows. The determinant
formula is Lemma~\ref{lem:frontier-determinant-parity} applied to
$V=V_g(\cA)$ and $d=d_g$.
\end{proof}

\begin{theorem}[Scalar fixed-point rigidity on a full scalar package and
genus-$1$ completeness;
\ClaimStatusProvedHere]
\label{thm:frontier-scalar-fixed-point-rigidity}
\index{scalar rigidity!self-holographic object|textbf}
\index{one-loop completeness!holographic scalar tower|textbf}
\textup{[}Type signature:
$(\mathrm{Open}\cap\mathrm{Cat},\mathrm{scalar},\mathrm{level}=5,
\mathrm{hypothesis\ package}=\mathrm{KSDual\ sublocus};
\mathrm{Theorem\ D\ uniform\text{-}weight\ scalar\ lane};
\mathrm{affine\ scalar\ duality\ law\ }(\alpha,\beta);
\mathrm{Genus\text{-}universality\ at\ all\ }g\ge 1)$\textup{]}
Assume $\cA$ lies in the KSDual sublocus and in a regime where the
full uniform-weight scalar package is known: scalar obstruction
classes and scalar free energies are governed by $\kappa(\cA)$ at
all genera, on the proved uniform-weight lane of Theorem~D
\textup{(}uniform conformal weights;
Definition~\textup{\ref{def:scalar-lane}}\textup{)}.
The scalar duality law is affine:
\begin{equation}\label{eq:frontier-affine-scalar-law}
\kappa(\cA^!)=\alpha-\beta\,\kappa(\cA)
\qquad
(\alpha,\beta\in\C).
\end{equation}
Then:
\begin{enumerate}[label=\textup{(\roman*)}]
\item If $\mathbb H_X(\cA)\simeq \cA$, then
 \begin{equation}\label{eq:frontier-fixed-point-kappa-hyperplane}
 (1+\beta)\kappa(\cA)=\alpha.
 \end{equation}
\item If $(\alpha,\beta)=(0,1)$, then every self-holographic fixed point is
 scalar-flat:
 \begin{equation}\label{eq:frontier-antisymmetric-fixed-point-flatness}
 \kappa(\cA)=0,
 \qquad
 \mathrm{obs}_g(\cA)=0,
 \qquad
 F_g(\cA)=0
 \quad \forall g\ge 1.
 \end{equation}
\item More generally, the genus-$1$ affine law propagates to all scalar genera:
 \begin{equation}\label{eq:frontier-all-genus-scalar-propagation}
 \mathrm{obs}_g(\cA^!)=
 \bigl(\alpha-\beta\kappa(\cA)\bigr)\lambda_g \qquad \textup{(}UNIFORM-WEIGHT\textup{}).
 \end{equation}
\end{enumerate}
\end{theorem}

\begin{proof}
The fixed-point relation $\mathbb H_X(\cA)\simeq \cA$ identifies $\cA$ with
$\cA^!$, so the affine law immediately gives
\eqref{eq:frontier-fixed-point-kappa-hyperplane}. If
$(\alpha,\beta)=(0,1)$, then $2\kappa(\cA)=0$, hence $\kappa(\cA)=0$ and the
assumed scalar genus tower vanishes. The all-genus propagation is
exactly the assumed scalar package:
\[
\mathrm{obs}_g(\cA^!)=\kappa(\cA^!)\lambda_g
=(\alpha-\beta\kappa(\cA))\lambda_g.
\]
The same substitution into the assumed scalar free-energy formula gives the vanishing in
part~(ii).
\end{proof}

\subsection{Linearized quantization: Heisenberg, Fock, Weyl, and the Fourier kernel}
\label{subsec:frontier-linearized-quantization}

The shifted cotangent realization equips the genus-$g$ ambient complex with a
linearized symplectic structure. Quantization proceeds as follows.

\begin{definition}[Genus-$g$ Heisenberg package]
\label{def:frontier-genus-g-heisenberg-package}
\index{Heisenberg algebra!genus-$g$ holographic|textbf}
\index{Fock object!genus-$g$ holographic|textbf}
For fixed $g\ge 1$ and modular Koszul $\cA$, write
\[
V:=V_g(\cA),
\qquad
W:=W_g(\cA)\simeq V^\vee.
\]
Define:
\begin{enumerate}[label=\textup{(\roman*)}]
\item the \emph{genus-$g$ Heisenberg dg Lie algebra}
 \begin{equation}\label{eq:frontier-heisenberg-dg-lie}
 \mathfrak h_g(\cA):=V\oplus W\oplus \C\mathbf 1,
 \end{equation}
 with $\mathbf 1$ central, $V$ and $W$ abelian, and bracket
 \begin{equation}\label{eq:frontier-heisenberg-bracket}
 [w,v]=w(v)\,\mathbf 1.
 \end{equation}
\item the \emph{boundary polynomial Fock object}
 \begin{equation}\label{eq:frontier-boundary-fock-object}
 \mathfrak F^{\partial}_g(\cA)
 :=\bigoplus_{n\ge 0}\operatorname{Sym}^n(V)
 \end{equation}
 and its symmetric-degree completion
 \begin{equation}\label{eq:frontier-completed-boundary-fock-object}
 \widehat{\mathfrak F}{}^{\partial}_g(\cA)
 :=\prod_{n\ge 0}\operatorname{Sym}^n(V).
 \end{equation}
\item the \emph{bulk polynomial Fock object}
 \begin{equation}\label{eq:frontier-bulk-fock-object}
 \mathfrak F^{\mathrm{bulk}}_g(\cA)
 :=\bigoplus_{n\ge 0}\operatorname{Sym}^n(W)
 \end{equation}
 and its completion
 \begin{equation}\label{eq:frontier-completed-bulk-fock-object}
 \widehat{\mathfrak F}{}^{\mathrm{bulk}}_g(\cA)
 :=\prod_{n\ge 0}\operatorname{Sym}^n(W).
 \end{equation}
\item the \emph{genus-$g$ Fourier kernel}
 \begin{equation}\label{eq:frontier-genus-g-fourier-kernel}
 \mathbf K_g^{\mathrm{hol}}(\cA)
 :=
 \exp\!\bigl(\operatorname{coev}_g(\cA)\bigr)
 \in
 \widehat{\mathfrak F}{}^{\partial}_g(\cA)
 \,\widehat\otimes\,
 \widehat{\mathfrak F}{}^{\mathrm{bulk}}_g(\cA),
 \end{equation}
 where $\operatorname{coev}_g(\cA)\in V\otimes W$ corresponds to the identity
 of~$V$.
\end{enumerate}
\end{definition}

\begin{theorem}[Heisenberg relations and Fourier transport;
\ClaimStatusProvedHere]
\label{thm:frontier-heisenberg-fourier-transport}
\index{Fourier kernel!intertwining theorem|textbf}
\index{Fock representation!Heisenberg theorem|textbf}
\textup{[}Type signature:
$(\mathrm{Open}\cap\mathrm{Cat},\mathrm{factorisation},
\mathrm{level}=4,
\mathrm{hypothesis\ package}=\mathrm{shifted\ cotangent\ realization\ of\ Theorem\ \ref{thm:frontier-transposition-cotangent}};
\mathrm{finite\text{-}rank\ perfect\ dg\ model\ for\ }V_g(\cA);
\mathrm{KSDual\ sublocus\ for\ canonical\ Fourier\ identification})$\textup{]}
Let $v\in V$ and $w\in W\simeq V^\vee$.
Let $m_v$ denote multiplication by $v$ on $\mathfrak F^{\partial}_g(\cA)$ and
let $\iota_w$ denote the derivation extending $w$.
Then:
\begin{enumerate}[label=\textup{(\roman*)}]
\item \emph{H-level Heisenberg commutators.}
 \begin{equation}\label{eq:frontier-heisenberg-commutators}
 [m_{v_1},m_{v_2}]=0,
 \qquad
 [\iota_{w_1},\iota_{w_2}]=0,
 \qquad
 [\iota_w,m_v]=w(v)\id.
 \end{equation}
\item \emph{H-level vacuum property.}
 The vacuum vector $|0\rangle:=1$ is annihilated by every $\iota_w$, and every
 vector of $\mathfrak F^{\partial}_g(\cA)$ is obtained by repeated action of
 multiplication operators on~$|0\rangle$.
\item \emph{H-level Fourier intertwining.}
 On the completed tensor product one has
 \begin{equation}\label{eq:frontier-fourier-intertwining-w}
 (\iota_w\otimes 1)\mathbf K_g^{\mathrm{hol}}(\cA)
 =
 (1\otimes m_w)\mathbf K_g^{\mathrm{hol}}(\cA),
 \end{equation}
 \begin{equation}\label{eq:frontier-fourier-intertwining-v}
 (m_v\otimes 1)\mathbf K_g^{\mathrm{hol}}(\cA)
 =
 (1\otimes \iota_v)\mathbf K_g^{\mathrm{hol}}(\cA),
 \end{equation}
 where on the bulk side $v$ acts by contraction through the dual pairing.
\item \emph{S-level plethystic reciprocity.}
 If $H^*(V)$ is finite-dimensional, then the cohomological Fock series satisfy
 \begin{equation}\label{eq:frontier-plethystic-reciprocity}
 \sum_{n\ge 0} q^n P\!\bigl(\operatorname{Sym}^n H^*(\mathbf Q_g(\cA^!));u\bigr)
 =
 \sum_{n\ge 0} q^n P\!\bigl(\operatorname{Sym}^n H^*(\mathbf Q_g(\cA));u^{-1}\bigr)
 u^{n d_g}.
 \end{equation}
\end{enumerate}
\end{theorem}

\begin{proof}
The first two parts are the standard creation--annihilation calculus on the
symmetric algebra. Because $W\simeq V^\vee$, each $w\in W$ defines a scalar
linear functional on $V$, hence a unique derivation $\iota_w$ of
$\operatorname{Sym}(V)$. A direct check on generators gives
$[\iota_w,m_v]=w(v)\id$.

For the intertwining identities, write
\[
\mathbf K_g^{\mathrm{hol}}(\cA)
=
\sum_{n\ge 0}\frac{1}{n!}\operatorname{coev}_g(\cA)^{\odot n}.
\]
Contraction on the left with $w$ removes one copy of the coevaluation tensor and
returns multiplication by $w$ on the right; multiplication by $v$ on the left is
dually contraction by $v$ on the right. This proves the transport identities.

For the plethystic formula, the spectral reciprocity
$\mathbf Q_g(\cA^!)\simeq V^\vee[-d_g]$ implies
\[
H^*\!\bigl(\mathbf Q_g(\cA^!)\bigr)
\cong
H^*(V)^\vee[-d_g].
\]
Symmetric powers commute with finite-dimensional duals, so taking Poincar\'e
polynomials termwise yields~\eqref{eq:frontier-plethystic-reciprocity}.
\end{proof}

\begin{definition}[Genus-$g$ Weyl object and Wick contraction]
\label{def:frontier-genus-g-weyl-wick}
\index{Weyl object!genus-$g$ holographic|textbf}
\index{Wick contraction!genus-$g$ holographic|textbf}
With $V$ and $W$ as above, define the completed Weyl object
\begin{equation}\label{eq:frontier-genus-g-weyl-object}
\widehat{\mathscr W}_g(\cA)
:=
\prod_{n\ge 0}\operatorname{Sym}^n(V\oplus W)[[\hbar]].
\end{equation}
Let
\[
\mu:\widehat{\mathscr W}_g(\cA)\widehat\otimes
\widehat{\mathscr W}_g(\cA)\to \widehat{\mathscr W}_g(\cA)
\]
denote ordinary graded-commutative multiplication.
The evaluation pairing $W\otimes V\to\C$ defines a continuous biderivation
\begin{equation}\label{eq:frontier-wick-contraction-operator}
\Gamma_g:
\widehat{\mathscr W}_g(\cA)\widehat\otimes
\widehat{\mathscr W}_g(\cA)
\longrightarrow
\widehat{\mathscr W}_g(\cA)\widehat\otimes
\widehat{\mathscr W}_g(\cA)
\end{equation}
characterized on generators by
\begin{equation}\label{eq:frontier-wick-contraction-on-generators}
\Gamma_g(w\otimes v)=w(v),
\qquad
\Gamma_g(v\otimes v')=\Gamma_g(v\otimes w)=\Gamma_g(w\otimes w')=0.
\end{equation}
Define the \emph{Wick star product}
\begin{equation}\label{eq:frontier-wick-star-product}
f\star_{g,\hbar} h
:=
\mu\circ e^{\hbar\Gamma_g}(f\otimes h).
\end{equation}
\end{definition}

\begin{theorem}[Associativity, PBW, and exact linear Weyl sewing;
\ClaimStatusProvedHere]
\label{thm:frontier-weyl-pbw-linear-sewing}
\index{PBW theorem!genus-$g$ Weyl object|textbf}
\index{Gaussian sewing!linear kernel|textbf}
\textup{[}Type signature:
$(\mathrm{Open}\cap\mathrm{Chain},\mathrm{chain\text{-}level},
\mathrm{level}=4,
\mathrm{hypothesis\ package}=\mathrm{Convention~\ref{conv:frontier-quadratic-gaussian-sector}\ \ \textup{(}finite\text{-}rank\ perfect\ dg\ model\textup{)}};
\mathrm{linear\ middle\ contraction};
\mathrm{KSDual\ sublocus\ for\ Verdier\text{-}consistent\ pairing})$\textup{]}
For each $g\ge 1$:
\begin{enumerate}[label=\textup{(\roman*)}]
\item \emph{H-level associativity.}
 The product $\star_{g,\hbar}$ is associative and unital on
 $\widehat{\mathscr W}_g(\cA)$.
\item \emph{M-level PBW property.}
 Filtering by total symmetric degree,
 \begin{equation}\label{eq:frontier-weyl-associated-graded}
 \operatorname{gr}\!\bigl(\widehat{\mathscr W}_g(\cA),\star_{g,\hbar}\bigr)
 \cong
 \operatorname{Sym}(V\oplus W)[[\hbar]].
 \end{equation}
\item \emph{H-level generator relations.}
 \begin{equation}\label{eq:frontier-weyl-generator-relations}
 v\star_{g,\hbar}v'=vv',
 \qquad
 w\star_{g,\hbar}w'=ww',
 \qquad
 v\star_{g,\hbar}w=vw,
 \qquad
 w\star_{g,\hbar}v=wv+\hbar w(v).
 \end{equation}
\item \emph{H-level exact linear sewing.}
 Let $\mathbf K_{ij}^{\mathrm{lin}}:=\exp(\operatorname{coev}_{ij})$ denote the
 linear Fourier kernel between the $i$-th $V$-copy and the $j$-th $W$-copy.
 Then middle contraction gives the exact sewing law
 \begin{equation}\label{eq:frontier-linear-gaussian-sewing}
 \langle \mathbf K_{12}^{\mathrm{lin}},\mathbf K_{23}^{\mathrm{lin}}\rangle_{(2)}
 =
 \mathbf K_{13}^{\mathrm{lin}}.
 \end{equation}
\end{enumerate}
\end{theorem}

\begin{proof}
Associativity is the same commuting-contractions argument as in the ordinary
Weyl algebra: on a triple tensor product the pairwise contraction operators
$\Gamma_{12},\Gamma_{13},\Gamma_{23}$ commute, so both possible parenthesizations
of a triple product equal
\[
\mu_{123}\circ
\exp\!\bigl(\hbar(\Gamma_{12}+\Gamma_{13}+\Gamma_{23})\bigr).
\]
The PBW statement follows because every application of $\Gamma_g$ lowers total
symmetric degree by $2$. The generator relations are immediate from the fact
that the only nontrivial contraction runs from a left $W$-generator to a right
$V$-generator. Finally,
\eqref{eq:frontier-linear-gaussian-sewing} is the special case of associativity
in which both kernels are pure exponentials of the coevaluation tensor and the
middle variables are integrated out by Wick contraction.
\end{proof}

\subsection{Gaussian kernels, Schur composition, and the metaplectic anomaly}
\label{subsec:frontier-metaplectic-gaussian-kernels}

The linear sewing law of
Theorem~\ref{thm:frontier-weyl-pbw-linear-sewing}(iv) sees only the coevaluation tensor.
Quadratic data introduces a determinant factor: the metaplectic correction.

\begin{convention}[Quadratic Gaussian sector]
\label{conv:frontier-quadratic-gaussian-sector}
For the rest of this subsection we work on a fixed finite-rank perfect dg model
for $V=V_g(\cA)$ and identify $W$ with $V^\vee$.
To keep the formulas readable, we state them in the parity-even sector.
Thus $V$ is represented by an ordinary finite-dimensional vector space over
$\C$, $W=V^\vee$, and all quadratic operators below have cohomological degree~$0$.
The same formulas hold on the graded model after inserting the standard Koszul
signs and replacing ordinary transpose by signed transpose.
\end{convention}

\begin{definition}[Gaussian triple and Gaussian kernel]
\label{def:frontier-gaussian-triple-kernel}
\index{Gaussian kernel!holographic|textbf}
\index{Schur composition!Gaussian triple|textbf}
Let $V$ be an $n$-dimensional complex vector space with coordinates
$x=(x_1,\dots,x_n)^t$ and dual coordinates $\eta=(\eta_1,\dots,\eta_n)^t$ on
$W=V^\vee$.
A \emph{Gaussian triple} is a triple $(B,A,C)$ consisting of
\begin{equation}\label{eq:frontier-gaussian-triple-data}
B\in\operatorname{Hom}(V,W),
\qquad
A\in\operatorname{End}(V),
\qquad
C\in\operatorname{Hom}(W,V),
\end{equation}
where $B$ and $C$ are symmetric in the sense that the matrices of $B$ and $C$
are symmetric in the chosen dual bases.
The corresponding \emph{Gaussian kernel} is
\begin{equation}\label{eq:frontier-gaussian-kernel}
\mathbf K(B,A,C)(x,\eta)
:=
\exp\!\left(
\frac12 x^t Bx + x^t A^t\eta + \frac12 \eta^t C\eta
\right)
\in
\widehat{\operatorname{Sym}}(V)\widehat\otimes \widehat{\operatorname{Sym}}(W).
\end{equation}
If $\mathbf K_1=\mathbf K(B_1,A_1,C_1)(x,\eta)$ and
$\mathbf K_2=\mathbf K(B_2,A_2,C_2)(\xi,y)$,
we define their \emph{middle contraction} by
\begin{equation}\label{eq:frontier-middle-contraction-definition}
\langle \mathbf K_1,\mathbf K_2\rangle_{(2)}
:=
\left.
\exp\!\left(
\frac12 \partial_\eta^t B_2\partial_\eta + y^t A_2\partial_\eta
+\frac12 y^t C_2 y
\right)
\mathbf K_1(x,\eta)
\right|_{\eta=0}.
\end{equation}
\end{definition}

\begin{lemma}[Formal Gaussian differential identity;
\ClaimStatusProvedHere]
\label{lem:frontier-formal-gaussian-differential-identity}
\index{Gaussian differential identity|textbf}
Let $S,T$ be symmetric $n\times n$ matrices over $\C$ and let $a,b\in\C^n$.
Assume $1-ST$ is invertible.
Then
\begin{align}
&\left.
\exp\!\left(\frac12 \partial_z^t T\partial_z + b^t\partial_z\right)
\exp\!\left(a^t z + \frac12 z^t S z\right)
\right|_{z=0}
\nonumber\\[4pt]
&\qquad=
\det(1-ST)^{-1/2}
\exp\!\left(
\frac12 b^t T(1-ST)^{-1}b
+ a^t(1-ST)^{-1}b
+ \frac12 a^t S(1-TS)^{-1}a
\right).
\label{eq:frontier-formal-gaussian-differential-identity}
\end{align}
\end{lemma}

\begin{proof}
Introduce the one-parameter family
\[
F(u,z)
:=
\exp\!\left(\frac{u}{2}\partial_z^t T\partial_z\right)
\exp\!\left(a^t z + \frac12 z^t S z\right).
\]
It satisfies the heat-type equation
\begin{equation}\label{eq:frontier-gaussian-heat-equation}
\partial_u F = \frac12\partial_z^t T\partial_z F,
\qquad
F(0,z)=\exp\!\left(a^t z + \frac12 z^t S z\right).
\end{equation}
Make the Gaussian ansatz
\begin{equation}\label{eq:frontier-gaussian-ansatz}
F(u,z)=
\exp\!\left(
\frac12 z^t S(u)z + a(u)^t z + c(u)
\right).
\end{equation}
Substituting into the differential equation and comparing coefficients of
quadratic, linear, and constant terms gives the ODE system
\begin{equation}\label{eq:frontier-gaussian-ode-system}
S'(u)=S(u)TS(u),
\qquad
a'(u)=S(u)Ta(u),
\qquad
c'(u)=\frac12 a(u)^tTa(u)+\frac12\operatorname{Tr}(TS(u)),
\end{equation}
with initial conditions $S(0)=S$, $a(0)=a$, $c(0)=0$.
The unique solutions are
\begin{align}
S(u)&=S(1-uTS)^{-1},
\label{eq:frontier-gaussian-ode-solution-S}\\
 a(u)&=(1-uST)^{-1}a,
\label{eq:frontier-gaussian-ode-solution-a}\\
 c(u)&=
 \frac12 a^tT(1-uST)^{-1}a
 -\frac12\log\det(1-uST).
\label{eq:frontier-gaussian-ode-solution-c}
\end{align}
Now translate by $b$ using $e^{b^t\partial_z}f(z)=f(z+b)$ and evaluate at
$z=0$. At $u=1$ this yields exactly
\eqref{eq:frontier-formal-gaussian-differential-identity}.
\end{proof}

\begin{theorem}[Gaussian composition, Schur complement, and determinant anomaly;
\ClaimStatusProvedHere]
\label{thm:frontier-gaussian-composition-schur-anomaly}
\index{Schur complement!Gaussian composition theorem|textbf}
\index{determinant anomaly!Gaussian composition theorem|textbf}
\textup{[}Type signature:
$(\mathrm{Open}\cap\mathrm{Chain},\mathrm{chain\text{-}level},
\mathrm{level}=4\to 5,
\mathrm{hypothesis\ package}=\mathrm{quadratic\ Gaussian\ sector\ \textup{(}Convention~\ref{conv:frontier-quadratic-gaussian-sector}\textup{)}};
\mathrm{invertibility\ of\ }1-C_1B_2;
\mathrm{KSDual\ sublocus\ for\ involution\text{-}consistent\ Schur\ symmetry})$\textup{]}
Let
\[
\mathbf K_1=\mathbf K(B_1,A_1,C_1)(x,\eta),
\qquad
\mathbf K_2=\mathbf K(B_2,A_2,C_2)(\xi,y)
\]
be Gaussian kernels and assume $1-C_1B_2$ is invertible.
Then:
\begin{enumerate}[label=\textup{(\roman*)}]
\item \emph{H/M-level exact composition formula.}
 Their middle contraction is again Gaussian:
 \begin{equation}\label{eq:frontier-gaussian-composition-main}
 \langle \mathbf K_1,\mathbf K_2\rangle_{(2)}
 =
 \det(1-C_1B_2)^{-1/2}\,
 \mathbf K(B_{12},A_{12},C_{12})(x,y),
 \end{equation}
 where
 \begin{align}
 B_{12}
 &:= B_1 + A_1^t B_2(1-C_1B_2)^{-1}A_1,
 \label{eq:frontier-schur-B}\\
 A_{12}^t
 &:= A_1^t(1-B_2C_1)^{-1}A_2^t,
 \label{eq:frontier-schur-A}\\
 C_{12}
 &:= C_2 + A_2 C_1(1-B_2C_1)^{-1}A_2^t.
 \label{eq:frontier-schur-C}
 \end{align}
\item \emph{H-level linear specialization.}
 For the pure linear kernels $(B_1,A_1,C_1)=(0,1,0)=(B_2,A_2,C_2)$,
 \eqref{eq:frontier-gaussian-composition-main} reduces to the exact linear
 sewing law
 \eqref{eq:frontier-linear-gaussian-sewing}.
\item \emph{S-level first nonlinear anomaly term.}
 Writing
 \begin{equation}\label{eq:frontier-alpha-bc}
 \alpha_g(C_1,B_2):=\det(1-C_1B_2)^{-1/2},
 \end{equation}
 one has the trace-log expansion
 \begin{equation}\label{eq:frontier-trace-log-anomaly-expansion}
 \log\alpha_g(C_1,B_2)
 =
 \frac12\sum_{m\ge 1}\frac{1}{m}\operatorname{Tr}\bigl((C_1B_2)^m\bigr)
 =
 \frac12\operatorname{Tr}(C_1B_2)
 + \frac14\operatorname{Tr}\bigl((C_1B_2)^2\bigr)+\cdots.
 \end{equation}
 The term $\tfrac14\operatorname{Tr}((C_1B_2)^2)$ is the first genuinely
 nonlinear correction to the linear sewing law.
\end{enumerate}
\end{theorem}

\begin{proof}
Apply Lemma~\ref{lem:frontier-formal-gaussian-differential-identity} with
\[
S=C_1,
\qquad
T=B_2,
\qquad
a=A_1x,
\qquad
b=A_2^t y.
\]
The resulting exponent is
\begin{align*}
&\frac12 x^t B_1x
+ \frac12 x^t A_1^t B_2(1-C_1B_2)^{-1}A_1x\\
&\qquad
+ x^t A_1^t(1-B_2C_1)^{-1}A_2^t y
+ \frac12 y^t C_2 y
+ \frac12 y^t A_2 C_1(1-B_2C_1)^{-1}A_2^t y,
\end{align*}
which is precisely the Gaussian kernel with coefficients
\eqref{eq:frontier-schur-B}--\eqref{eq:frontier-schur-C}.
This proves the exact composition formula.
The linear specialization is immediate.
The trace-log expansion is the standard power-series identity
\[
-\frac12\log\det(1-T)
=
\frac12\sum_{m\ge 1}\frac{1}{m}\operatorname{Tr}(T^m),
\]
applied to $T=C_1B_2$.
\end{proof}

\begin{remark}[Interpretation of
Theorem~\textup{\ref{thm:frontier-gaussian-composition-schur-anomaly}}]
\label{rem:frontier-gaussian-composition-interpretation}
The formulas
\eqref{eq:frontier-schur-B}--\eqref{eq:frontier-schur-C} are the Schur
complements of the middle quadratic block. The determinant factor is the
precise algebraic shadow of integrating out the intermediate holographic phase
space. The linear kernel of the previous subsection is the
$B=C=0$ specialization.
\end{remark}

\begin{definition}[Determinant anomaly line and metaplectic extension]
\label{def:frontier-determinant-anomaly-line}
\index{determinant anomaly line|textbf}
\index{metaplectic extension!Gaussian semigroup|textbf}
Let $(B_1,A_1,C_1)$ and $(B_2,A_2,C_2)$ be composable Gaussian triples.
\begin{enumerate}[label=\textup{(\roman*)}]
\item The \emph{determinant anomaly line} of the pair is
 \begin{equation}\label{eq:frontier-determinant-anomaly-line}
 \mathfrak a_g\bigl((B_1,A_1,C_1),(B_2,A_2,C_2)\bigr)
 :=
 \Det\!\bigl(\operatorname{cone}(\id_V-C_1B_2)\bigr)^{-1}.
 \end{equation}
 On the acyclic finite-dimensional locus this line carries the canonical
 torsion trivialization whose scalar is $\det(1-C_1B_2)^{-1}$.
\item A \emph{metaplectic lift} of the pair is a square root line
 \begin{equation}\label{eq:frontier-metaplectic-lift-line}
 \mathfrak m_g\bigl((B_1,A_1,C_1),(B_2,A_2,C_2)\bigr)
 \end{equation}
 together with an isomorphism
 \begin{equation}\label{eq:frontier-metaplectic-square-root}
 \mathfrak m_g\bigl((B_1,A_1,C_1),(B_2,A_2,C_2)\bigr)^{\otimes 2}
 \xrightarrow{\sim}
 \mathfrak a_g\bigl((B_1,A_1,C_1),(B_2,A_2,C_2)\bigr).
 \end{equation}
\item The \emph{Gaussian metaplectic semigroup} is the central extension whose
 multiplication law is twisted by the scalar cocycle
 $\alpha_g(C_1,B_2)=\det(1-C_1B_2)^{-1/2}$ on the acyclic locus.
\end{enumerate}
\end{definition}

\begin{theorem}[Cocycle law and metaplectic strictification;
\ClaimStatusProvedHere]
\label{thm:frontier-metaplectic-cocycle-strictification}
\index{cocycle!metaplectic Gaussian extension|textbf}
\index{strictification!metaplectically corrected sewing|textbf}
\textup{[}Type signature:
$(\mathrm{Open}\cap\mathrm{Cat},\mathrm{factorisation},
\mathrm{level}=4\to 5,
\mathrm{hypothesis\ package}=\mathrm{Theorem~\ref{thm:frontier-gaussian-composition-schur-anomaly}};
\mathrm{acyclic\ Gaussian\ sub\text{-}semigroup\ for\ part\ (iii)};
\mathrm{KSDual\ sublocus\ for\ Verdier\text{-}consistent\ central\ extension};
\mathrm{modularity\ here\ is\ the\ trace\ }\Tr_\cC\ \mathrm{on\ the\ open\ category},\ \mathrm{not\ a\ closed\text{-}algebra\ property})$\textup{]}
Let $G_i=(B_i,A_i,C_i)$ be three Gaussian triples such that all pairwise and
iterated compositions are defined.
Then:
\begin{enumerate}[label=\textup{(\roman*)}]
\item \emph{H/M-level $2$-cocycle.}
 If $G_1\sharp G_2$ denotes the Schur-complement triple produced by
 Theorem~\ref{thm:frontier-gaussian-composition-schur-anomaly}, then the scalar
 anomaly factors satisfy
 \begin{equation}\label{eq:frontier-metaplectic-2-cocycle}
 \alpha_g(C_1,B_2)\,\alpha_g\bigl(C_{12},B_3\bigr)
 =
 \alpha_g(C_2,B_3)\,\alpha_g\bigl(C_1,B_{23}\bigr),
 \end{equation}
 where $(B_{12},A_{12},C_{12})=G_1\sharp G_2$ and
 $(B_{23},A_{23},C_{23})=G_2\sharp G_3$.
\item \emph{H-level central extension.}
 Consequently, the multiplication law
 \begin{equation}\label{eq:frontier-metaplectic-central-extension-law}
 (\lambda_1,G_1)\cdot(\lambda_2,G_2)
 :=
 \bigl(\lambda_1\lambda_2\alpha_g(C_1,B_2),\; G_1\sharp G_2\bigr)
 \end{equation}
 defines an associative central extension of the Gaussian semigroup by
 $\C^{\times}$ on the acyclic locus.
\item \emph{H-level strictification after a metaplectic splitting.}
 If a chosen subsemigroup admits a $1$-cochain $\mu_g$ satisfying
 \begin{equation}\label{eq:frontier-metaplectic-splitting-equation}
 \alpha_g(C_1,B_2)
 =
 \mu_g(G_1)\mu_g(G_2)\mu_g(G_1\sharp G_2)^{-1},
 \end{equation}
 then the corrected kernels
 \begin{equation}\label{eq:frontier-corrected-gaussian-kernel}
 \widetilde{\mathbf K}(G):=\mu_g(G)^{-1}\mathbf K(G)
 \end{equation}
 compose strictly:
 \begin{equation}\label{eq:frontier-strict-corrected-sewing}
 \langle \widetilde{\mathbf K}(G_1),\widetilde{\mathbf K}(G_2)\rangle
 =
 \widetilde{\mathbf K}(G_1\sharp G_2).
 \end{equation}
\item \emph{S-level anomaly series.}
 Every such splitting is governed by the same universal trace-log series:
 \begin{equation}\label{eq:frontier-metaplectic-anomaly-series}
 \log\alpha_g(C_1,B_2)
 =
 \log\mu_g(G_1)+\log\mu_g(G_2)-\log\mu_g(G_1\sharp G_2)
 =
 \frac12\operatorname{Tr}(C_1B_2)
 + \frac14\operatorname{Tr}\bigl((C_1B_2)^2\bigr)+\cdots.
 \end{equation}
\end{enumerate}
\end{theorem}

\begin{proof}
Associativity of actual middle contraction gives
\[
\langle \langle \mathbf K(G_1),\mathbf K(G_2)\rangle,\mathbf K(G_3)\rangle
=
\langle \mathbf K(G_1),\langle \mathbf K(G_2),\mathbf K(G_3)\rangle\rangle.
\]
Expanding both sides with
Theorem~\ref{thm:frontier-gaussian-composition-schur-anomaly} and comparing the
unique Gaussian normal forms yields the cocycle identity
\eqref{eq:frontier-metaplectic-2-cocycle}. This is exactly the associativity
condition for the central extension law
\eqref{eq:frontier-metaplectic-central-extension-law}.
If $\mu_g$ splits the cocycle, substituting
\eqref{eq:frontier-corrected-gaussian-kernel} into the composition formula makes
the scalar anomaly factors cancel, giving strict sewing.
The anomaly series is the logarithm of the cocycle split and therefore begins
with the same trace-log expansion as
\eqref{eq:frontier-trace-log-anomaly-expansion}.
\end{proof}

\begin{corollary}[The first nonlinear holographic anomaly;
\ClaimStatusProvedHere]
\label{cor:frontier-first-nonlinear-holographic-anomaly}
\index{nonlinear anomaly!first holographic term|textbf}
Let $G_1,G_2$ be composable Gaussian triples. The deviation from the linear
sewing law is measured by the class
\begin{equation}\label{eq:frontier-first-nonlinear-holographic-anomaly}
\mathfrak N_g(G_1,G_2)
:=
\frac14\operatorname{Tr}\bigl((C_1B_2)^2\bigr).
\end{equation}
This vanishes precisely on the square-zero locus $(C_1B_2)^2=0$, in which case
Gaussian composition is controlled by the one-loop term
$\tfrac12\operatorname{Tr}(C_1B_2)$ alone.
\end{corollary}

\begin{proof}
This is the quadratic term of the universal trace-log anomaly series
\eqref{eq:frontier-trace-log-anomaly-expansion}. If $(C_1B_2)^2=0$, all higher
powers vanish and only the linear trace term survives.
\end{proof}

\subsection{The full protected picture and the remaining frontier}
\label{subsec:frontier-full-protected-picture}

\begin{remark}[Summary of proved results]
\label{rem:frontier-what-is-now-proved}
The previous subsections establish the following package within the
proved modular Koszul core.
\begin{enumerate}[label=\textup{(\roman*)}]
\item The protected holographic transform $\mathbb H_X$ is a contravariant
 anti-involution recovering the Koszul dual boundary package.
\item The genus-$g$ ambient complex is canonically a shifted cotangent object of
 the deformation branch, with the dual branch realized as its holographic
 cofiber.
\item Self-holographic objects satisfy palindromicity, even-genus Euler
 cancellation, determinant-parity constraints, and scalar fixed-point
 rigidity.
\item The linearized phase space admits canonical Heisenberg, Fock, Fourier, and
 Weyl quantizations.
\item Quadratic Gaussian kernels compose by Schur complement, and their sewing is
 twisted by a determinant anomaly whose logarithm is the universal trace-log
 series.
\item A metaplectic splitting strictifies this twisted sewing, upgrading
 the linear Fourier kernel to a nonlinear Gaussian holographic
 calculus.
\end{enumerate}
The first nonlinear holographic correction is proved within the algebraic core.
\end{remark}

\begin{conjecture}[Bulk enhancement: 3d geometric realization;
\ClaimStatusConjectured]

\label{conj:frontier-bulk-enhancement-universal-mc}
\index{bulk factorization algebra!holographic enhancement conjecture|textbf}
\index{Maurer--Cartan class!universal nonlinear modular conjecture|textbf}
\textup{[}Type signature:
$(\mathrm{Open}\cap\mathrm{Cat},\mathrm{factorisation},
\mathrm{level}=2\to 5,
\mathrm{hypothesis\ package}=\mathrm{modular\ Koszul\ }\cA\mathrm{\ on\ }(X,D,\tau);
\mathrm{Costello\text{--}Gwilliam\ globalisation\ data};
\mathrm{QME\ +\ descent\ +\ anomaly\ cancellation\ +\ locality\ \textup{(}MA\text{-}10\textup{)}};
\mathrm{KSDual\ sublocus\ for\ self\text{-}adjoint\ comparison\ map};
\mathrm{modularity\ realised\ as\ }\mathrm{Tr}_\cC\ +\ \mathrm{clutching\ on\ open\ category};
\mathrm{Vol\ III\ }(\Sigma_2,C)\text{-specialisation\ of\ Drinfeld\text{-}double\ comparison};
\mathrm{tree\text{-}level\ }+\ \mathrm{quartic\ obstruction\ scope;\ not\ a\ closed\text{-}algebra\ statement})$\textup{]}
Let $\cA$ be a modular Koszul chiral algebra on a smooth projective curve~$X$.
Then there exists an enhanced bulk package
\textup{(}Note: the algebraic objects $\Theta_{\cA}$,
the metaplectic anomaly, the scalar genus tower on its stated lane,
and the Gaussian sewing calculus are proved in the boundary algebraic
core. The conjectural content is a typed $3$d bulk object together
with comparison maps whose boundary images are those proved
objects.\textup{)}
\[
\mathfrak{Bulk}(\cA)
\]
with the following properties.
\begin{enumerate}[label=\textup{(\roman*)}]
\item \emph{H-level bulk realization.}
 $\mathfrak{Bulk}(\cA)$ is a factorization algebra on a genuine
 holomorphic-topological bulk space with boundary~$X$, equipped with a
 protected boundary map
 \[
 \rho_{\partial}\colon
 \mathfrak{Bulk}(\cA)|_X
 \longrightarrow
 \mathbb H_X(\cA)
 \]
 in $D^{\mathrm{co}}(\operatorname{Fact}(X))$. The conjecture asserts
 that $\rho_{\partial}$ is an equivalence on the protected sector.
\item \emph{H-level nonlinear genus control.}
 The deformation complex of $\mathfrak{Bulk}(\cA)$ carries a bulk
 Maurer--Cartan element $\Theta_{\mathfrak{Bulk}}$ and a comparison map
 \[
 \rho_{\mathrm{MC}}\colon
 \Def(\mathfrak{Bulk}(\cA))\longrightarrow \gAmod
 \]
 satisfying $\rho_{\mathrm{MC}}(\Theta_{\mathfrak{Bulk}})=\Theta_\cA$.
 Its scalar boundary shadow is
 $\kappa(\cA)\sum_{g\ge 1}\lambda_g$ on the scalar lane, and its
 quadratic boundary image is the metaplectic trace-log anomaly series of
 Theorem~\ref{thm:frontier-metaplectic-cocycle-strictification}.
\item \emph{M-level sewing.}
 After choosing a compatible metaplectic structure, the Gaussian semigroup of
 this section is the quadratic boundary truncation of a nonlinear kernel
 calculus closed under modular operadic gluing.
\item \emph{S-level shadow.}
 The resulting scalar shadow agrees with genus universality, the determinant
 shadow agrees with
 Corollary~\ref{cor:frontier-spectral-reciprocity-palindromicity}, and the
 first nonlinear anomaly agrees with
 Corollary~\ref{cor:frontier-first-nonlinear-holographic-anomaly}.
\end{enumerate}
\end{conjecture}

\begin{remark}[Platonic summary]
\label{rem:frontier-platonic-summary}
The protected algebraic core admits a first-principles reading:
bar records logarithmic collision data, cobar exponentiates it back to
the boundary algebra, Verdier duality exchanges deformation and
obstruction, and the genus-$g$ ambient complex is a shifted cotangent
phase space. The two algebraic branches are dual Lagrangian
polarizations of this boundary phase space. A physical bulk enters only
after a comparison map to a holomorphic-topological or chiral
Hochschild closed-sector object has been supplied. The Heisenberg and Weyl packages
quantize the linearized phase space; the Gaussian kernels are its
canonical wavefunctions; the determinant anomaly is the first
obstruction to naive sewing. The metaplectic correction encodes the
interaction of determinant-line geometry, Weyl quantization, and
gluing topology.
\end{remark}


%% ====================================================================
%% THE HOLOGRAPHIC MODULAR KOSZUL DATUM
%% ====================================================================
\subsection{The holographic modular Koszul datum}
\label{subsec:holographic-modular-koszul-datum}
\index{holographic modular Koszul datum|textbf}
\index{modular homotopy theory!holographic datum|textbf}

Every algebraic structure in this monograph is a Maurer--Cartan element
in a convolution algebra; the same holds for
holomorphic-topological holography.
The universal MC element $\Theta_\cA \in \MC(\gAmod)$
(Theorem~\ref{thm:mc2-bar-intrinsic}) is bar-intrinsic: it exists for
every modular Koszul chiral algebra~$\cA$, and its finite-order
projections are the shadow obstruction tower. The holographic lift
requires a typed comparison functor out of the HT open/closed
deformation problem. Without that functor the proved object is
$\Theta_\cA$ on the boundary algebra, not a physical bulk field.

Concretely, HT holography
literature (Costello--Dimofte--Gaiotto~\cite{CDG20},
Dimofte--Niu--Py~\cite{DNP25}, Gaiotto--Zeng~\cite{GZ26},
Khan--Zeng~\cite{KhanZeng25}, and
Latyntsev~\cite{Latyntsev23}) identifies the comparison surfaces of
the holographic datum below.  The datum is written here with the bar-dual
coalgebra exposed, because Verdier duality acts on that coalgebra and
does not identify it with the line algebra before finite-type or
completed continuous duality is supplied.  The content of the present
subsection is that these components are not independent data: they are
successive constructions from the bar, Verdier, Hochschild, collision,
MC, and connection packages attached to~$\cA$.

\begin{definition}[Holographic modular Koszul datum]

\label{def:holographic-modular-koszul-datum}
\index{holographic modular Koszul datum!definition}
Let $X$ be a smooth projective curve and let $T$ be a
holomorphic-topological system on $X\times\R_{\ge 0}$ whose boundary
chiral algebra~$\cA$ lies on the modular Koszul locus
\textup{(}Definition~\ref{def:koszul-locus}\textup{)}. On the
boundary-linear and Drinfeld--Sokolov comparison lanes of Vol~II, the
same notation is used for the resulting HT theory. The
\emph{holographic modular Koszul datum} of~$T$ is the image of the
universal MC element $\Theta_\cA \in \MC(\gAmod)$ under the holographic
shadow functor:
\begin{equation}\label{eq:holographic-datum}
\mathcal{H}(T)
\;=\;
\bigl(\cA,\;\cA^i,\;\cA^!,\;\mathcal{C},\;r(z),\;\Theta_\cA,\;
 \nabla^{\mathrm{hol}}_{\bullet,\bullet}\bigr).
\end{equation}
The seven entries carry the following content.
\begin{enumerate}[label=\textup{(\roman*)}]
\item \emph{Boundary chiral algebra.}
 $\cA$ is the boundary chiral algebra on~$X$: an augmented
 $\Eone$/chiral $A_\infty$-algebra, and the primary algebraic
 input from which all other entries are derived.
\item \emph{Bar-dual coalgebra.}
 $\cA^i:=H^*(B^{\mathrm{ch}}(\cA))$ is the cohomology of the
 chiral bar coalgebra.  It is a coalgebraic invariant of the bar
 construction, not the line-operator algebra.
\item \emph{Verdier/Koszul-dual branch.}
 $\cA^!$ is the strict Verdier dual algebra obtained from $\cA^i$
 on the finite-type modular Koszul locus, or from the corresponding
 completed continuous dual when the Mittag--Leffler and separatedness
 hypotheses are present.  Off that locus the canonical object supplied
 by the bar/Verdier formalism is the factorization companion
 $\cA^{!}_\infty$; strictification to a line algebra is an additional
 Verdier/Koszul rectification, not the bar-cobar inversion
 $\Omega^{\mathrm{ch}}B^{\mathrm{ch}}(\cA)\simeq\cA$.
\item \emph{Bulk/line comparison slot.}
 $\mathcal{C}$ denotes the third slot of the datum: the object through
 which bulk and line operators are compared. On the open/closed side one
 may take
 $\mathcal{C} = \mathcal{C}^{\bullet}_{\mathrm{ch}}(\cA,\cA)
 \simeq \cZ^{\mathrm{der}}_{\mathrm{ch}}(\cA)$
 by the chiral Swiss-cheese theorem
 \textup{(}Theorem~\ref{thm:thqg-swiss-cheese}\textup{)}. On the
 chirally Koszul line side, when a compact-generator comparison is
 available, the same slot is presented by the dg-category of line
 operators
 $\mathcal{C}_{\mathrm{line}} \simeq
 \cA^!_{\mathrm{line}}\text{-}\mathsf{mod}$
 \textup{(}Vol~II, theorem ``Lines as Modules for the Open-Colour
 Dual''\textup{)}. The
 global identification of these two presentations is precisely the
 global-triangle problem. The categorical annulus trace and a physical
 bulk, when present, enter through typed maps
 \[
 \mathrm{Tr}_{\mathcal C_{\mathrm{line}}}
 \longrightarrow
 \mathcal{C}^{\bullet}_{\mathrm{ch}}(\cA,\cA)
 \simeq
 \cZ^{\mathrm{der}}_{\mathrm{ch}}(\cA),
 \qquad
 \mathfrak{Bulk}(T)|_X
 \longrightarrow
 \cZ^{\mathrm{der}}_{\mathrm{ch}}(\cA).
 \]
 These maps are not automatic equivalences. The base curve~$X$ is
 ambient data and is not counted as a slot of~$\mathcal{H}(T)$.
\item \emph{Genus-zero binary shadow.}
 $r(z) \in \cA^! \,\widehat{\otimes}\, \cA^!\bigl[\![z^{-1},z]\!\bigr]$
 is the degree-one Maurer--Cartan element controlling the binary
 line-operator OPE. It satisfies the $A_\infty$ Yang--Baxter equation
 and identifies with the bar-cobar twisting morphism
 (Definition~\ref{def:dg-shifted-yangian};
 Remark~\ref{rem:yangian-twisting-morphism}). Off the strict dual
 locus the same statement lives in the completed companion
 $\cA^{!}_\infty$.
\item \emph{Modular Maurer--Cartan class.}
 $\Theta_\cA \in
 \mathrm{MC}\!\bigl(\Defcyc(\cA) \,\widehat{\otimes}\,
 R\Gamma(\overline{\mathcal{M}}_{\bullet,\bullet}^{\mathrm{st}},\,
 \mathbb{Q})\bigr)$
 is the full modular characteristic class
 (Theorem~\ref{thm:universal-MC};
 Theorem~\ref{thm:mc2-bar-intrinsic}). Its scalar shadow is
 $\kappa(\cA) \cdot \lambda_g$ (Theorem~\ref{thm:genus-universality}),
 and the finite-order projections are the shadow obstruction tower
 (Definition~\ref{def:shadow-postnikov-tower}).
 Note: $\Theta_\cA$ itself is proved; the conjectural content of
 this definition is the HT holographic packaging, the assertion
 that the seven-slot datum $\mathcal{H}(T)$ defines a bulk/boundary
 system.
\item \emph{Modular shadow connection.}
 $\nabla^{\mathrm{hol}}_{g,n}$ is the flat connection on derived
 coinvariants / protected states defined by the shadow of~$\Theta_\cA$;
 on benchmark comparison surfaces, it specializes to the
 corresponding derived conformal-block package
 \textup{(}Definition~\ref{def:modular-shadow-connection} below;
 Theorem~\ref{thm:sphere-reconstruction}\textup{)}.
\end{enumerate}
The bar-intrinsic Maurer--Cartan element $\Theta_\cA$ exists in the
generality of Theorem~\ref{thm:mc2-bar-intrinsic}; the additional
content of Definition~\ref{def:holographic-modular-koszul-datum} is the
typed assembly of the seven slots above into a single HT holographic
package. The Vol~I theorem proves the boundary modular Koszul
constructions. The HT assembly is conditional on the boundary-linear
and DS comparison maps recorded in Volume~II, and is conjectural as a
uniform statement for arbitrary HT inputs.
The critical discriminant $\Delta(\cA) = 8\kappa(\cA)\,S_4(\cA)$
(Theorem~\ref{thm:single-line-dichotomy}) is a derived invariant
of~$\mathcal{H}(T)$:
$\Delta = 0$ iff the scalar quartic line vanishes on the
single-line standard surface (classes $\mathbf{G}$ and $\mathbf{L}$);
$\Delta \neq 0$ detects the mixed Virasoro/$\mathcal W$ scalar line
(class~$\mathbf{M}$ on the standard surface). This statement is
scalar-line only. It does not assert vanishing of cross-channel
corrections or of the all-weight tower.
For the $\cH_k$ holographic system: $\Delta = 0$.
For the Virasoro holographic system:
$\Delta = 40/(5c{+}22) \neq 0$.

The four principles organizing~$\mathcal{H}(T)$ are:
\textup{(P1)}~HT holography compares derived chiral algebras, not only
their cohomology or planar OPE shadow;
\textup{(P2)}~the protected algebraic part of holography is a typed
Koszul-duality comparison between boundary protected operators and a
universal defect/boundary algebra;
\textup{(P3)}~the exact non-planar quantum algebra is the line-operator
category with meromorphic tensor product, encoded by $r(z)$/$R(z)$-type
data; and
\textup{(P4)}~HT holography is not complete until one has a
deformation-quantization functor from classical PVA/coisson/factorization
data to the quantum chiral/defect package.
\end{definition}

\begin{remark}[Relation to the modular cumulant transform]
\label{rem:holographic-cumulant-comparison}
\index{modular cumulant transform!holographic comparison}
The holographic datum $\mathcal{H}(T)$ is a
\emph{physics-packaged} version of the modular cumulant transform
$M(\cA) = (\hat{T}^c(sQ(\cA)^\vee), D_\cA, \tau_\cA, r_\cA(z),
\Theta_\cA)$
(Equation~\eqref{eq:modular-cumulant-transform} in the
concordance). The correspondence is:
\begin{center}
\small
\begin{tabular}{lcl}
\textbf{Holographic} && \textbf{Algebraic (cumulant)} \\
\hline
boundary algebra $\cA$ && source algebra \\[2pt]
defect algebra $\cA^!$ &&
 $\operatorname{Cum}_c(\cA)^\vee$ (cumulant dual) \\[2pt]
line category $\cA^!\text{-mod}$ &&
 $\operatorname{Cum}_c(\cA)^\vee\text{-mod}$ \\[2pt]
$r(z)$ && $\tau_\cA|_{\deg 2}$ (ordered $E_1$ face) \\[2pt]
$\Theta_\cA$ && connected stable-graph logarithm \\[2pt]
$\nabla^{\mathrm{hol}}$ && shadow projection of $\Theta_\cA$
\end{tabular}
\end{center}
The cumulant language adds three structures not visible in the
holographic packaging: the primitive spectrum $Q(\cA)$ and
generating series $G_\cA(t)$, which are \emph{kinematic} invariants
independent of OPE data; the completion entropy $h_K(\cA)$, which
measures boundary complexity; and the reduced-weight windows
$K_q(\cA)$, which make every jet computation finite. The
holographic modular Koszul datum $\mathcal{H}(T)$ is the
physics-side packaging of $M(\cA)$ after the strict Verdier dual,
Hochschild comparison, and completion hypotheses have been supplied;
the completion kinematics is its algebraic backbone.
\end{remark}

\begin{remark}[Bi-based Ran-space factorisation datum for $K3 \times E$ and Kodaira nearby cycles]
\label{rem:frontier-bi-based-ran-K3E}
\ClaimStatusConjectured
\index{Ran space!bi-based for K3 elliptic family}
\index{Kodaira fibre!nearby cycles}
\index{averaging morphism!Torelli-Kuga-Satake-Kodaira}
\index{factorisation infty-category!Francis-Gaitsgory holonomic}
The conjectural K3 import asks for a
\emph{bi-based Ran-space factorisation datum} when the boundary curve
$X$ is upgraded to the genus-one elliptic factor of an algebraic
$K3$-fibration $\pi\colon \cX \to A_2$ over the affine plane base
parametrising complex structures on the elliptic factor. The two bases
are:
\begin{enumerate}[label=\textup{(\roman*)}]
\item \emph{Chiral base.}
 $\operatorname{Ran}(E^{\mathrm{nod},\mathrm{sm}}_{24})$, the Ran
 space of the universal elliptic curve specialised to the
 $24$ Kodaira~$I_1$ nodes that record the
 Euler-characteristic factorisation
 $\chi(K3) = 24$
 \textup{(}each $I_1$ a smooth-vanishing cycle, contributing $1$
 to the Euler total\textup{)}; the chiral base carries the
 nearby-cycles $D$-modules at each $I_1$ node, recording the
 monodromy of the local family around the singular fibre.
\item \emph{Parameter base.}
 $\overline{A_2}$, the compactified $A_2$-base of the elliptic
 family, equipped with the Francis--Gaitsgory factorisation
 $\infty$-category in holonomic $D$-modules
 \textup{(}\cite{FG12}\textup{)} regular-singular along the
 hyperplane arrangement $H_1 \cup H_4$
 (the locus of geometric multiplicity~$1$ and~$4$ degenerations
 in the Kodaira list).
\end{enumerate}
The candidate averaging morphism on the bi-based datum is a typed
pushforward on the indicated factorization categories, written
symbolically as
\begin{equation}\label{eq:frontier-bi-based-averaging}
\operatorname{av}
\;=\;
\operatorname{Torelli}_{K3} \;\circ\;
\operatorname{KugaSatake} \;\circ\;
j_{\mathrm{Kodaira}},
\end{equation}
the composite of the Kodaira nearby-cycles inclusion
$j_{\mathrm{Kodaira}}\colon
E^{\mathrm{nod},\mathrm{sm}}_{24} \hookrightarrow E^{\mathrm{univ}}$,
the Kuga--Satake construction sending an elliptic family to its
attached abelian variety
$\operatorname{KugaSatake}\colon
\pi_*\mathbb Q_E \to A^{\mathrm{KS}}$,
and the Torelli map for $K3$ surfaces
$\operatorname{Torelli}_{K3}\colon A^{\mathrm{KS}} \to
\Omega_{K3}/\Gamma_{K3}$
landing in the period domain quotiented by the orthogonal monodromy
group~$\Gamma_{K3} \subset O(2,19)$. The composite~\eqref{eq:frontier-bi-based-averaging}
is the geometric incarnation of the averaging map
$\operatorname{av}\colon \fg^{E_1} \to \fg^{\mathrm{mod}}$
of the chiral programme: the Kodaira-nearby-cycles step
extracts the $E_1$ ordered factorisation data from the boundary curve;
the Kuga--Satake step couples that data to the Hodge-theoretic side
of the K3 family; the Torelli step lands the result in the symmetric
modular convolution algebra~$\fg^{\mathrm{mod}}$, where
$\Theta_\cA$~lives. The bi-based assembly
$(\operatorname{Ran}, \overline{A_2})$ would make the K3-coupled
chiral algebra into a genuine bulk/boundary factorisation
$\infty$-system, with the holographic
datum~$\mathcal{H}(T)$ of
Definition~\ref{def:holographic-modular-koszul-datum}
arising as the fibre at a generic basepoint of $\overline{A_2}$ over
the smooth Kodaira locus. The proved Vol~I input on each chamber is
only the modular Koszul boundary calculus. The K3 global gluing,
nearby-cycle comparison, Kuga--Satake lift, and Torelli descent are
the conjectural content of this remark.
\end{remark}

\begin{definition}[Modular shadow connection: the central formula]
\label{def:modular-shadow-connection}
\index{modular shadow connection|textbf}
\index{derived coinvariants!modular shadow connection}
\index{protected states!modular shadow connection}
For each $(g,n)$ with $2g - 2 + n > 0$, let
$\mathcal{V}_{g,n}(\cA)$ denote the derived coinvariant /
protected-state package at genus~$g$ with $n$ insertions. On
benchmark comparison surfaces, this specializes to the
corresponding derived conformal-block package. The
\emph{modular shadow connection} is
\begin{equation}\label{eq:modular-shadow-connection}
\boxed{%
\nabla^{\mathrm{hol}}_{g,n}
\;:=\;
d \;-\; \operatorname{Sh}_{g,n}(\Theta_\cA).}
\end{equation}
This is the central formula of the chapter: it converts the abstract
MC element~$\Theta_\cA$ into a concrete flat connection on the
derived coinvariant / protected-state package. Here
$\operatorname{Sh}_{g,n}\colon
\mathrm{MC}\bigl(\Defcyc(\cA)\,\widehat{\otimes}\,
R\Gamma(\overline{\mathcal{M}}_{\bullet,\bullet},\mathbb{Q})\bigr)
\to \operatorname{End}\bigl(\mathcal{V}_{g,n}(\cA)\bigr)$
is the shadow representation: project~$\Theta_\cA$ to the
$(g,n)$-sector and act on this package via the cyclic deformation
complex. Flatness $(\nabla^{\mathrm{hol}}_{g,n})^2 = 0$ is the
$(g,n)$-projection of the MC equation for~$\Theta_\cA$.
At $g=0$ it reduces, on the established sphere comparison
surfaces, to the commuting differential constraints determining
sphere correlators exactly~\cite{GZ26};
at higher genus it is the boundary shadow equation for protected
correlators. It becomes a physical holographic field equation only
after the comparison functor from the HT bulk problem to this boundary
package has been constructed.
\end{definition}

\subsubsection{The master equation and compatibility laws}

The entire datum is subordinated to a single equation:
\begin{equation}\label{eq:holographic-master-equation}
d\Theta_\cA
\;+\;
\sum_{m \ge 2} \frac{1}{m!}\,
\ell_m(\Theta_\cA, \dotsc, \Theta_\cA)
\;=\; 0,
\end{equation}
where $\{\ell_m\}_{m \ge 2}$ are the higher brackets of the
modular $L_\infty$ convolution algebra~$\gAmod$
(Theorem~\ref{thm:modular-homotopy-convolution}). In the strict dg~Lie
model $\Convstr$, this collapses to
$D\Theta_\cA + \tfrac{1}{2}[\Theta_\cA,\Theta_\cA] = 0$,
which is proved (Theorem~\ref{thm:mc2-bar-intrinsic};
Theorem~\ref{thm:convolution-d-squared-zero}).

Two compatibility laws constrain~$\Theta_\cA$ beyond the MC equation:
\begin{enumerate}[label=\textup{(\alph*)}]
\item \emph{Collision residue.}
\begin{equation}\label{eq:collision-residue-identification}
\operatorname{Res}^{\mathrm{coll}}_{0,2}(\Theta_\cA)
\;=\; r(z).
\end{equation}
The binary genus-zero collision residue of the full modular class
recovers the dg-shifted Yangian element. The trivalent collision
residue
$\operatorname{Res}^{\mathrm{coll}}_{0,3}(\Theta_\cA)$
encodes the $A_\infty$ Yang--Baxter relation:
the two-point stratum gives the binary MC element, and the
three-point boundary intersections give coassociativity and
Yang--Baxter.
\item \emph{Clutching / factorization.}
\begin{equation}\label{eq:clutching-factorization-frontier}
\xi^*(\Theta_\cA)
\;=\;
\Theta_\cA \star \Theta_\cA,
\end{equation}
where $\star$ is the operadic composition along stable boundary
strata $\xi\colon \overline{\mathcal{M}}_{g_1,n_1+1} \times
\overline{\mathcal{M}}_{g_2,n_2+1} \to \overline{\mathcal{M}}_{g,n}$
(already proved: Theorem~\ref{thm:universal-MC}(ii)).
\end{enumerate}
The collision claim used below is:
\emph{the dg-shifted Yangian is the binary genus-zero shadow of the
modular class~$\Theta_\cA$.}
That is, $r(z)$ is the first visible collision residue of the full
modular deformation problem. The 2026 exact sphere-correlator
match~\cite{GZ26} is the genus-zero matrix-element shadow of the
pair~$(\mathcal{C}, r)$.

\subsubsection{Five theorem targets: projections of~$\Theta_\cA$}

Theorems A, B, C, D, and H of the algebraic engine are different
typed projections of the bar-intrinsic class~$\Theta_\cA$. Only
Theorem~D is the scalar projection through~$\kappa(\cA)$. The bar
projection gives A/B, the Verdier-Lagrangian projection gives C, and
the chiral Hochschild projection gives H. The five targets below are
the analogous restrictions to a collision stratum, genus sector, or
degree truncation of the convolution algebra~$\gAmod$.

\begin{conjecture}[Boundary--defect realization: modular lift; \ClaimStatusConjectured]
\label{conj:boundary-defect-realization}
\index{boundary-defect realization!holographic}
\textup{[}Type signature:
$(\mathrm{Open}\cap\mathrm{Cat},\mathrm{factorisation},
\mathrm{level}=2\to 4,
\mathrm{hypothesis\ package}=\mathrm{large\ transversal\ boundary\ condition\ }B;
\mathrm{KK\text{-}reduction\ of\ HT\ bulk};
\mathrm{KSDual\ sublocus\ for\ canonical\ identification\ of\ }\cA^!_\mathrm{univ};
\mathrm{Vol\ III\ }(\Sigma_{d-1},C)\text{-specialisation\ of\ stage\text{-}2\ shadow};
\mathrm{modularity\ realised\ as\ }\Tr_\cC\ +\ \mathrm{clutching\ on\ open\ category},\ \mathrm{not\ as\ closed\text{-}algebra\ property})$\textup{]}
For a large transversal boundary condition~$B$, the boundary algebra
of the KK-reduced HT theory is equivalent, as a derived chiral algebra,
to the universal defect algebra of the original theory:
\begin{equation}\label{eq:boundary-defect-realization}
\cA_\partial^{\mathrm{transv}}
\;\simeq\;
\cA_{\mathrm{univ}}^!.
\end{equation}
The mathematical upgrade over existing work~\textup{\cite{GZ26}} is to
prove this in the modular category where $\Theta_\cA$ lives, not merely
at the level of cohomology or planar OPEs.
\end{conjecture}

\begin{remark}[Reduction of boundary--defect realization]
\label{rem:boundary-defect-reduction}
At genus~$0$ and the level of derived categories, the identification
$\cA^!_{\mathrm{univ}} \simeq \Omega(\barBch(\cA^!))$ is
bar-cobar inversion (Theorem~B) applied to~$\cA^!$:
Theorem~\ref{thm:bar-cobar-isomorphism-main} recovers the
boundary algebra from its bar coalgebra.
The protected dual-transform theorem
\textup{(}Theorem~\textup{\ref{thm:frontier-protected-bulk-antiinvolution})}
then gives $\mathbb{H}_X(\cA) \simeq \cA^!$. What remains
conjectural is the identification of the \emph{geometric} KK-reduced
boundary algebra (produced by dimensional reduction of the 3d HT
bulk) with the \emph{algebraic} Koszul dual, at the level of the
modular category carrying $\Theta_\cA$. This requires a 3d
geometric input: constructing the bulk factorization algebra on
$X \times \mathbb{R}_{\geq 0}$ whose boundary restriction is~$\cA^!$.
\end{remark}

\begin{theorem}[Yangian-shadow theorem; \ClaimStatusConditional]
\label{thm:yangian-shadow-theorem}
\index{Yangian!shadow theorem|textbf}
\index{collision residue!Yangian shadow}
\textup{[}Type signature:
$(\mathrm{Open}\cap\mathrm{Cat},\mathrm{factorisation},
\mathrm{level}=2\to 4,
\mathrm{hypothesis\ package}=\mathrm{modular\ Koszul};
\mathrm{KSDual\ sublocus\ for\ canonicity\ of\ }r(z);
\mathrm{Verdier\ window\ for\ part\ (ii)})$\textup{]}
Let $\cA$ be a modular Koszul chiral algebra. Then the binary
collision residue of~$\Theta_\cA$ canonically produces a dg-shifted
Yangian structure on~$\cA^!$:
\begin{equation}\label{eq:yangian-shadow}
\operatorname{Res}^{\mathrm{coll}}_{0,2}(\Theta_\cA) = r(z),
\qquad
\operatorname{Res}^{\mathrm{coll}}_{0,3}(\Theta_\cA)
\;\Longrightarrow\;
A_\infty\text{-}\mathrm{YBE}.
\end{equation}
\end{theorem}

\begin{proof}
The two assertions are the content of the two proved recovery theorems
in \S\ref{subsec:collision-filtration-recovery}.

\smallskip\noindent
\emph{Binary collision $\Rightarrow$ $r$-matrix.}\;
Theorem~\ref{thm:collision-residue-twisting} identifies the
genus-zero, two-point collision residue of~$\Theta_\cA$ with the
universal twisting morphism
$\pi_\cA \in \mathrm{Tw}(\barBch(\cA), \cA)$.
Passing first to the intrinsic bar-dual coalgebra
$\cA^i := H^*(\barBch(\cA))$ and then, in a finite Verdier window,
to coordinates in $(\cA^!)^\vee$ produces the spectral
$r$-matrix $r_\cA(z)$ of
Definition~\ref{def:dg-shifted-yangian}; this is
Theorem~\ref{thm:collision-residue-twisting}(ii).

\smallskip\noindent
\emph{Ternary collision $\Rightarrow$ $A_\infty$-YBE.}\;
The MC equation restricted to the genus-zero, four-point sector and
projected to collision depth~$2$ gives the infinitesimal braid
relation (Theorem~\ref{thm:collision-depth-2-ybe}), which is
equivalent to the classical Yang--Baxter equation for $r_\cA(z)$
via partial fractions on $\overline{\mathcal{M}}_{0,4}$.
The mechanism is the Arnold relation
$\eta_{12}\wedge\eta_{23} + \eta_{23}\wedge\eta_{31}
+ \eta_{31}\wedge\eta_{12} = 0$ tensored with the
structure-constant data of~$\cA$.

\smallskip\noindent
No quasi-linearity hypothesis is needed: the collision residue theorem
applies to any modular Koszul~$\cA$ with well-defined bar complex.
\end{proof}

\begin{remark}[Relation to earlier proof strategy]
\label{rem:yangian-shadow-strategy}
The proof above executes the strategy of restricting the modular MC
equation to two- and three-vertex collision strata. The pattern is
visible in the DNP formulas~\cite{DNP25} and the bar-cobar
identification of $r(z)$ as twisting morphism
(Remark~\ref{rem:yangian-twisting-morphism}).\end{remark}

\begin{remark}[Hall--Drinfeld double, Manin-pair signature, and Schiffmann--Vasserot shuffle]
\label{rem:frontier-hall-drinfeld-vs-yangian}
\ClaimStatusConditional
\index{Hall algebra!Drinfeld double obstruction}
\index{Drinfeld double!Yangian distinction}
\index{Manin pair!signature obstruction}
\index{Schiffmann--Vasserot!shuffle presentation}
\index{Yangian!collision residue origin}
The dg-shifted Yangian produced by the binary collision residue
$r(z) = \operatorname{Res}^{\mathrm{coll}}_{0,2}(\Theta_\cA)$
\textup{(}Theorem~\ref{thm:yangian-shadow-theorem}\textup{)} is
constitutionally distinct from the Drinfeld double of a Hall
algebra, despite both objects carrying spectral parameters and
trigonometric/rational~$R$-matrices. The distinction is invariant
data of the underlying Manin pair, not a presentation-dependent
choice. Concretely:
\begin{enumerate}[label=\textup{(\roman*)}]
\item \emph{Manin-pair signature.}
 A Hall--Drinfeld double $D(H)$ of a hereditary abelian
 category~$\mathsf A$ over a Lagrangian sub-Hopf-algebra of itself
 produces a self-dual Manin pair of signature $(n,n)$:
 the canonical bilinear form on $D(H)$ is split, and the
 Lagrangian decomposition is part of the data. The dg-shifted
 Yangian
 $Y_\hbar(\fg) \subset \cA \otimes \cA^!$
 produced by~$r(z)$ on the modular Koszul locus carries the
 \emph{Verdier} pairing
 $\cA \otimes \cA^! \to \mathbf 1$
 of signature $(2,1)$ (two algebra slots paired against one bulk
 slot via the chiral derived centre): the natural pairing has no
 Lagrangian decomposition because the Verdier involution
 $\sigma\colon \cA \leftrightarrow \cA^!$ is not part of the data
 of a sub-Hopf-algebra. The signature obstruction is irreducible.
\item \emph{No Lagrangian decomposition.}
 The Lagrangian decomposition
 $D(H) = H \oplus H^*$
 of the Drinfeld double is forced by the bilinear form of
 signature $(n,n)$. On our side the scalar complementarity class
 $\varkappa(\cA) = \kappa(\cA) + \kappa(\cA^!)$
 is the numerical shadow of the obstruction to such a Lagrangian
 splitting:
 $\varkappa = 0$ recovers a self-dual Lagrangian on the
 Heisenberg/affine/free-field families; $\varkappa \neq 0$
 obstructs every such splitting on Virasoro/$\mathcal{W}_3$/BP.
 The Yangian is therefore an extension class, not a doubled
 sub-Hopf-algebra.
\item \emph{Shuffle presentation.}
 On the cohomological-Hall side, the
 Schiffmann--Vasserot shuffle algebra
 \textup{(}see~\cite{Latyntsev23} for the chiral-side translation,
 and~\cite{Jindal-Kaubrys-Latyntsev2603} for the critical CoHA
 version\textup{)} presents the Drinfeld double~$D(H_{\mathsf A})$
 of a curve-class CoHA via shuffle products on a polynomial
 representation. The chiral-side dg-shifted Yangian of
 \textup{Definition~\ref{def:dg-shifted-yangian}} admits a
 different presentation: as the deformation quantization of the
 collision residue~$r(z)$ at the prefactor $\hbar = 1/(k+h^\vee)$
 \textup{(}Theorem~\ref{thm:yangian-sklyanin-quantization}\textup{)}.
 The two presentations agree on the shared Heisenberg face
 \textup{(}both reduce to the abelian shuffle on a
 polynomial Fock space\textup{)} but diverge on the
 affine Kac--Moody face, where the chiral version inherits
 the trace-form prefactor~$k$ from the bar-complex normalization
 while the Hall double inherits a quiver-dependent
 normalization from the Euler form on~$\mathsf A$.
\end{enumerate}
This distinction matters at every point in
Part~\ref{part:standard-landscape} where the Yangian-shadow
identification appears: the dg-shifted Yangian is the binary
genus-zero shadow of the modular MC element~$\Theta_\cA$, not the
Drinfeld double of a Hall algebra; the absence of a Lagrangian
decomposition is the algebraic content of $\varkappa(\cA) \neq 0$.
The point is proved at the level of the Manin-pair signature. Its
upgrade to a universal obstruction class is conditional on the
shadow-tower TQFT interpretation of
Remark~\ref{rem:shadow-tower-tqft-taylor}; the explicit dictionary between the
Schiffmann--Vasserot shuffle and the Vol~I bar-cobar twisting morphism
on a single comparison surface (e.g.\ rational curves with two
punctures) is the content
of~\cite{Latyntsev23} and is the chiral-side translation of the
CoHA-vertex correspondence.
\end{remark}

\begin{theorem}[Sphere flatness and reconstruction on established comparison surfaces; \ClaimStatusConditional]
\label{thm:sphere-reconstruction}
\index{sphere correlators!reconstruction from modular shadow|textbf}
\textup{[}Type signature:
$(\mathrm{Open}\cap\mathrm{Cat},\mathrm{factorisation},
\mathrm{level}=4\to 5,
\mathrm{hypothesis\ package}=\mathrm{modular\ Koszul};
\mathrm{boundary\ chiral\ Ward\ identity\ on\ comparison\ surface};
\mathrm{KZ\ analytic\ SDR\ on\ }\cM_{0,n};
\mathrm{flat\text{-}section\ existence})$\textup{]}
Let $\cA$ be a modular Koszul chiral algebra and let
$V_1,\dotsc,V_n$ be $\cA$-modules. Then:
\begin{enumerate}[label=\textup{(\roman*)}]
\item The genus-$0$, $n$-point shadow connection
\begin{equation}\label{eq:sphere-reconstruction}
\nabla^{\mathrm{hol}}_{0,n}
\;=\;
d \;-\; \operatorname{Sh}_{0,n}(\Theta_\cA)
\end{equation}
is flat: $(\nabla^{\mathrm{hol}}_{0,n})^2 = 0$.
\item On established genus-$0$ comparison surfaces where the chiral
Ward identity is known \textup{(}affine Kac--Moody, Virasoro, and
their $\mathcal{W}$-algebra descendants via DS reduction\textup{)},
the sphere correlators
$\langle \phi_1(z_1)\dotsm\phi_n(z_n)\rangle_{0,n}$
are the joint horizontal sections of\/ $\nabla^{\mathrm{hol}}_{0,n}$.
\item For $\cA = \widehat{\fg}_k$:
$\nabla^{\mathrm{hol}}_{0,n} = \nabla^{\mathrm{KZ}}$
\textup{(}Theorem~\textup{\ref{thm:shadow-connection-kz})}.
For $\cA = \mathrm{Vir}_c$:
$\nabla^{\mathrm{hol}}_{0,n}$ gives the Virasoro conformal Ward
connection, and on the degenerate-representation comparison surface
the higher collision residues yield the BPZ null-vector equations
\textup{(}Proposition~\textup{\ref{prop:shadow-connection-bpz})}.
\end{enumerate}
\end{theorem}

\begin{proof}
\emph{Flatness.}\;
The shadow connection is constructed from the MC element
$\Theta_\cA \in \mathrm{MC}(\gAmod)$
(Theorem~\ref{thm:mc2-bar-intrinsic}). The MC equation
$D_\cA\Theta_\cA + \tfrac{1}{2}[\Theta_\cA,\Theta_\cA]=0$
projected to genus~$0$ at $n$~points gives
$(\nabla^{\mathrm{hol}}_{0,n})^2 = 0$: the curvature of the shadow
connection is the MC failure, which vanishes.

\smallskip\noindent
\emph{Horizontal sections = correlators.}\;
For algebras where the chiral Ward identity is established (affine
Kac--Moody algebras, the Virasoro algebra, and their $\mathcal{W}$-algebra
descendants via DS reduction), the Ward identity on $\mathcal{M}_{0,n}$
identifies the $n$-point correlator as a multivalued function annihilated
by the connection $d - \operatorname{Sh}_{0,n}(\Theta_\cA)$.
The shadow representation
$\operatorname{Sh}_{0,n}\colon \gAmod \to
\operatorname{End}(\mathcal{V}_{0,n})$
maps the MC element to the connection form acting on the space
of conformal blocks on the corresponding comparison surface,
realized there by the specialized package
$\mathcal{V}_{0,n}(\cA;V_1,\dotsc,V_n)$.
The horizontal-section condition
$\nabla^{\mathrm{hol}}_{0,n}\,\Phi = 0$
is therefore exactly the system of differential equations satisfied
by $n$-point genus-$0$ correlators.

\emph{Recovery of classical connections.}\;
Part~(iii) is the content of
Theorem~\ref{thm:shadow-connection-kz} (affine KM $\Rightarrow$ KZ)
and
Proposition~\ref{prop:shadow-connection-bpz}
(Virasoro $\Rightarrow$ BPZ).
\end{proof}

\begin{theorem}[GZ26 commuting differentials from the MC element;
\ClaimStatusConditional]
\label{thm:gz26-commuting-differentials}
\index{Gaiotto--Zeng!commuting differentials|textbf}
\index{shadow connection!commuting Hamiltonians}
\textup{[}Type signature:
$(\mathrm{Open}\cap\mathrm{Cat},\mathrm{factorisation},
\mathrm{level}=4\to 5,
\mathrm{hypothesis\ package}=\mathrm{modular\ Koszul};
\mathrm{boundary\ chiral\ Ward\ identity\ for\ parts\ (iii)\text{--}(iv)};
\mathrm{Prochazka\ }\Walg_N\mathrm{\ closure\ +\ KZ\ analytic\ SDR\ for\ part\ (v)})$\textup{]}
Let\/ $\cA$ be a modular Koszul chiral algebra and
$V_1,\dotsc,V_n$ be $\cA$-modules with conformal blocks
$\cV_{0,n}(\cA;V_1,\dotsc,V_n)$ on~$\cM_{0,n}$.
Write $z_1,\dotsc,z_n$ for the marked points with
$z_{ij}=z_i-z_j$.
\begin{enumerate}[label=\textup{(\roman*)}]
\item \textup{(Hamiltonian decomposition.)}
 The shadow connection
 $\nabla^{\mathrm{hol}}_{0,n}=d-\operatorname{Sh}_{0,n}(\Theta_\cA)$
 decomposes into commuting Hamiltonians:
 \begin{equation}\label{eq:gz26-hamiltonian-decomposition}
 \operatorname{Sh}_{0,n}(\Theta_\cA)
 \;=\;
 \sum_{i=1}^{n}H_i\,dz_i,
 \qquad
 [H_i,\,H_j]=0
 \quad\text{for all }i,j.
 \end{equation}
 Commutativity is the flatness
 $(\nabla^{\mathrm{hol}}_{0,n})^2=0$, the genus-$0$ projection of the MC equation
 \textup{(}Theorem~\textup{\ref{thm:sphere-reconstruction})}.
\item \textup{(Collision-depth expansion.)}
 Each Hamiltonian decomposes by collision depth:
 \begin{equation}\label{eq:hamiltonian-collision-depth}
 H_i
 \;=\;
 \sum_{k=1}^{k_{\max}}
 \sum_{j\neq i}
 \frac{
 \operatorname{Res}^{\mathrm{coll}}_{0,k}(\Theta_\cA)
 \big|_{(i,j)}
 }{z_{ij}^k},
 \end{equation}
 where $k_{\max}=p_{\max}-1$ by $d\log$ absorption; see
 Chapter~\ref{ch:three-invariants}.
\item \textup{(Affine KM: KZ Hamiltonians.)}
 For $\cA=\widehat{\fg}_k$, $k_{\max}=1$:
 $H_i = (k+h^\vee)^{-1}\sum_{j\neq i}\Omega_{ij}/z_{ij}$,
 the KZ Hamiltonians.
\item \textup{(Virasoro: BPZ operators.)}
 For $\cA=\mathrm{Vir}_c$, $k_{\max}=3$:
 $H_i = \sum_{j\neq i}[h_j/z_{ij}^2 + \partial_{z_j}/z_{ij}]$
 acting on weight-$h_j$ primaries; the depth-$3$ central term
 $(c/2)/z_{ij}^3$ vanishes by the global $\SL(2)$ Ward identity.
\item \textup{($\cW_N$ prediction; finite-spin endpoint hypotheses required.)}
 For $\cA=\cW_N$ with $k_{\max}=2N-1$, the
 $H_i$ are differential operators of order $2N-2$, and the existence
 and commutativity follow from the MC equation. \emph{Statement of the
 endpoint package}: identifying these $H_i$ with the principal-$\cW_N$
 quantum higher Hamiltonians at the level of operator coefficients
 requires the Prochazka closure of $\Walg_N$ at finite~$N$
 \textup{(}strong generation by spin-$2$ through spin-$N$ currents,
 finite OPE\textup{)}, the Yamada generic-level finite presentation,
 and the analytic SDR for the KZ propagator on $\cM_{0,n}$. The
 finite-spin generation alone does not entail the full theorem: term-by-term
 GZ\textup{26} comparison is a separate normalization problem and
 fails without the endpoint package.
\end{enumerate}
\end{theorem}

\begin{proof}
Part~(i): the shadow connection is a flat connection on
$\cM_{0,n}\cong\Conf_n(\bP^1)/\PGL_2$ valued in $\End(\cV_{0,n})$
\textup{(}Theorem~\textup{\ref{thm:sphere-reconstruction}})}. The
connection form depends on positions only through differences $z_{ij}$
because the OPE data defining $\Theta_\cA$ is translation invariant
and the bar propagator $\mathrm d\log(z_i-z_j)$ depends only on
$z_{ij}$. Flatness gives
$\partial_i H_j - \partial_j H_i + [H_i,H_j] = 0$, which reduces to
$[H_i,H_j]=0$.

Part~(ii): the collision residue at depth $k$ acts as a differential
operator of order $k-1$ in $\partial_{z_j}$, by the mode expansion
$a_{(n)}b\to\partial_{z_j}^{n-k}/(n-k)!$.

Parts~(iii)--(iv): direct computation from
Theorem~\textup{\ref{thm:shadow-connection-kz}} and
Proposition~\textup{\ref{prop:shadow-connection-bpz}}. The vanishing
of the depth-$3$ central term $(c/2)/z_{ij}^3$ on primary states
follows because it is a scalar multiple of the identity, and the
global conformal Ward identity $\sum_i\partial_{z_i}=0$ on conformal
blocks makes $\sum_{j\neq i}(c/2)/z_{ij}^3$ an $\SL(2)$-coboundary
that drops out on primary states.

Part~(v): for $\cW_N$, the maximal OPE pole between the spin-$N$
generator and itself is $2N$, giving $k_{\max}=2N-1$ and operator
order $2N-2$.
\end{proof}

\begin{remark}[Term-by-term comparison with GZ\textup{26} for $\cW_N$ is conjectural;
\ClaimStatusConjectured]
\label{rem:gz26-wn-comparison-conjectural}
The existence and commutativity of $\cW_N$ Hamiltonians from
$\Theta_{\cW_N}$ are proved by Theorem~\textup{\ref{thm:gz26-commuting-differentials}}(v).
The conjectural content is the term-by-term identification with the
operators of \textup{\cite{GZ26}}, which depends on a normalization
convention match that is not yet established.
\end{remark}

\begin{remark}[Scope of the GZ26 comparison]
\label{rem:gz26-scope}
Theorem~\textup{\ref{thm:gz26-commuting-differentials}} derives the
\emph{existence} and \emph{algebraic structure} of the commuting
Hamiltonians from $\Theta_\cA$. The 2026 interface paper~\cite{GZ26}
constructs these operators specifically for interface minimal model
holography. The KZ and BPZ recoveries (parts~iii--iv) are exact;
part~(v) is a concrete prediction.
\end{remark}

\begin{theorem}[Gaudin--Yangian identification; \ClaimStatusConditional]
\label{thm:gaudin-yangian-identification}
\index{Gaudin model!Yangian identification|textbf}
\index{dg-shifted Yangian!Gaudin model}
\textup{[}Type signature:
$(\mathrm{Open}\cap\mathrm{Cat},\mathrm{factorisation},
\mathrm{level}=4,
\mathrm{hypothesis\ package}=\mathrm{modular\ Koszul};
\mathrm{KSDual\ sublocus\ for\ Verdier\ identification};
\mathrm{Drinfeld\ Yangian\ classical\ limit\ at\ }\hbar=(k+h^\vee)^{-1})$\textup{]}
Let\/ $\cA$ be a modular Koszul chiral algebra with dg-shifted Yangian
$Y^{\mathrm{dg}}_\cA$.
\begin{enumerate}[label=\textup{(\roman*)}]
\item For $\cA = \widehat{\fg}_k$ with $k_{\max}=1$, the commuting
 Hamiltonians of Theorem~\textup{\ref{thm:gz26-commuting-differentials}}
 are the Gaudin Hamiltonians of the Yangian $Y(\fg)$:
 $H_i^{\mathrm{GZ}} = (k+h^\vee)^{-1} H_i^{\mathrm{Gaudin}}$
 where $H_i^{\mathrm{Gaudin}} = \sum_{j\neq i}\Omega_{ij}/(z_i-z_j)$.
\item \textup{(Higher Gaudin Hamiltonians.)}
 For general\/ $\cA$ with $k_{\max}>1$, the higher collision residues
 give a higher Gaudin system:
 the depth-$k$ Hamiltonian
 $\sum_{j\neq i}\operatorname{Res}^{\mathrm{coll}}_{0,k}(\Theta_\cA)|_{(i,j)}/z_{ij}^k$
 is the transferred $A_\infty$ operation $m_k$ on
 $\cA^!_{\mathrm{line}}$ acting at spectral separation $z_{ij}$.
\item The quantization parameter is $\hbar = 1/(k+h^\vee)$.
\end{enumerate}
\end{theorem}

\begin{proof}
Part~(i): substituting the Casimir collision residue
$r(z) = \Omega/((k+h^\vee)z)$
\textup{(}Theorem~\textup{\ref{thm:yangian-shadow-theorem}}\textup{)}
into the depth-$1$ Hamiltonian formula yields exactly the Gaudin
Hamiltonian of \cite{FFR94} multiplied by $1/(k+h^\vee)$.

Part~(ii): by Theorem~\textup{\ref{thm:gz26-commuting-differentials}}(ii),
the depth-$k$ contribution to $H_i$ is the collision residue at depth
$k$. Each such residue is the $m_k$ operation transferred via the
Homological Perturbation Lemma to the cohomology, viewed as a
spectral-parameter deformation of the classical $m_2$.

Part~(iii): the prefactor $1/(k+h^\vee)$ is the Drinfeld quantization
parameter \cite{Drinfeld85}.
\end{proof}

\begin{theorem}[Three-parameter identification of $\hbar$]
\label{thm:yangian-sklyanin-quantization}
\index{quantization parameter!three-parameter identification|textbf}
\index{Sklyanin bracket!Yangian quantization}
\index{Khan--Zeng!Sklyanin quantization}
\textup{[}Type signature:
$(\mathrm{Open}\cap\mathrm{Cat},\mathrm{factorisation},
\mathrm{level}=4,
\mathrm{hypothesis\ package}=\cA=\widehat{\fg}_k,\ k\ne -h^\vee;
\mathrm{Drinfeld\ Yangian\ classical\ limit};
\mathrm{Khan\text{--}Zeng\ sigma\text{-}model\ gauge\ invariance};
\mathrm{DNP\ loop\ parameter\ identification};
\mathrm{KSDual\ sublocus\ for\ Verdier\text{-}consistent\ Casimir\ rescaling};
\mathrm{Vol\ III\ }(\Sigma_2,C)\text{-specialisation\ at\ }d=3,\ C\ \mathrm{the\ KZ\ bar\text{-}propagator\ curve})$\textup{]}
Let\/ $\cA = \widehat{\fg}_k$ and let\/ $Y_\hbar(\fg)$ be the Yangian
quantization of $U(\fg[u])$ in the sense of Drinfeld~\cite{Drinfeld85}.
The classical Sklyanin--Poisson structure on $(\fg^!)^*$ from the
level-$k$ rational $r$-matrix $r^{\mathrm{cl}}(z) = k\Omega/z$
(vanishing at $k=0$) is the classical
limit $\hbar\to 0$ of the Yangian commutator
\textup{(}Semenov-Tian-Shansky~\cite{STS83}, Drinfeld~\cite{Drinfeld85}\textup{)}.
The new content of this theorem is the three-parameter identification:
\begin{equation}\label{eq:three-parameter-hbar}
\hbar
\;=\;
\underbrace{\frac{1}{k+h^\vee}}_{\substack{\text{KZ\textup{25} sigma-}\\\text{model coupling}}}
\;=\;
\underbrace{\frac{1}{k+h^\vee}}_{\substack{\text{DNP\textup{25} loop}\\\text{parameter}}}
\;=\;
\underbrace{\frac{1}{k+h^\vee}}_{\substack{\text{collision-residue}\\\text{prefactor}}}.
\end{equation}
The three a priori distinct parameters, defined in three different
papers via three different constructions, are the same scalar
because they all arise from the same algebraic datum: the level $k$ of
the affine algebra $\widehat{\fg}_k$, which controls the bar-complex
normalization.
\end{theorem}

\begin{proof}
The classical limit identification is Drinfeld's theorem
\cite{Drinfeld85}; the Sklyanin bracket interpretation is
\cite{STS83}. These are classical results, used here without reproof.

The new content is the three-parameter identification. The KZ\textup{25}
sigma-model coupling on $\bC\times\bR_{\geq 0}$ is fixed by gauge
invariance to $1/(k+h^\vee)$ \textup{(\cite[\S\,2.2]{KhanZeng25})}.
The DNP\textup{25} perturbative expansion of line-operator OPE is organized
by the same parameter $1/(k+h^\vee)$ \textup{(\cite[\S\,3]{DNP25})}.
The collision residue
$r(z) = \operatorname{Res}^{\mathrm{coll}}_{0,2}(\Theta_\cA)
= k\,\Omega_{\mathrm{tr}}/z$
in the trace-form normalization. The KZ normalization is
$r^{\mathrm{KZ}}(z)=\Omega/((k+h^\vee)z)$ after Casimir and coupling
rescaling; it is the same finite collision datum in a different
normalization, not a literal equality of rational functions
\textup{(}Theorem~\textup{\ref{thm:collision-residue-twisting}}\textup{)}.

The three independent constructions agree because they measure the
same algebraic datum: the level $k$ controls the normalization of the
chiral bracket $J^a_{(1)}J^b = k\delta^{ab}$, which propagates through
the Sugawara construction to the bar complex on our side, the 1-loop
determinant on the KZ\textup{25} side, and the line-operator OPE residue on
the DNP\textup{25} side.
\end{proof}

\begin{theorem}[OPE pole order trichotomy for GZ\textup{26} Hamiltonians;
\ClaimStatusConditional]
\label{thm:shadow-depth-operator-order}
\index{operator-order trichotomy|textbf}
\index{three invariants $p_{\max},k_{\max},r_{\max}$}
\textup{[}Type signature:
$(\mathrm{Open}\cap\mathrm{Cat},\mathrm{factorisation},
\mathrm{level}=4\to 5,
\mathrm{hypothesis\ package}=\mathrm{modular\ Koszul};
\mathrm{three\ invariants}\ p_{\max},k_{\max},r_{\max}\
\mathrm{computed\ on\ a\ chosen\ minimal\ strong\ generating\ set})$\textup{]}
The trichotomy below is a \emph{generator-OPE} statement at fixed strong
generators; it does not by itself decide the full $\cW_N$ Hamiltonian
package, which requires Prochazka finite-OPE closure and the KZ
analytic SDR.
Let\/ $\cA$ be a modular Koszul chiral algebra with three invariants
$p_{\max},k_{\max},r_{\max}$ \textup{(}Chapter~\textup{\ref{ch:three-invariants})}.
\begin{enumerate}[label=\textup{(\roman*)}]
\item \textup{($k_{\max}=0$: trivial Hamiltonians.)}
 If $p_{\max}\leq 1$ \textup{(}only $\beta\gamma$ in the standard
 landscape\textup{)}, then $k_{\max}=0$ and the GZ\textup{26} Hamiltonians are
 trivial: there is no nonzero contribution at degree~$2$.
\item \textup{($k_{\max}=1$: multiplication operators.)}
 If $p_{\max}=2$ \textup{(}Heisenberg, affine KM\textup{)}, then $k_{\max}=1$
 and the GZ\textup{26} Hamiltonians are multiplication operators on
 $\cV_{0,n}$, acting by Casimir-type matrices without derivatives.
\item \textup{($k_{\max}\geq 3$: differential operators.)}
 If $p_{\max}\geq 4$ \textup{(}Virasoro: $p_{\max}=4$;
 $\cW_N$: $p_{\max}=2N$\textup{)}, then $k_{\max}\geq 3$ and the
 GZ\textup{26} Hamiltonians are genuine differential operators of order
 $k_{\max}-1=p_{\max}-2$.
\end{enumerate}
The operator order is a genus-$0$ invariant depending only on the
generating OPE \textup{(}via $p_{\max}$\textup{)}, while the shadow
depth $r_{\max}$ may exceed it: for $\beta\gamma$, $p_{\max}=1$ but
$r_{\max}=4$, and the higher-degree shadow data arises from composite-field
channels, not from the generator OPE. This is the sharp form of
the separation between generator pole order, collision depth, and
shadow depth.
\end{theorem}

\begin{proof}
Each case follows from Theorem~\textup{\ref{thm:gz26-commuting-differentials}}(ii).
At depth $k$, the collision residue acts on $V_j$ by a differential
operator of order $k-1$ \textup{(}or by zero if $k=0$\textup{)}. The
trichotomy reflects the standard-landscape values
$p_{\max}\in\{1,2\}\cup\{4,5,\ldots\}$:
the value $p_{\max}=3$ is excluded because no bosonic generator self-OPE
has maximal pole order $3$ \textup{(}Remark~\textup{\ref{rem:k-max-2-missing}}\textup{)}.
The independence of $r_{\max}$ from $p_{\max}$ for $\beta\gamma$ is
established in Chapter~\ref{ch:three-invariants}.
\end{proof}

\begin{remark}[Scope of the GZ26 commuting differentials]
\label{rem:sphere-reconstruction-gz26}
The 2026 interface paper~\cite{GZ26} constructs commuting differential
operators on the sphere conformal block space. The theorem above gives
the Vol~I boundary construction from $\Theta_\cA$ on the stated
comparison surfaces. The comparison with~\cite{GZ26} is exact for KZ
and BPZ normalizations; for $\mathcal W_N$ it remains the
normalization problem recorded in
Remark~\ref{rem:gz26-wn-comparison-conjectural}.
\end{remark}

\subsubsection{Khan--Zeng classical-to-quantum bridge}

\begin{theorem}[Classical-to-quantum bridge: proved algebraic content;
\ClaimStatusConditional]
\label{thm:kz-classical-quantum-bridge}
\index{Khan--Zeng!classical-to-quantum bridge|textbf}
\index{deformation quantization!proved algebraic content}
\textup{[}Type signature:
$(\mathrm{Open}\cap\mathrm{Cat},\mathrm{factorisation},
\mathrm{level}=2\to 3,
\mathrm{hypothesis\ package}=\mathrm{freely\ generated\ PVA\ with\ }
\lambda\mathrm{\text{-}Jacobi};
\mathrm{Genus\text{-}1\ BV\ obstruction\ for\ part\ (ii)};
\mathrm{Yamada\ analytic\ closure\ at\ all\ genera};
\mathrm{KZ\ analytic\ SDR\ on\ }\cM_{0,n})$\textup{]}
The classical PVA datum alone is not the all-loop quantum HT theory;
genus-$1$ requires a separate BV obstruction calculation, and the
all-genus identification requires the Yamada finite-presentation
hypothesis together with the KZ analytic SDR.
Let\/ $\cP$ be a freely generated Poisson vertex algebra with
$\lambda$-bracket $\{a_\lambda b\}=\sum_{n\geq 0}\lambda^{(n)}c_n$
satisfying the $\lambda$-Jacobi identity. Let $\cA$ be a modular
Koszul chiral algebra with $\cP\cong H^*(\cA,Q)/(\hbar)$ as
associated classical shadow \textup{(}cf.\ the examples part of Vol.~II\textup{)}.
\begin{enumerate}[label=\textup{(\roman*)}]
\item \textup{(Genus-$0$ seed.)} The $\lambda$-bracket determines the
 genus-$0$ component of the bar differential. Quadratic duality
 \textup{(}\cite{GLZ22}\textup{)} produces a classical $r$-matrix
 $r^{\mathrm{cl}}(z)$ satisfying CYBE; its quantum deformation is the
 collision residue
 $r(z)=r^{\mathrm{cl}}(z)+\hbar r^{(1)}(z)+\cdots = \operatorname{Res}^{\mathrm{coll}}_{0,2}(\Theta_\cA)$
 \textup{(}Theorem~\textup{\ref{thm:collision-residue-twisting}})}.
\item \textup{(Genus-$1$ obstruction.)} The BV operator $\Delta_{\mathrm{cyc}}$
 controls the lift from genus~$0$ to genus~$1$; for $\cW_3$ the
 obstruction vanishes \textup{(}Theorem~\textup{\ref{thm:w3-genus1-curvature}}\textup{)}.
\item \textup{(All-genera completion.)} The full MC element
 $\Theta_\cA$ exists at all genera by
 Theorem~\textup{\ref{thm:mc2-bar-intrinsic}}.
\item \textup{(Gauge invariance $=$ Jacobi; genus~$0$.)} The KZ\textup{25}
 sigma-model gauge invariance \textup{(\cite[\S\,2.2]{KhanZeng25})} is
 the $\lambda$-Jacobi identity, equivalent to $d^2_{\barB}=0$ via the
 Arnold relation. At higher genus, the geometric-algebraic
 identification is conjectural \textup{(}see
 Conjecture~\textup{\ref{conj:master-bv-brst}}\textup{)}.
\end{enumerate}
\end{theorem}

\begin{proof}
Part~(i) follows from Theorem~\textup{\ref{thm:collision-residue-twisting}}
combined with~\cite{GLZ22}. Parts~(ii) and~(iii) are
Theorem~\textup{\ref{thm:w3-genus1-curvature}} and
Theorem~\textup{\ref{thm:mc2-bar-intrinsic}} respectively. Part~(iv) at
genus~$0$ follows from the Arnold relation
\textup{(}Theorem~\textup{\ref{thm:bar-nilpotency-complete}}\textup{)};
the higher-genus extension is the standing
Conjecture~\textup{\ref{conj:master-bv-brst}}.
\end{proof}

\begin{theorem}[Quartic resonance obstruction; \ClaimStatusConditional]
\label{thm:quartic-resonance-obstruction}
\index{quartic resonance class!obstruction to modular extension|textbf}
\textup{[}Type signature:
$(\mathrm{Open}\cap\mathrm{Cat},\mathrm{factorisation},
\mathrm{level}=3\to 4,
\mathrm{hypothesis\ package}=\mathrm{constructive\ shadow\ tower};
\mathrm{vanishing\ of\ class\text{-}}M\,o_4\mathrm{\ requires\ KZ\ analytic\ SDR\ +\
weight\text{-}completed\ ambient\ for\ Class\ M})$\textup{]}
The first genuinely nonlinear obstruction to extending the
constructive shadow obstruction tower $\Theta_\cA^{\leq 3}$ to
$\Theta_\cA^{\leq 4}$ is the quartic resonance class
$R^{\mathrm{mod}}_{4,g,n}(\cA)$
\textup{(}Definition~\textup{\ref{def:nms-modular-quartic-resonance-class}};
Appendix~\textup{\ref{app:nonlinear-modular-shadows})}:
\begin{equation}\label{eq:quartic-resonance-obstruction}
R^{\mathrm{mod}}_{4,g,n}(\cA) = 0
\quad\Longleftrightarrow\quad
\Theta_\cA^{\leq 3}\text{ extends to }
\Theta_\cA^{\leq 4}.
\end{equation}
The clutching law
for~$R^{\mathrm{mod}}_{4,g,n}$ is established in
Appendix~\textup{\ref{app:nonlinear-modular-shadows}}.
The class $R^{\mathrm{mod}}_{4,g,n}(\cA)$ is canonically identified
with the $L_\infty$ obstruction class $[o_4(\cA)]$ of
Theorem~\textup{\ref{thm:quartic-obstruction-linf}}.
The obstruction vanishes automatically for shadow-depth classes
$G$, $L$, $C$; for class~$M$ it is generically nonzero.
\end{theorem}

\begin{proof}
This is Theorem~\ref{thm:quartic-obstruction-linf} restated in
holographic language. The identification of the quartic resonance
class with the $L_\infty$ obstruction to extending
$\Theta_\cA^{\leq 3}$ is given there via the clutching formula
\[
\xi^*\mathfrak{Q}
= H \star \mathfrak{Q}
+ \mathfrak{C}\star\mathfrak{C}
+ \mathfrak{Q}\star H.
\]
Class-by-class vanishing:
for class~$G$, all higher shadows vanish
(Theorem~\ref{thm:heisenberg-exact-linearity});
for class~$L$, Jacobi kills the quartic
(Corollary~\ref{cor:affine-postnikov-termination});
for class~$C$, $o_4 = 0$ but $\mathfrak{Q}\neq 0$
(the contact class is well-defined);
for class~$M$, $o_4 \neq 0$ generically
(Theorem~\ref{thm:w-virasoro-quintic-forced}).

\smallskip\noindent
The bar-intrinsic construction
(Theorem~\ref{thm:mc2-bar-intrinsic}) proves $\Theta_\cA$ to all
orders, bypassing the constructive degree-by-degree approach; the present
theorem retains value as a characterization of when the
\emph{constructive} shadow obstruction tower extends at finite order.
\end{proof}

\begin{conjecture}[Singular-fiber descent; \ClaimStatusConjectured]
\label{conj:singular-fiber-descent}
\index{minimal model!singular-fiber descent}
\index{coderived!minimal-model descent}
\textup{[}Type signature:
$(\mathrm{Open}\cap\mathrm{Cat},\mathrm{factorisation},
\mathrm{level}=2\to 4,
\mathrm{hypothesis\ package}=\mathrm{generic\ Virasoro/}\Walg_N\mathrm{\ family\ on\ KSDual\ sublocus};
\mathrm{coderived\ specialisation\ at\ rational\ minimal\text{-}model\ point};
\mathrm{Yamada\ analytic\ closure\ for\ Prochazka\text{-}closed\ }\Walg_N\mathrm{\ data};
\mathrm{Vol\ III\ }(\Sigma_{N-1},C)\text{-specialisation\ at\ stage\text{-}2\ shadow\ at\ rational\ point})$\textup{]}
Rational minimal models are obtained by derived specialization from
the resolved interface family. The parent generic family
\textup{(}generic Virasoro/$\Walg_N$\textup{)} carries the
modular Koszul structure; the rational minimal-model point is a
singular fiber where bar spectral sequences cease to collapse
cleanly and coderived methods become necessary
\textup{(}\S\textup{\ref{subsec:coderived-ran})}.
\end{conjecture}

\begin{remark}[Motivation for singular-fiber descent]
\label{rem:singular-fiber-motivation}
The 2026 paper~\cite{GZ26} works with a resolved/interface parent
geometry. This monograph gives modular Koszul theorems on generic
Virasoro/$\Walg_N$-type families
(Theorem~\ref{thm:algebraic-family-rigidity}), while explicitly
warning that minimal-model quotients are generally not modular Koszul
(Remark~\ref{rem:four-levels}). The correct platonic
theorem therefore lives first on the generic/resolved family, and
the minimal model appears as a singular specialization.
\end{remark}

\begin{remark}[Five targets as typed shadows of $\Theta_\cA$]
\label{rem:five-targets-as-projections}
The five theorems of the algebraic engine are not all literal
projections of a single displayed component. They are typed shadows of
the same bar-intrinsic modular package: Theorem~D is the scalar trace
through~$\kappa$, Theorems~A--B are the bar--cobar and chiral
Positselski comparison attached to the same differential, Theorem~C is
the Verdier-Lagrangian splitting, and Theorem~H is the chiral
Hochschild projection. The holographic targets below are the analogous
restrictions to collision strata, genus sectors, and degree
truncations of~$\gAmod$:
\begin{center}
\begin{tabular}{lll}
\textbf{Target} & \textbf{Stratum of $\gAmod$} & \textbf{Projection}\\
\hline
Boundary--defect & genus~0, full degree
 & typed Verdier/Koszul map $\mathbb H_X(\cA)\to\iota_X(\cA^!)$\\
Yangian--shadow & genus~0, binary/ternary collision & $\operatorname{Res}^{\mathrm{coll}}_{0,2}(\Theta)=r(z)$\\
Sphere reconstruction & genus~0, $n$-point & flat sections of $\nabla^{\mathrm{hol}}_{0,n}$\\
Quartic obstruction & degree~4, all genera & $[\Theta^{\le 3}]_{\mathrm{wt}\,4} = \mathfrak{R}^{\mathrm{mod}}_4$\\
Singular-fiber descent & degeneration locus & coderived specialization of $\Theta$\\
\end{tabular}
\end{center}
\end{remark}

\subsubsection{The factorisation quantum-group envelope}

The categorical home for the full quantum algebra is one level beyond
the dg-shifted Yangian. Latyntsev~\cite{Latyntsev23} develops
factorisation and vertex quantum groups whose representation theory is
controlled by spectral $R$-matrices, and proves equivalences between
vertex quantum groups and factorisation op-quantum groups. The DNP
line category is meromorphic and not genuinely braided; the factorisation
quantum-group formalism supplies the nearest categorical envelope:
\begin{equation}\label{eq:fqg-envelope}
\Ydg_\cA
\;\subset\;
\text{meromorphic HT shadow}
\;\rightsquigarrow\;
\text{factorisation quantum-group envelope}.
\end{equation}
The target theorem should identify the DNP line-operator category
with the non-braided/meromorphic boundary shadow of a factorisation
quantum-group object, recovering braided structure only after
passing to the correct completed or op-dual factorization envelope.
This is the categorical home for the full quantum-algebra claim and
the correct target for
Conjecture~\ref{conj:derived-drinfeld-kohno} and the Bridge
Theorem (Remark~\ref{rem:bridge-theorem-programme}).

\subsubsection{The deformation-quantization bridge}

\begin{conjecture}[HT deformation quantization: geometric bridge; \ClaimStatusConjectured]
\label{conj:ht-deformation-quantization}
\index{deformation quantization!HT holographic}
\index{Poisson vertex algebra!HT quantization}
\textup{[}Type signature:
$(\mathrm{Open}\cap\mathrm{Cat},\mathrm{factorisation},
\mathrm{level}=2\to 3,
\mathrm{hypothesis\ package}=\mathrm{classical\ modular\ coisson\ data};
\mathrm{Penkava\text{--}Vanhaecke\ PVA\text{-}to\text{-}quantum\ bridge\ \textup{(}MA\text{-}12\ endpoint\textup{)}};
\mathrm{Khan\text{--}Zeng\ 3d\ Poisson\ sigma\text{-}model\ realization};
\mathrm{KZ\ analytic\ SDR\ on\ }\cM_{0,n};
\mathrm{KSDual\ sublocus\ for\ Verdier\text{-}consistent\ classical\ }r\text{-}\mathrm{matrix};
\mathrm{Vol\ III\ }(\Sigma_2,C)\text{-specialisation\ at\ }d=3;
\mathrm{classical\text{-}to\text{-}quantum\ promotion\ requires\ named\ analytic\ endpoint:\ not\ guaranteed\ by\ PVA\ Jacobi\ alone})$\textup{]}
There exists a deformation-quantization functor
\begin{equation}\label{eq:ht-deformation-quantization}
Q_{\mathrm{HT}}
\;\colon\;
\{\text{classical modular coisson data}\}
\;\longrightarrow\;
\{\text{quantum modular Koszul data}\}
\end{equation}
with the property that $Q_{\mathrm{HT}}(\text{binary classical
bracket}) = r(z)$ and $Q_{\mathrm{HT}}(\text{full modular classical
class}) = \Theta_\cA$. The classical bulk PVA/$\lambda$-bracket
package is promoted to 3d HT gauge data via the 3d Poisson sigma
model~\textup{\cite{KhanZeng25}}, gauge invariance is equivalent to the Jacobi
identity, and the boundary vertex algebra is the deformation
quantization of the Poisson vertex algebra.
\end{conjecture}

\begin{remark}[Proved formal algebraic content]
\label{rem:ht-deformation-quantization-formal}
The formal algebraic deformation-quantization theory is proved in
Volume~II. The formal structure theorem identifies the genus-$0$
deformation complex; the universal genus-$1$ obstruction class and its
higher-genus recursive extensions are computed; and the
Virasoro/$\mathcal{W}_3$ quantization via Drinfeld--Sokolov reduction
is proved by direct calculation.
The remaining conjectural content of
Conjecture~\textup{\ref{conj:ht-deformation-quantization}} is
the \emph{geometric interpretation}: that $Q_{\mathrm{HT}}$ admits
a 3d Poisson sigma-model realization in the sense
of~\textup{\cite{KhanZeng25}}, and that the formal algebraic functor
extends to an equivalence of modular categories.
\end{remark}

\begin{proposition}[PVA degree constraint and
the inevitability of $2{+}1$ dimensions; \ClaimStatusProvedElsewhere]
\label{prop:pva-degree-constraint}
\index{Khan--Zeng!degree constraint}
\index{dimensional jump!PVA degree constraint}
Let $\cP$ be a freely generated Poisson vertex algebra
on $d$-dimensional base, with $\lambda$-bracket of pole order~$k$.
The mixed holomorphic-topological gauge theory of\/
\textup{\cite{KhanZeng25}} has $(d+k)$-dimensional spacetime with
$d$~holomorphic and $k$~topological directions. For a
$\lambda$-bracket with a single-variable $\lambda$-pole,
the constraint $d + k = d + 1$ forces $k = 1$, yielding spacetime
$\bC^d \times \bR$. For a one-variable PVA \textup{(}$d=1$\textup{)},
the unique target dimension is~$3$:
\[
 \cP \text{ a one-variable PVA}
 \quad\Longrightarrow\quad
 \text{3d HT theory on } \bC \times \bR.
\]
If $\cP$ contains an inner Virasoro element
$T \in \cP$ with $\{T_\lambda a\} = \partial a + (\Delta_a+1)\lambda a$,
then holomorphic translation becomes gauge-exact in the
BRST complex, upgrading the $\mathrm{HT}$ symmetry to a fully
topological symmetry.
\end{proposition}

\begin{proof}
The degree constraint is~\cite[\S\,2.2, equation~(2.24)]{KhanZeng25}:
the action functional has form degree $d + k$ on
$\bC^d \times \bR^k$, and gauge invariance requires the
$\lambda$-bracket to supply exactly $k$ spectral/topological
directions. A single-variable $\lambda$-bracket has $k = 1$.
The Virasoro enhancement
is~\cite[\S\,4.3]{KhanZeng25}: the generator $T$ defines a
gauge transformation $\delta_T A = (d_t + \bar\partial)\varepsilon$
that makes holomorphic translation BRST-exact.
\end{proof}

The bridge operates on three levels:

\emph{Genus-$0$ seed.}
The classical coisson bracket
$\{a_\lambda b\}=\sum_{n\ge 0}\lambda^n c_n$
determines the genus-$0$ component of the bar differential.
Quadratic duality (Gui--Li--Zeng~\cite{GLZ22}) produces
the classical $r$-matrix
$r^{\mathrm{cl}}(z)\in\cA^!\otimes\cA^!((\!( z^{-1})\!))$
satisfying CYBE\@. The quantum deformation
$r(z)=r^{\mathrm{cl}}(z)+\hbar\,r^{(1)}(z)+\cdots$
is the collision residue of~$\Theta_\cA$:
\[
r(z)
\;=\;
\operatorname{Res}^{\mathrm{coll}}_{0,2}(\Theta_\cA).
\]
This is proved
(Theorem~\ref{thm:collision-residue-twisting}).

\emph{Genus-$1$ lift.}
The genus-$1$ obstruction is the BV operator
$\Delta_{\mathrm{cyc}}$
(cf.\ Theorem~\ref{thm:w3-genus1-curvature}).
A genus-$0$ Maurer--Cartan element $\Theta_0$ extends to
genus~$1$ if and only if
$\Delta_{\mathrm{cyc}}(\Theta_0)=0$.
For $\mathcal W_3$: the obstruction vanishes
(Theorem~\ref{thm:w3-genus1-curvature}), and the unique genus-$1$ lift
direction is the central-parameter direction $\partial_c P_c$
(Theorem~\ref{thm:w3-genus1-complementarity}).

\emph{Higher-genus completion.}
The full genus tower is the MC equation
$D\Theta_\cA+\tfrac12[\Theta_\cA,\Theta_\cA]=0$
(Theorem~\ref{thm:mc2-bar-intrinsic}). The genus-$g$
partition function $F_g(\cA)=\kappa\cdot\tfrac{2^{2g-1}-1}{2^{2g-1}}\cdot\tfrac{|B_{2g}|}{(2g)!}$
at the scalar level; the nonlinear corrections are the
quartic contact shadow $\mathfrak Q$, the Hessian
$\delta_H^{(1)}$, and higher planted-forest terms.

The deformation-quantization bridge $Q_{\mathrm{HT}}$ maps:
\begin{center}
\renewcommand{\arraystretch}{1.15}
\begin{tabular}{ll}
\textbf{Classical (PVA)} & \textbf{Quantum (vertex algebra)} \\
\hline
$\lambda$-bracket $\{a_\lambda b\}$
 & chiral bracket $a_{(n)}b$ \\
Poisson bivector $P$
 & genus-$0$ bar differential $d_{\barB}$ \\
classical $r$-matrix
 & dg-shifted Yangian $r(z)$ \\
Jacobi identity
 & $d_{\barB}^2=0$ (Arnold + Borcherds) \\
trivial curvature
 & $d_{\mathrm{fib}}^2=\kappa\cdot\omega_g$ \\
PVA deformation space
 & modular homotopy type $\mathcal M^{\mathrm{mod}}_\cA$
\end{tabular}
\end{center}

\subsubsection{The mature platonic theorem}

\begin{remark}[The mature statement]
\label{rem:mature-platonic-statement}
\index{holographic modular Koszul datum!platonic theorem}
The target theorem is conditional:
\emph{for an HT holographic pair whose boundary comparison functor has
been constructed, the protected boundary sector is controlled by a
holographic modular Koszul datum~$\mathcal{H}(T)$}. Here $\cA$ is the
boundary chiral algebra, $\cA^!$ is the typed Verdier/Koszul branch,
$\mathcal C$ is the bulk/line comparison slot, $r(z)$ is the binary
genus-zero shadow of~$\Theta_\cA$, and $\Theta_\cA$ is the all-genus
modular MC class. The flat sections of
$\nabla^{\mathrm{hol}}_{0,n}$ are protected sphere correlators on the
established comparison surfaces. Higher genus is governed algebraically
by clutching, trace, and Verdier duality; its physical HT interpretation
requires the additional bulk comparison maps.

The proof obligations are:
\begin{enumerate}[label=\textup{(\arabic*)}]
\item Prove theorems on the generic resolved interface family, not
 at the minimal-model quotient. Generic Virasoro/$\Walg_N$ families
 are the first all-genus modular laboratory.
\item Construct $\cA$, $\cA^!$, $\mathcal{C}$ simultaneously as one
 package.
\item Construct $\Theta_\cA^{\leq 4}$ with binary shadow $r(z)$,
 trivalent shadow yielding $A_\infty$-YBE, and quartic obstruction
 $R^{\mathrm{mod}}_{4,g,n}(\cA)$. \textup{(}The full
 $\Theta_\cA$ is already proved by
 Theorem~\ref{thm:mc2-bar-intrinsic}; the constructive approach here
 is computational refinement extracting explicit finite-order
 shadows.\textup{)}
\item Show that the commuting sphere Hamiltonians in~\textup{\cite{GZ26}}
 are the scalar images of $\operatorname{Sh}_{0,n}(\Theta_\cA)$.
\item Pass to genus~$1$: the first higher-genus holographic check is
 the identification of genus-$1$ curvature/anomaly with the first
 modular correction to the sphere shadow.
\item Descend to rational minimal-model points by coderived
 specialization.
\end{enumerate}
The 2026 non-planar sphere algebra is therefore a comparison surface
for the binary shadow of a modular Maurer--Cartan class, not a
replacement for the all-genus problem.
\end{remark}


%% ====================================================================
%% THREE STRUCTURAL TARGET THEOREMS
%% ====================================================================
\subsection{Three structural target theorems}%
\label{subsec:three-target-theorems}%
\index{target theorems!structural}

\begin{conjecture}[Twisted-pair deformation enhancement;
\ClaimStatusConjectured]%
\label{conj:twisted-pair-deformation}%
\index{twisted pair!deformation enhancement}%
\textup{[}Type signature:
$(\mathrm{Open}\cap\mathrm{Cat},\mathrm{factorisation},
\mathrm{level}=2\to 3,
\mathrm{hypothesis\ package}=\mathrm{dualizable\ effective\ chiral\ QLS\ data};
\mathrm{curved\ }L_\infty\mathrm{\text{-}envelope\ of\ MC\ equation};
\mathrm{KSDual\ sublocus\ for\ self\text{-}adjoint\ basepoint\ }S=\Theta_\cA)$\textup{]}
For dualizable effective chiral QLS data, the functor of
deformations of chiral morphisms out of~$\cA$ is
represented by the twisted Maurer--Cartan functor of its
dual twisted pair\textup:
\begin{equation}\label{eq:twisted-pair-mc}
\mathrm{Def}_{\mathrm{ch}}(\cA \to \cC)
\;\simeq\;
\MC_{T_\cA}(\cC)
\;:=\;
\bigl\{\alpha \in \Gamma(X,\cB\otimes\cC)^{-1}
\;\big|\;
\mu\bigl((S{+}\alpha)\boxtimes(S{+}\alpha)\bigr) = 0
\bigr\}/{\sim},
\end{equation}
where $S \in \mathrm{MC}(\gAmod)$ is the basepoint MC element
(the universal shadow class $\Theta_\cA$ restricted to the given
deformation),
as an equivalence of deformation functors \textup{(}not
merely a set-level bijection\textup{)}. The enhancement
requires a curved $L_\infty$-envelope of the twisted-pair
MC equation.
\end{conjecture}

\begin{conjecture}[Lattice Fourier--Koszul identification;
\ClaimStatusConjectured]%
\label{conj:lattice-fourier-koszul}%
\index{lattice VOA!Fourier--Koszul}%
\index{Fourier--Mukai!Koszul identification}%
\textup{[}Type signature:
$(\mathrm{Open}\cap\mathrm{Cat},\mathrm{factorisation},
\mathrm{level}=2\to 5,
\mathrm{hypothesis\ package}=\mathrm{lattice\ chiral\ algebra}\ \cA_\theta\ \mathrm{with}\ \theta=(\kappa,\lambda,c);
\mathrm{Fourier\text{--}Mukai\ Picard\ data};
\mathrm{KSDual\ sublocus\ on\ self\text{-}dual\ even\ lattices};
\mathrm{modularity\ via\ trace\ +\ clutching\ on\ }\overline{\cM}_{1,n},\ \mathrm{not\ a\ closed\text{-}lattice\text{-}VOA\ property})$\textup{]}
For a lattice chiral algebra~$\cA_\theta$ with
$\theta$-datum~$\theta = (\kappa, \lambda, c)$
\textup{(}Definition~\textup{\ref{def:analytic-theta-datum}}\textup{)},
the holographic modular Koszul datum
$\mathcal{H}(\cA_\theta)$ determines~$\theta$
functorially, and the genus-$1$ categorical trace of
the line category is naturally identified with the
Fourier--Mukai transform of the corresponding
Picard-line-bundle data. The lattice sector is where
the three worlds (local Heisenberg, global
theta/Fourier, and dual/MC algebra) first meet in a
single proved theorem.
\end{conjecture}

%% ====================================================================
%% PROVED CONSTRUCTIONS FROM THE COLLISION FILTRATION
%% ====================================================================
\subsection{The collision filtration and proved recovery theorems}
\label{subsec:collision-filtration-recovery}
\index{collision filtration|textbf}
\index{collision residue!proved recovery theorems}

The core identifications of the holographic modular Koszul datum (the collision residue recovering
the Yangian, the affine genus-$0$ shadow/KZ comparison on the
affine Kac--Moody comparison surface, the quartic obstruction governing
extensions) are formal consequences
of the proved core. This subsection extracts them.

\subsubsection{The collision filtration}

\begin{definition}[Collision filtration on the genus-zero convolution algebra]
\label{def:collision-filtration}
\index{collision filtration!definition}
Let $\gAmod = \prod_{2g-2+n>0}
\Hom_{\Sigma_n}(C_*(\overline{\mathcal{M}}_{g,n}),
\operatorname{End}_\cA(n))$
be the modular convolution dg~Lie algebra
(Definition~\ref{def:modular-convolution-dg-lie}).
Define the \emph{collision filtration} on the genus-zero sector
$\mathfrak{g}_{\cA}^{\mathrm{mod},\,g{=}0} = \prod_{n \geq 3}
\Hom_{\Sigma_n}(C_*(\overline{\mathcal{M}}_{0,n}),
\operatorname{End}_\cA(n))$ by
\begin{equation}\label{eq:collision-filtration}
F^d_{\mathrm{coll}}\, \mathfrak{g}_{\cA}^{\mathrm{mod},\,g{=}0}
\;:=\;
\prod_{n \geq 3}
\Hom_{\Sigma_n}\bigl(
 C_*(\overline{\mathcal{M}}_{0,n}^{\geq d}),\;
 \operatorname{End}_\cA(n)\bigr),
\end{equation}
where $\overline{\mathcal{M}}_{0,n}^{\geq d}$ denotes the union of
boundary strata of codimension $\geq d$ in
$\overline{\mathcal{M}}_{0,n}$. The associated graded
$\operatorname{gr}^d_{\mathrm{coll}}$ is supported on
strata of \emph{exactly} codimension~$d$: dual trees with
exactly $d$~internal edges, corresponding to exactly
$d$~collision events.
\end{definition}

\begin{remark}[Collision depth vs.\ degree weight]
\label{rem:collision-vs-degree}
The collision filtration is \emph{not} the degree/weight filtration of
the shadow obstruction tower. Degree controls the order of the shadow
truncation~$\Theta_\cA^{\leq r}$; collision depth controls the
\emph{boundary stratum depth} within a fixed $n$-point moduli space.
At collision depth~$0$ (the open stratum $\mathcal{M}_{0,n}$), the
datum is the ``bulk'' OPE with all points spread out.
At $d = 1$ (codimension-$1$ boundary), binary collisions produce the
$r$-matrix. At $d = 2$, ternary collisions produce Yang--Baxter.
\end{remark}

\subsubsection{The collision residue theorem}

\begin{theorem}[Collision residue = universal twisting morphism;
\ClaimStatusConditional]
\label{thm:collision-residue-twisting}
\index{collision residue!equals twisting morphism|textbf}
\index{twisting morphism!from collision residue}
% Type signature follows below within the theorem body
\textup{[}Type signature:
$(\mathrm{Open}\cap\mathrm{Cat},\mathrm{factorisation},
\mathrm{level}=2\to 4,
\mathrm{hypothesis\ package}=\mathrm{modular\ Koszul};
\mathrm{Verdier\ window\ for\ part\ (ii)};\mathrm{KSDual\ sublocus\ if\ }
r_\cA(z)\mathrm{\ is\ to\ be\ self\text{-}adjoint\ in\ the\ Verdier\ pairing})$\textup{]}
Let $\cA$ be a modular Koszul chiral algebra on a smooth projective
curve~$X$, and let
$\Theta_\cA = D_\cA - d_0 \in
\mathrm{MC}(\Defcyc(\cA) \,\widehat{\otimes}\, \mathcal{G}_{\mathrm{mod}})$
be the bar-intrinsic MC element
\textup{(}Theorem~\textup{\ref{thm:mc2-bar-intrinsic})}. Then:
\begin{enumerate}[label=\textup{(\roman*)}]
\item The restriction of\/ $\Theta_\cA$ to the genus-zero, two-point
 collision stratum recovers the universal twisting morphism:
 \begin{equation}\label{eq:collision-residue-equals-twisting}
 \operatorname{Res}^{\mathrm{coll}}_{0,2}(\Theta_\cA)
 \;=\;
 \pi_\cA
 \;\in\;
 \mathrm{Tw}(\barBch(\cA), \cA),
 \end{equation}
 where $\pi_\cA\colon \barBch(\cA) \to \cA$ is the universal
 twisting morphism
 \textup{(}Corollary~\textup{\ref{cor:three-bijections})}, realized
 geometrically by propagator extraction
 \textup{(}Proposition~\textup{\ref{prop:twisting-morphism-propagator})}:
 \begin{equation}\label{eq:twisting-propagator-frontier}
 \pi_\cA([s^{-1}\bar{a}])(z)
 \;=\;
 \operatorname{Res}_{w = z}\,
 \bar{a}(w) \cdot \eta(z,w),
 \end{equation}
 with $\eta(z,w) = d\log(z - w)$ the logarithmic propagator.
\item After passing to the intrinsic bar-dual coalgebra
 $\cA^i := H^*(\barBch(\cA))$ and, in a finite Verdier window, to
 coordinates in $(\cA^!)^\vee$, the twisting morphism produces the
 spectral $r$-matrix
 $r_\cA(z)$
 of Definition~\textup{\ref{def:dg-shifted-yangian}}.
\end{enumerate}
\end{theorem}

\begin{proof}
The bar-intrinsic MC element is $\Theta_\cA = D_\cA - d_0$.
At genus zero, $d_\cA^{(0)} = d_0$ is the standard bar differential on
$\overline{C}_n(X)$. The collision residue along the diagonal
$D_{ij} = \{z_i = z_j\} \subset \overline{C}_n(X)$ extracts the
leading pole of the bar differential, which is the chiral bracket
$\mu^{\mathrm{ch}}(a_i, a_j)$ evaluated against the propagator
$\eta(z_i, z_j)$. This is exactly the geometric realization of the
universal twisting morphism
(Proposition~\ref{prop:twisting-morphism-propagator}).
For~(ii), the three bijections
(Corollary~\ref{cor:three-bijections}) give the canonical identification
with the $r$-matrix on the evaluation locus.
\end{proof}

\begin{theorem}[MC at collision depth~$2$ gives $A_\infty$-YBE;
\ClaimStatusProvedHere]
\label{thm:collision-depth-2-ybe}
\index{Yang--Baxter equation!from collision depth 2}
\index{Arnold relation!gives Yang--Baxter}
\textup{[}Type signature:
$(\mathrm{Open}\cap\mathrm{Chain},\mathrm{chain\text{-}level},
\mathrm{level}=2\to 4,
\mathrm{hypothesis\ package}=\Theta_\cA\mathrm{\ \textup{(}Theorem~\ref{thm:mc2-bar-intrinsic}\textup{)}};
\mathrm{Arnold\ relation\ on\ }\mathrm{Conf}_4(X);
\mathrm{KSDual\ sublocus\ for\ self\text{-}adjoint\ }r(z))$\textup{]}
The MC equation restricted to the genus-zero four-point sector
$\overline{\mathcal{M}}_{0,4}$ and projected to collision depth~$2$
gives the infinitesimal braid relation \textup{(}IBR\textup{)} for
$r_\cA(z) = \Omega/(k{+}h^\vee) z$:
\begin{equation}\label{eq:ibr-from-collision}
[\Omega_{ij},\; \Omega_{ik} + \Omega_{jk}] \;=\; 0
\qquad \text{for all distinct } i,j,k.
\end{equation}
This is equivalent to the classical Yang--Baxter equation for the
spectral-parameter $r$-matrix via partial fractions on
$\overline{\mathcal{M}}_{0,4}$.
\end{theorem}

\begin{proof}
The boundary $\partial\overline{\mathcal{M}}_{0,4}$ has three divisors
$D_s, D_t, D_u$, each isomorphic to
$\overline{\mathcal{M}}_{0,3} \times
\overline{\mathcal{M}}_{0,3} = \mathrm{pt}$.
The MC equation at four points, evaluated at triple collision, reads
$d(\Theta_{0,4}) + \sum_{\text{channels}}
[\Theta_{0,3}, \Theta_{0,3}]_{\text{channel}} = 0$.
After collision residue extraction, the three channel contributions
give the IBR\@. The mechanism is the Arnold relation
(Theorem~\ref{thm:arnold-relations})
$\eta_{12} \wedge \eta_{23} + \eta_{23} \wedge \eta_{31} +
\eta_{31} \wedge \eta_{12} = 0$
tensored with the structure-constant data of~$\cA$. The equivalence
between $d_{\mathrm{bar}}^2 = 0$ and the Arnold relation
(Theorem~\ref{thm:bar-nilpotency-complete}) completes the proof.
\end{proof}

\subsubsection{Recovery of classical connections}

\begin{theorem}[Shadow/KZ comparison on the affine Kac--Moody surface;
\ClaimStatusConditional]
\label{thm:shadow-connection-kz}
\index{Knizhnik--Zamolodchikov connection!from shadow connection|textbf}
\index{shadow connection!affine Kac--Moody comparison}
\textup{[}Type signature:
$(\mathrm{Open}\cap\mathrm{Cat},\mathrm{factorisation},
\mathrm{level}=4\to 5,
\mathrm{hypothesis\ package}=\cA=\widehat{\fg}_k,\ k\ne -h^\vee;
\mathrm{chiral\ Ward\ identity\ on\ }\cM_{0,n};
\mathrm{KSDual\ sublocus\ for\ self\text{-}adjoint\ }r(z);
\mathrm{Vol\ III\ }(\Sigma_{\mathrm{rank}\,\fg-1},C)\text{-specialisation\ at\ stage\text{-}2})$\textup{]}
Let $\cA = \widehat{\fg}_k$ at level $k \neq -h^\vee$, with
$\kappa(\widehat{\fg}_k) = \dim(\fg)(k{+}h^\vee)/(2h^\vee)$.
On the affine Kac--Moody comparison surface, the genus-$0$ shadow
connection on $n$-point conformal blocks
$\mathcal{V}_{0,n}(\widehat{\fg}_k; V_1, \dotsc, V_n)$ identifies with the
Knizhnik--Zamolodchikov connection:
\begin{equation}\label{eq:shadow-equals-kz}
\nabla^{\mathrm{hol}}_{0,n}
\;=\;
d \;-\; \frac{1}{k + h^\vee}\,
\sum_{1 \leq i < j \leq n}
\Omega_{ij}\, d(z_i - z_j)/(z_i - z_j)
\;=\; \nabla^{\mathrm{KZ}},
\end{equation}
where $\Omega_{ij} = \sum_a J^a_i \otimes J^a_j$ is the Casimir.
In particular, on this affine comparison surface, flatness
$(\nabla^{\mathrm{hol}}_{0,n})^2 = 0$ is the KZ flatness equation, and
the corresponding conformal blocks are horizontal sections.
\end{theorem}

\begin{proof}
By Theorem~\ref{thm:collision-residue-twisting}, the collision residue
of~$\Theta_\cA$ recovers the twisting morphism $\pi_\cA$, which for
$\widehat{\fg}_k$ extracts the chiral bracket
$\mu^{\mathrm{ch}}(J^a(z), J^b(w)) =
[J^a, J^b](w)/(z{-}w) + k(a,b)/(z{-}w)^2$.
The shadow representation on conformal blocks produces
$\operatorname{Sh}_{0,n}(\Theta_\cA) =
(k{+}h^\vee)^{-1} \sum_{i<j} \Omega_{ij}\, d(z_i{-}z_j)/(z_i{-}z_j)$,
where $(k{+}h^\vee)^{-1}$ is the propagator normalization
against the Killing form.
The connection $\nabla^{\mathrm{hol}}_{0,n} =
d - \operatorname{Sh}_{0,n}(\Theta_\cA)$ therefore identifies on this
affine comparison surface with the KZ connection. Flatness follows from the MC equation
(Theorem~\ref{thm:collision-depth-2-ybe}).
The bar-complex comparison with KZ
is also established in Proposition~\ref{prop:kz-from-bar}.
\end{proof}

\begin{corollary}[Heisenberg shadow connection;
\ClaimStatusProvedElsewhere]
\label{cor:shadow-connection-heisenberg}
\index{Heisenberg algebra!shadow connection}
For rank-$1$ Heisenberg $\mathcal{H}_\kappa$ \textup{(}class~$G$,
tower terminates at degree~$2$;
Corollary~\textup{\ref{cor:heisenberg-postnikov-termination})}, the
shadow connection is
\begin{equation}\label{eq:heisenberg-shadow-connection}
\nabla^{\mathrm{hol}}_{0,n}
=
d - \kappa
\sum_{i < j} q_i q_j\, d(z_i - z_j)/(z_i - z_j),
\end{equation}
with flat sections
$\langle \phi_{q_1}(z_1) \dotsm \phi_{q_n}(z_n) \rangle
= \delta_{\sum q_i,0}\,
\prod_{i<j} (z_i - z_j)^{\kappa q_i q_j}$.
Flatness is automatic: the Heisenberg is abelian.
\end{corollary}

\begin{proposition}[Shadow connection for Virasoro and BPZ on the degenerate-representation surface;
\ClaimStatusProvedHere]
\label{prop:shadow-connection-bpz}
\index{BPZ equations!from higher collision residues}
For $\cA = \mathrm{Vir}_c$, the genus-$0$ shadow connection on
$n$-point blocks of primaries with conformal weights
$h_1, \dotsc, h_n$ acts as
\[
\nabla^{\mathrm{hol}}_{0,n}
=
d -
\sum_{i=1}^n \left(
 \sum_{j \neq i} \frac{h_j}{(z_i - z_j)^2}
 +
 \frac{\partial_{z_j}}{z_i - z_j}
\right) dz_i.
\]
The BPZ equations for degenerate representations are the conditions
that higher collision residues
$\operatorname{Res}^{\mathrm{coll}}_{0,k}(\Theta_{\mathrm{Vir}})$
for $k \geq 3$ (encoding the infinite shadow obstruction tower of class~$M$) produce
null-vector differential operators acting on the flat sections.
\end{proposition}

\begin{proof}
The Virasoro OPE gives the collision residue. The shadow
representation on primaries extracts the conformal Ward identity.
The infinite tower
(Theorem~\ref{thm:w-virasoro-mixed-shadow};
$r_{\max} = \infty$, Theorem~\ref{thm:w-virasoro-quintic-forced})
produces higher-order differential constraints that for degenerate
representations truncate to the finite-order BPZ null-vector equations.
\end{proof}

\subsubsection{The quartic obstruction theorem}

\begin{theorem}[Quartic obstruction = $L_\infty$ obstruction;
\ClaimStatusConditional]
\label{thm:quartic-obstruction-linf}
\index{quartic resonance class!equals L-infinity obstruction}
\textup{[}Type signature:
$(\mathrm{Open}\cap\mathrm{Chain},\mathrm{chain\text{-}level},
\mathrm{level}=3\to 4,
\mathrm{hypothesis\ package}=\mathrm{constructive\ shadow\ tower};
\mathrm{five\text{-}archetype\ classification\ G/L/C/M};
\mathrm{class\text{-}M\ vanishing\ at\ }o_4\ \mathrm{requires\ KZ\ analytic\ SDR\ +\ weight\text{-}completed\ ambient};
\mathrm{KSDual\ sublocus\ for\ Verdier\text{-}consistent\ obstruction\ class})$\textup{]}
The obstruction to extending $\Theta_\cA^{\leq 3}$ to
$\Theta_\cA^{\leq 4}$ is
\begin{equation}\label{eq:quartic-obstruction-class}
[o_4(\cA)]
=
\bigl[
 D_\cA \Theta_\cA^{\leq 3}
 + \tfrac{1}{2}[\Theta_\cA^{\leq 3},\,
 \Theta_\cA^{\leq 3}]
\bigr]_{\mathrm{wt}\,4}
\in
H^2(\cA^{\mathrm{sh}}_{4,0}),
\end{equation}
canonically identified with the modular quartic resonance class
$\mathfrak{R}^{\mathrm{mod}}_{4,\bullet,\bullet}(\cA)$
\textup{(}Definition~\textup{\ref{def:nms-modular-quartic-resonance-class})}.
In particular: $\mathfrak{R}^{\mathrm{mod}}_4 = 0$ if and only if\/
$\Theta^{\leq 3}$ extends to $\Theta^{\leq 4}$.
For class~$G,L$: $o_4 = 0$ automatically.
For class~$C$: $o_4 = 0$ but $\mathfrak{Q} \neq 0$.
For class~$M$: $o_4 \neq 0$ generically; the bar-intrinsic
construction bypasses the tower entirely
\textup{(}Theorem~\textup{\ref{thm:recursive-existence})}.
\end{theorem}

\begin{proof}
Standard $L_\infty$ obstruction theory
(Kontsevich~\cite{Kontsevich97}, Manetti~\cite{Manetti99}) applied to
$\gAmod$ with the weight filtration: the obstruction at weight~$r{+}1$
is the class of the MC failure projected to weight~$r{+}1$. At
weight~$4$, the clutching formula
$\xi^*\mathfrak{Q} = H \star \mathfrak{Q} +
\mathfrak{C} \star \mathfrak{C} + \mathfrak{Q} \star H$
identifies this with
$\mathfrak{R}^{\mathrm{mod}}_{4,g,n}(\cA)$. The class-by-class
vanishing follows from
Theorem~\ref{thm:heisenberg-exact-linearity} (class~$G$),
Corollary~\ref{cor:affine-postnikov-termination} (class~$L$), and
Theorem~\ref{thm:w-virasoro-quintic-forced} (class~$M$).
\end{proof}

\subsubsection{The collision-filtration spectral sequence}

\begin{construction}[Collision-filtration spectral sequence]
\label{constr:holographic-spectral-sequence}
\index{spectral sequence!collision filtration|textbf}
The collision filtration
(Definition~\ref{def:collision-filtration}) induces a spectral
sequence:
\begin{equation}\label{eq:holographic-spectral-sequence}
E_1^{d,q}
=
H^{d+q}\!\bigl(
 \operatorname{gr}^d_{\mathrm{coll}}\, \mathfrak{g}_{\cA}^{\mathrm{mod},\,g{=}0},\,
 d_{\mathrm{int}}\bigr)
\Longrightarrow
H^{d+q}(\mathfrak{g}_{\cA}^{\mathrm{mod},\,g{=}0},\, D).
\end{equation}
The $E_1$~page decomposes by collision depth:
$E_1^{0,*}$ = bulk data on $\mathcal{M}_{0,n}$;
$E_1^{1,*}$ = binary collision = tangent complex of the Yangian;
$E_1^{2,*}$ = ternary collision = YBE deformation complex;
$E_1^{d,*}$ for $d \geq 3$ = higher nonlinear shadows.
The $d_1$ differential propagates binary $\to$ ternary $\to$
quartic via clutching. Collapse classifies shadow depth:
\begin{itemize}
\item Class~$G$: $E_1$ collapses at $d = 0$.
\item Class~$L$: $E_2$ collapses at $d = 1$
 \textup{(}Jacobi kills quartic\textup{)}.
\item Class~$C$: $E_3$ collapses at $d = 2$.
\item Class~$M$: no finite collapse.
\end{itemize}
At genus~$g \geq 1$, the Hodge curvature $\kappa \cdot \omega_g$
deforms the genus-$0$ spectral sequence. For affine KM, the genus-$1$
deformation produces the KZB connection. The first non-trivial
higher-genus differential is controlled by the Hessian correction
$\delta H^{(1)}_\cA$.
\end{construction}

\begin{computation}[The $E_1$ page for $\widehat{\mathfrak{sl}}_2$;
\ClaimStatusConditional]
\label{comp:holographic-ss-sl2}
\index{spectral sequence!holographic!sl2@$\mathfrak{sl}_2$}
For $\widehat{\mathfrak{sl}}_2$ at level~$k$, $n=4$ insertions:
\[
E_1^{0,*}
\;=\;
H^*\!\bigl(\mathfrak{sl}_2^{\otimes 4},\,d_{\mathrm{CE}}\bigr),
\qquad
E_1^{1,*}
\;=\;
\bigoplus_{i<j}
H^*\!\bigl(\mathfrak{sl}_2^{\otimes 3},\,d_{\mathrm{CE}}\bigr)
\;\otimes\;
[\delta_{ij}],
\]
\[
E_1^{2,*}
\;=\;
\bigoplus_{i<j<k}
H^*\!\bigl(\mathfrak{sl}_2^{\otimes 2},\,d_{\mathrm{CE}}\bigr)
\;\otimes\;
[\delta_{ijk}].
\]
The $d_1$ differential $E_1^{0,*}\to E_1^{1,*}$ is the binary
OPE residue: for a Chevalley--Eilenberg cocycle
$\phi\in C^*(\mathfrak{sl}_2^{\otimes 4})$, the image is the
sum over binary collisions
\[
d_1(\phi)
\;=\;
\sum_{i<j}
\operatorname{Res}_{z_i=z_j}\,
\phi(J^a(z_i),\,J^b(z_j),\,\cdots)
\;\cdot\;
[\delta_{ij}].
\]
At the $E_2$ page:
$E_2^{0,*}=\ker d_1$,
$E_2^{1,*}=\operatorname{coker}(d_1|_{E_1^0})
\cap\ker(d_1|_{E_1^1})$.

For $\widehat{\mathfrak{sl}}_2$, strict formality
(Theorem~\ref{thm:km-strictification}) implies
$d_1\colon E_1^{1,*}\to E_1^{2,*}$ is zero:
the ternary collision produces a class in
$E_1^{2,*}$ that is killed by the Jacobi identity on
$\mathfrak{sl}_2$. The spectral sequence collapses at~$E_2$:
$E_\infty = E_2$. Shadow depth~$3$ (Lie/tree class~L).
\end{computation}

\begin{computation}[The $E_1$ page for $\mathrm{Vir}_c$;
\ClaimStatusProvedHere]
\label{comp:holographic-ss-vir}
\index{spectral sequence!holographic!Virasoro}
For $\mathrm{Vir}_c$, $n=4$:
\[
E_1^{0,*}
\;=\;
H^*\!\bigl(L_0^{\otimes 4},\,d_{\mathrm{Vir}}\bigr),
\qquad
E_1^{1,*}
\;=\;
\bigoplus_{i<j}
H^*\!\bigl(L_0^{\otimes 3}\bigr)
\;\otimes\;
[\delta_{ij}],
\]
where $L_0$ denotes the relevant conformal-weight-$0$
sector. The $d_1$ differential maps binary to ternary
collisions via the Virasoro OPE:
\[
d_1\colon
E_1^{1,*}\;\longrightarrow\;E_1^{2,*}
\]
is \emph{nonzero}: the quartic OPE contact term
$\Theta_4^{\mathrm{contact}}$ survives as a class in
$E_2^{2,*}$. Its cohomology class is the quartic contact
invariant:
\[
[d_1]_{E_2^{2,*}}
\;=\;
\frac{10}{c(5c+22)}\,[\Lambda]
\;\neq\;0
\quad(\text{for }c\neq 0,\,-22/5).
\]
The spectral sequence does not collapse at any finite page.
At each page $E_r$, a new obstruction class $o^{(r+2)}_{\mathrm{Vir}}$
appears as the $d_r$ differential. The quintic obstruction
$o^{(5)}\neq 0$ is the $d_3$ differential; its
non-vanishing was computed through
Theorem~\ref{thm:nms-virasoro-quintic-forced}.
Shadow depth~$\infty$ (mixed class~M).
\end{computation}

\begin{computation}[The $E_1$ page for $\beta\gamma$;
\ClaimStatusProvedHere]
\label{comp:holographic-ss-betagamma}
\index{spectral sequence!holographic!beta-gamma@$\beta\gamma$}
For $\beta\gamma$ at $n=4$: the $d_1$ differential
$E_1^{0,*}\to E_1^{1,*}$ is zero. The simple pole
$\beta_{(0)}\gamma=1$ is the vacuum scalar, so the
field-valued binary bracket is zero. The $d_1$ differential
$E_1^{1,*}\to E_1^{2,*}$ is zero (no cubic bracket).
But $E_1^{2,*}$ is nonzero: the quartic contact term
$\Theta_4^{\mathrm{contact}}\neq 0$ appears as a
\emph{permanent cycle} in $E_2^{2,*}$. However,
$\mu_{\beta\gamma}=0$
(Corollary~\ref{cor:nms-betagamma-mu-vanishing}): the
quartic class is a coboundary in $E_3$.
The spectral sequence collapses at~$E_3$: shadow depth~$4$
(contact class~C).

The $\beta\gamma$ system is the unique example where
the shadow depth ($4$) exceeds the Lie depth ($1$, since the
OPE is abelian): the contact quartic is a purely
weight-changing phenomenon, not a Lie-theoretic one.
\end{computation}

\begin{computation}[The $E_1$ page for $\mathcal W_3$;
\ClaimStatusProvedHere]
\label{comp:holographic-ss-w3}
\index{spectral sequence!holographic!W3@$\mathcal W_3$}
For $\mathcal W_3$ at $n=4$: both the cubic and quartic
differentials are nonzero. The $E_1^{1,*}$ term carries the
cubic bracket $[\ell_3^{(0)}]$ from the $TW$-OPE, and the
$E_1^{2,*}$ term carries the quartic contact
$\frac{16}{22+5c}\,\Lambda$.

At the $E_2$ page: the quartic class survives because
$\ell_3^{(0)}\neq 0$ does not kill it (the cubic and quartic
obstructions are independent). The $d_2$ differential
propagates the quartic to the quintic, producing a class in
$E_2^{3,*}$ controlled by the $WW$-OPE at sixth-order pole.

The spectral sequence for $\mathcal W_3$ carries more data
than in the Virasoro case: the deformation space is $2$-dimensional
($T$ and $W$ directions), the genus-$1$ Hessian is a
$2\times 2$ matrix
(Computation~\ref{comp:w3-genus1-hessian}),
and the $d_r$ differentials at each page carry $2\times 2$
matrix data. The spectral discriminant
$\Delta_{\mathcal W_3}(x)$ has two zeros, reflecting the
two eigenvalues of the Hessian.
Shadow depth~$\infty$ (mixed class~M, rank~$2$).
\end{computation}

\subsubsection{The five shadows of $\Theta_\cA$}

\begin{remark}[One object, five projections]

\label{rem:five-shadows-synthesis}
\index{five shadows|textbf}
The single object $\Theta_\cA \in \mathrm{MC}(\gAmod)$ controls the
holographic datum through five proved projections:
\begin{center}
\renewcommand{\arraystretch}{1.3}
\begin{tabular}{@{}clll@{}}
\toprule
\textbf{Shadow} & \textbf{Projection} & \textbf{Classical object}
 & \textbf{Status} \\
\midrule
Scalar $\kappa$ & $\operatorname{tr}(\Theta_\cA)$
 & Genus tower $\kappa\lambda_g$
 & \textbf{Proved} (Thm~D) \\
$r$-matrix & $\operatorname{Res}^{\mathrm{coll}}_{0,2}(\Theta_\cA)$
 & Yangian/$r$-matrix
 & \textbf{Proved} (Thm~\ref{thm:collision-residue-twisting}) \\
Connection & $d - \operatorname{Sh}_{g,n}(\Theta_\cA)$
 & KZ on affine KM comparison surfaces,
 BPZ on Virasoro degenerate-representation surfaces
 & \textbf{Proved} at $g{=}0$ \\
Quartic $R^{\mathrm{mod}}_4$ & $\pi_4(\Theta_\cA)$
 & Obstruction class
 & \textbf{Proved}
 (Thm~\ref{thm:quartic-obstruction-linf}) \\
Clutching & $\xi^*(\Theta_\cA) = \Theta_\cA \star \Theta_\cA$
 & Modular operadic
 & \textbf{Proved} (Thm~\ref{thm:universal-MC}) \\
\bottomrule
\end{tabular}
\end{center}
\emph{One object, one equation, five projections, everything else is
a corollary.} The scalar shadow is Theorem~D\@. The $r$-matrix
shadow is the Yangian. The connection shadow is the modular
shadow connection, with shadow/KZ comparison on the affine Kac--Moody
comparison surfaces and BPZ on the Virasoro degenerate-representation
comparison surfaces. The
quartic shadow is the first nonlinear obstruction. The clutching
shadow is the modular operad structure. Each of the five classical
theories (genus expansion, quantum groups, conformal-block
comparison packages,
deformation obstruction, sewing) is a single projection of the same
MC element.

In particular, $\mathcal{H}(T)$ is not a list of independent objects.
The MC entry controls the shadow tower, while the bar-dual coalgebra,
Verdier line algebra, derived centre, collision residue, and connection
are the corresponding bar, Verdier, Hochschild, residue, and connection
constructions attached to~$\cA$.
\end{remark}

\subsubsection{Complete Heisenberg datum}

\begin{computation}[Complete holographic datum for Heisenberg;
\ClaimStatusProvedHere]
\label{comp:heisenberg-holographic-datum}
\index{Heisenberg algebra!complete holographic datum}
For rank-$1$ Heisenberg $\mathcal{H}_\kappa$ at level $\kappa > 0$:
\begin{enumerate}[label=\textup{(\roman*)}]
\item $\cA = \mathcal{H}_\kappa$; $J(z)J(w) \sim \kappa/(z{-}w)^2$.
\item $\cA^i=H^*(B^{\mathrm{ch}}(\mathcal H_\kappa))$ is the
 bar-dual coalgebra of the finite Heisenberg Fock model.
\item $\cA^! = \operatorname{Sym}^{\mathrm{ch}}(V^*)$ after
 finite-piece Verdier duality, with
 $\kappa(\cA^!) = -\kappa$
 \textup{(}same modular characteristic as $\mathcal{H}_{-\kappa}$,
 but a different algebra; the Koszul dual and the negative-level
 algebra are categorically distinct\textup{)}.
\item $\mathcal{C} = \mathrm{Vect}^{\mathbb{Z}}$; no non-trivial lines
 beyond Fock modules.
\item $r_{\mathcal{H}}(z) = \kappa/z$ \textup{(}scalar; collision residue of the $d\log$ kernel\textup{)}.
\item $\Theta_{\mathcal{H}} = \kappa \cdot \eta \otimes \Lambda$;
 tower terminates at degree~$2$
 \textup{(}Corollary~\textup{\ref{cor:heisenberg-postnikov-termination})}.
\item $\nabla^{\mathrm{hol}}_{0,n} =
 d - \kappa \sum_{i<j} q_i q_j\, d(z_i{-}z_j)/(z_i{-}z_j)$;
 flat sections $= \prod_{i<j}(z_i{-}z_j)^{\kappa q_i q_j}$.
\end{enumerate}
The finite Heisenberg Fock model determines all entries from
$\kappa$ and the chosen rank-one lattice. The spectral sequence
collapses at~$E_1$.
\end{computation}

\begin{computation}[Holographic datum for affine Kac--Moody;
\ClaimStatusProvedHere]
\label{comp:affine-holographic-datum}
\index{Kac--Moody!holographic datum}
For $\widehat{\mathfrak g}_k$ at non-critical level $k\neq -h^\vee$:
\begin{enumerate}[label=\textup{(\roman*)}]
\item $\cA = \widehat{\mathfrak g}_k$.
\item $\cA^! = \widehat{\mathfrak g}_{-k-2h^\vee}$
 (Feigin--Frenkel duality).
\item $\cC \simeq \cO^{\mathrm{sh}}_q(\widehat{\mathfrak g})$,
 $q = e^{i\pi/(k+h^\vee)}$.
\item $r(z) = k\,\Omega_{\mathfrak g}/z$, the Casimir $r$-matrix
 with $\Omega = \sum e_a \otimes e^a$.
\item $\Theta_k$ has shadow depth~$3$ (Lie/tree archetype);
 $\kappa(\widehat{\mathfrak g}_k)
 = \dim\mathfrak g\cdot(k+h^\vee)/(2h^\vee)$.
\item $\nabla^{\mathrm{hol}}_{0,n}$ is the Knizhnik--Zamolodchikov
 connection
 $d - \frac{1}{k+h^\vee}\sum_{i<j}\Omega_{ij}\,d(z_i{-}z_j)/(z_i{-}z_j)$.
\end{enumerate}
Complementarity: $\kappa_k + \kappa_{-k-2h^\vee} = 0$.
At $k = -h^\vee$: the center
$\mathfrak z(\widehat{\mathfrak g}_{-h^\vee})$ is the
Feigin--Frenkel algebra of $\mathfrak g^\vee$-opers.

\emph{Specialisation to $\mathfrak{sl}_2$.}
$\dim\mathfrak{sl}_2=3$, $h^\vee=2$.
$\kappa = 3(k+2)/4$.
$r(z) = (e\otimes f + f\otimes e + \tfrac{1}{2}h\otimes h)/z$.
$\nabla^{\mathrm{hol}}_{0,n} =
d - \frac{1}{k+2}\sum_{i<j}\Omega_{ij}\,d(z_i{-}z_j)/(z_i{-}z_j)$
(the $\widehat{\mathfrak{sl}}_2$ KZ connection).
Genus-$1$ Hessian: $H^{(1)}=\kappa=3(k+2)/4$ (scalar,
one-dimensional deformation space via Sugawara).
Spectral discriminant:
$\Delta(x) = (1-kx)(1-(k{+}4)x)/(1-2x)$;
zeros at $x=1/k$ (level~$k$) and $x=1/(k{+}4)$ (dual level);
pole at $x=1/2$ (critical level $k=-2$).
\end{computation}

\begin{computation}[Derivation of the KZ connection from the
graph-sum formula; \ClaimStatusProvedHere]
\label{comp:kz-from-graph-sum}
\index{KZ connection!derivation from graph sum}
\index{shadow connection!KZ derivation}
The genus-$0$, $n$-point shadow connection
$\nabla^{\mathrm{hol}}_{0,n}=d-\mathrm{Sh}_{0,n}(\Theta_\cA)$
for $\cA=\widehat{\mathfrak{sl}}_2{}_k$ is derived from the
graph-sum formula~\eqref{eq:preface-theta} as follows.

\emph{Contributing graphs.}
At genus~$0$ with $n$ external legs, the contributing stable
graphs are trees with $n$ leaves. The weight-$1$ component
$\Theta^{[1]}_\cA|_{(0,n)}$ sums over trees with one internal
edge, that is, binary trees with two vertices. For $n$ insertions at
points $z_1,\dots,z_n$, each binary tree $T_{ij}$ has one
internal edge connecting vertex sets $\{i\}$ and $\{j\}$.

\emph{Vertex labels.}
At each vertex, the transferred chiral operation is the OPE
residue. For $\widehat{\mathfrak{sl}}_2$, the binary vertex
label is
\[
\Phi_{ij}
\;=\;
\sum_{a} J^a_i \otimes J^a_j
\;=\;
\Omega_{ij}
\;=\;
e_i f_j + f_i e_j + \tfrac{1}{2}h_i h_j.
\]

\emph{Edge weight.}
The FM propagator on $\overline{C}_n(\mathbb{A}^1)$ for the
edge connecting $z_i$ and $z_j$ is
\[
W_{T_{ij}}
\;=\;
\frac{1}{k+h^\vee}\,d\log(z_i-z_j)
\;=\;
\frac{1}{k+2}\,\eta_{ij}.
\]
The factor $1/(k+h^\vee)$ is the normalisation of the
complementarity propagator $P_\cA=H_\cA^{-1}$: the inner
product on $\widehat{\mathfrak{sl}}_2$ at level~$k$ is
$(k+h^\vee)$ times the Killing form, so the inverse
propagator scales as $1/(k+h^\vee)$.

\emph{Graph sum.}
\begin{align*}
\mathrm{Sh}_{0,n}(\Theta_\cA)
&\;=\;
\sum_{i<j}
W_{T_{ij}}\otimes\Phi_{ij}
\\
&\;=\;
\frac{1}{k+2}\sum_{i<j}
\Omega_{ij}\,d(z_i-z_j)/(z_i-z_j).
\end{align*}
The shadow connection is
\[
\boxed{
\nabla^{\mathrm{hol}}_{0,n}
\;=\;
d
\;-\;
\frac{1}{k+2}\sum_{i<j}
\Omega_{ij}\,d(z_i-z_j)/(z_i-z_j),}
\]
which is the Knizhnik--Zamolodchikov connection on
$n$-point conformal blocks of $\widehat{\mathfrak{sl}}_2$
at level~$k$. Flatness
$(\nabla^{\mathrm{hol}}_{0,n})^2=0$ is the
$(0,n)$-projection of the MC equation
$D\Theta_\cA+\tfrac12[\Theta_\cA,\Theta_\cA]=0$.
The three terms in the flatness expansion are:
$\sum_{i<j}\Omega_{ij}\wedge\Omega_{ij}=0$ (the Casimir is
central), $\sum_{i<j<k}[\Omega_{ij},\Omega_{ik}]\,
\eta_{ij}\wedge\eta_{ik}=0$ (the Jacobi identity on
$\mathfrak{sl}_2$), and the Arnold relation kills the
form-theoretic contribution.

At higher genus, the shadow connection acquires corrections
from the period structure:
$\nabla^{\mathrm{hol}}_{g,n}
=d-\frac{1}{k+2}\sum_{i<j}\Omega_{ij}\,\varpi_{ij}
-\kappa\,\omega_g-\cdots$,
where $\varpi_{ij}$ is the prime-form kernel and $\omega_g$
is the first Chern class of the Hodge bundle.
\end{computation}

\begin{computation}[Genus-$1$ Hessian for
$\widehat{\mathfrak{sl}}_3$; \ClaimStatusProvedHere]
\label{comp:sl3-genus1-hessian}
\index{Kac--Moody!genus-1 Hessian!sl3@$\mathfrak{sl}_3$}
\index{spectral discriminant!sl3@$\mathfrak{sl}_3$}
For $\widehat{\mathfrak{sl}}_3$ at level~$k$, the genus-$1$
deformation space has $\dim H^1_{\mathrm{cyc,prim}}=2$
(the Sugawara direction $\partial_k T_{\mathrm{Sug}}$ and the
$\mathcal W_3$-current direction $\partial_k W$).
The genus-$1$ Hessian in this basis is
\[
H^{(1)}_{\widehat{\mathfrak{sl}}_3{}_k}
\;=\;
\begin{pmatrix}
\dfrac{8k}{k+3} & 0 \\[6pt]
0 & \dfrac{8k}{k+3}+\dfrac{48}{5c(k)+22}
\end{pmatrix},
\]
where $c(k)=2-24(k{+}2)^2/(k{+}3)$.
The off-diagonal vanishing
follows from conformal weight selection: $T$ has weight~$2$
and $W$ has weight~$3$. The spectral discriminant is
\[
\Delta_{\widehat{\mathfrak{sl}}_3}(x)
\;=\;
\Bigl(1-\frac{8k}{k+3}\,x\Bigr)
\Bigl(1-\Bigl(\frac{8k}{k+3}+\frac{48}{5c(k)+22}\Bigr)x\Bigr).
\]
Under DS reduction
$\mathrm{DS}\colon\widehat{\mathfrak{sl}}_3{}_k
\to\mathcal W_3{}_{c(k)}$, the Hessian transforms as
\[
H^{(1)}_{\widehat{\mathfrak{sl}}_3}
\;\longmapsto\;
H^{(1)}_{\mathcal W_3}
\;+\;
H^{(1)}_{\mathrm{ghost}},
\]
where $H^{(1)}_{\mathrm{ghost}}$ is the ghost contribution
from the three $bc$-pairs. The eigenvalues of the
affine Hessian are \emph{not} the eigenvalues of the
$\mathcal W_3$ Hessian: DS reduction mixes the Sugawara and
$W$-current directions through the ghost sector, and the
individual eigenvalues shift. What is preserved is the
determinant:
$\det H^{(1)}_{\widehat{\mathfrak{sl}}_3}
=\det H^{(1)}_{\mathcal W_3}\cdot\det H^{(1)}_{\mathrm{ghost}}$.

At the critical level $k=-3$: both eigenvalues diverge;
the Hessian degenerates and the spectral discriminant has
a double pole. The formal moduli problem
$\mathcal M^{\mathrm{mod}}_{\widehat{\mathfrak{sl}}_3{}_{-3}}$
acquires a non-reduced structure (the tangent space has a
Jordan block), reflecting the appearance of the
Feigin--Frenkel center.
\end{computation}

\begin{computation}[Holographic datum for Virasoro;
\ClaimStatusProvedHere]
\label{comp:virasoro-holographic-datum}
\index{Virasoro algebra!holographic datum}
For $\mathrm{Vir}_c$:
\begin{enumerate}[label=\textup{(\roman*)}]
\item $\cA = \mathrm{Vir}_c$.
\item $\cA^! = \mathrm{Vir}_{26-c}$. Self-duality at $c=13$.
\item $\cC = \mathrm{Vir}_c\text{-mod}^{\mathrm{dg}}$.
\item $r_c(z) = \frac{c/2}{z^3} + \frac{2T}{z}$.
\item $\Theta_c$ has infinite shadow depth (mixed archetype);
 $\kappa(\mathrm{Vir}_c) = c/2$ (the dual curvature $(26-c)/2$ vanishes at $c = 26$). The quartic contact
 invariant
 $\mathfrak Q^{\mathrm{contact}} = 10/[c(5c+22)]$.
 The genus-$1$ Hessian correction
 $\delta_H^{(1)} = 120/[c^2(5c+22)]\,x^2$.
\item $\nabla^{\mathrm{hol}}_c =
 d - \frac{c}{2}\,\omega_g
 - \frac{10}{c(5c+22)}\,\delta_4
 - \frac{120}{c^2(5c+22)}\,\delta_H - \cdots$.
\end{enumerate}
At $c=0$: $\kappa=0$ and the bar complex is uncurved.
At $c=26$: the \emph{dual} curvature $(26-c)/2$ vanishes.
At $c=0$, $c=-22/5$: the quartic diverges (resonance points).

\emph{Derivation of the shadow connection from the graph sum.}
The weight-$1$ contribution at genus~$0$ with $n$ insertions is
\[
\mathrm{Sh}^{[1]}_{0,n}(\Theta_c)
\;=\;
\frac{c/2}{}\sum_{i<j}
\frac{2}{(z_i-z_j)^2}\,dz_j
\;+\;
\sum_{i<j}\frac{\partial T(z_j)}{z_i-z_j}\,dz_j,
\]
the sum over binary trees with vertex label
$\Phi_{ij}=\frac{c/2}{(z_i-z_j)^2}+\frac{2T(z_j)}{z_i-z_j}$
(the Virasoro OPE). These are the BPZ differential equations
for $n$-point functions on~$\mathbb{P}^1$.

The weight-$2$ contribution involves the quartic contact stratum
of $\overline{C}_4(X)$:
\[
\mathrm{Sh}^{[2]}_{0,4}(\Theta_c)
\;\supset\;
\frac{10}{c(5c+22)}\;\Lambda_{ijkl}\cdot\delta_4,
\]
where $\Lambda=\mathopen{:}TT\mathclose{:}-\tfrac{3}{10}\partial^2 T$
and $\delta_4$ is the codimension-$2$ contact class on
$\overline{C}_4$. At genus~$g\ge 1$, the curvature
$\kappa=c/2$ enters through the Hodge bundle class
$\omega_g=c_1(\mathbb E)$:
\[
\nabla^{\mathrm{hol}}_c
\;=\;
d
\;-\;\frac{c}{2}\,\omega_g
\;-\;\frac{10}{c(5c+22)}\,\delta_4
\;-\;\frac{120}{c^2(5c+22)}\,\delta_H
\;-\;\cdots\,.
\]
The leading term is the genus expansion; the quartic term is
the first planted-forest correction; the Hessian term is the
second-order genus-$1$ variation. At $c=0$: $\kappa=0$ and
the leading term vanishes (the Witt algebra is uncurved).
At $c=26$: the \emph{dual} curvature $\kappa(\mathrm{Vir}_0)=0$
vanishes, so the dual scalar term disappears; this does not by itself
collapse the dual higher-degree shadow obstruction tower.

\emph{Genus-$2$ shell decomposition.}
\begin{align*}
\Theta^{(2)}_c
&= \Theta^{(2)}_{\mathrm{loop}^2}
+ \Theta^{(2)}_{\mathrm{sep}\circ\mathrm{loop}}
+ \Theta^{(2)}_{\mathrm{pf}},\\[4pt]
\Theta^{(2)}_{\mathrm{loop}^2}
&= \kappa^2\cdot(\text{iterated BV on }
H^0(\overline{\cM}_{2,0})),\\
\Theta^{(2)}_{\mathrm{sep}\circ\mathrm{loop}}
&= \kappa\cdot\ell_3^{(0)}\cdot(\text{separating
boundary of }\overline{\cM}_{2,0}),\\
\Theta^{(2)}_{\mathrm{pf}}
&= \frac{10}{c(5c{+}22)}\cdot(\text{planted-forest
correction from log-FM strata}).
\end{align*}
Heisenberg sees only the first; affine algebras the first two;
Virasoro all three. The quintic obstruction
$o^{(5)}_{\mathrm{Vir}}\neq 0$ forces infinite shadow depth.

\emph{Spectral discriminant.}
The Hessian eigenvalue is the curvature $\kappa = c/2$, giving
$\Delta_{\mathrm{Vir}}(x)=1-\frac{c}{2}\,x$;
single branch point at $x=2/c$, diverging at $c=0$ (critical level).
\end{computation}

\begin{computation}[Holographic datum for $\beta\gamma$;
\ClaimStatusProvedHere]

\label{comp:betagamma-holographic-datum}
\index{beta-gamma@$\beta\gamma$!holographic datum}
The $\beta\gamma$ system (conformal weights $(\lambda,1{-}\lambda)$)
is the contact/quartic archetype: the cubic bracket
$\ell_3^{(0)}=0$ because the simple pole is scalar, while the
quartic contact invariant $\mu_{\beta\gamma}=0$ by rank-one rigidity
(Corollary~\ref{cor:nms-betagamma-mu-vanishing}).
\begin{enumerate}[label=\textup{(\roman*)}]
\item $\cA = \beta\gamma_\lambda$.
\item $\cA^! = \beta\gamma_{1-\lambda}$.
 Self-duality at $\lambda=1/2$.
\item The closed genus-zero collision residue and the closed
 genus-zero curvature vanish:
 \[
 r_{\beta\gamma}^{\mathrm{cl}}(z)
 :=
 \operatorname{Res}^{\mathrm{coll}}_{0,2}
 (\Theta_{\beta\gamma_\lambda})_{\mathrm{cl}}
 =0,
 \qquad
 \Theta^{\mathrm{cl}}_{0,\mathrm{curv}}(\beta\gamma_\lambda)=0.
 \]
 The simple contact/OPE residue is the separate free-field kernel
 \[
 r_{\beta\gamma}^{\mathrm{OPE}}(z)
 =
 \frac{\beta\otimes\gamma+\gamma\otimes\beta}{z}.
 \]
\item The scalar characteristic is
 \[
 \kappa(\beta\gamma_\lambda)=6\lambda^2-6\lambda+1.
 \]
 Thus $\kappa(\beta\gamma_1)=1$, while the symplectic-boson
 specialization at $\lambda=1/2$ has
 $\kappa(\mathrm{Sb}_{1/2})=-1/2$.
 Shadow depth~$4$ (contact archetype).
\item The simple contact/OPE connection is
 $\nabla^{\mathrm{hol}}_{0,n} =
 d + \sum_{i<j}(\beta_i\gamma_j+\gamma_i\beta_j)\,
 d(z_i{-}z_j)/(z_i{-}z_j)$.
\end{enumerate}
Genus-$2$ shell profile: only
$\Theta^{(2)}_{\mathrm{loop}^2}$ survives; the separating
and planted-forest shells vanish because $\ell_3^{(0)}=0$
and $\mu_{\beta\gamma}=0$.
Spectral discriminant on the symplectic-boson branch-operator
normalization: $\Delta_{\mathrm{Sb}}(x) = 1+2x$;
single branch point at $x=-1/2$. This is not the scalar
Theorem~D polynomial $1-\kappa x$.

The symplectic-boson specialization is the standard free-field lane
where the one-pair scalar is negative: the genus tower carries
$F_g(\mathrm{Sb}_{1/2}) = -\tfrac12\lambda_g^{\mathrm{FP}}$, and the
$\hat A$-series is evaluated at imaginary argument. The Burns-space
lane $\lambda=1$ has $\kappa=1$ per pair and does not have this sign.
\end{computation}

\begin{computation}[Holographic datum for $\mathcal W_3$;
\ClaimStatusProvedHere]
\label{comp:w3-holographic-datum}
\index{W3@$\mathcal W_3$!holographic datum}
For $\mathcal W_3$ at central charge~$c$, generated by $T$
(weight~$2$) and $W$ (weight~$3$):
\begin{enumerate}[label=\textup{(\roman*)}]
\item $\cA = \mathcal W_3{}_c$.
\item the companion algebra has central charge $100-c$ and satisfies
 $\kappa((\mathcal W_3{}_c)^!)=\kappa(\mathcal W_3{}_{100-c})$;
 the same-family identification
 $(\mathcal W_3{}_c)^!\simeq\mathcal W_3{}_{100-c}$ is conditional
 on DS/bar transport.
\item $\cC = \cC_c^{W_3}$, a module category for the dg-shifted
 $\mathcal W_3$-Yangian.
\item $r_c^{W_3}(z)$: the $\mathcal W_3$ $r$-matrix.
\item $\kappa(\mathcal W_3{}_c) = 5c/6$.
 The Hessian eigenvalue $(c-50)/2$ governs the spectral discriminant.
 The quartic contact shadow is controlled by
 $\Lambda = {:}TT{:} - \tfrac{3}{10}\partial^2 T$
 with $[\ell_4^{(0)}] = \frac{16}{22+5c}\,\Lambda$.
 The genus-$1$ Hessian is
 $H^{(1)} = \operatorname{diag}\bigl(
 \kappa_c,\;\kappa_c + 48/(5c+22)\bigr)$
 in the basis $(\partial_c T,\,\partial_c W)$.
 Spectral discriminant
 $\Delta(x) = (1-\kappa_c x)(1-(\kappa_c+48/(5c+22))x)$.
\item $\nabla^{\mathrm{hol}}_c =
 d - \frac{5c}{6}\,\omega_g
 - \frac{16}{22+5c}\,\Lambda\cdot\delta_4 - \cdots$.
\end{enumerate}
At $c=0$: $\kappa=0$ and the bar complex is uncurved.
At $c=100$: the companion curvature $\kappa(\mathcal W_3{}_{100-c})=0$.
All three genus-$2$ shells present.
\end{computation}

\begin{computation}[Holographic datum for subregular
$W_k(\mathfrak{sl}_3,f_{\mathrm{sub}})$]
\label{comp:subregular-sl3-holographic-datum}
\index{W-algebra@$\mathcal W$-algebra!subregular!holographic datum}
\index{non-principal W-algebra!holographic datum}
The subregular $W$-algebra
$W_k(\mathfrak{sl}_3,f_{\mathrm{sub}})$ is obtained by
quantum DS reduction at the subregular nilpotent
$f_{\mathrm{sub}}\in\mathfrak{sl}_3$. It is generated by
$T$ (weight~$2$), $G^\pm$ (weight~$3/2$), and $J$ (weight~$1$);
it is the $\mathcal N=2$ superconformal algebra at
$c = 3(2k+1)/(k+3)$.
\begin{enumerate}[label=\textup{(\roman*)}]
\item $\cA = W_k(\mathfrak{sl}_3,f_{\mathrm{sub}})$.
\item $\cA^!$: conjectured to be
 $W_{k^\vee}(\mathfrak{sl}_3,f_{\mathrm{sub}})$
 at the dual level $k^\vee=-k-6$
 (the type-A transport-to-transpose conjecture,
 Conjecture~\ref{conj:type-a-transport-to-transpose-frontier}).
\item $\kappa(W_k(\mathfrak{sl}_3,f_{\mathrm{sub}}))
 = \kappa(\widehat{\mathfrak{sl}}_3{}_k)
 - \kappa_{\mathrm{ghost}}
 = \frac{8k}{k+3} - \kappa_{\mathrm{ghost}}$,
 where $\kappa_{\mathrm{ghost}}$ is the modular characteristic
 of the two $bc$-ghost pairs at the Jacobson--Morozov heights of
 the two nilradical roots. From
 $c = 3(2k+1)/(k+3)$ and the Sugawara embedding this gives
 $\kappa = c/2 = 3(2k+1)/[2(k+3)]$ on the Virasoro subalgebra;
 the full $\kappa$ on the $\mathcal{N}{=}2$ SCA requires the
 contribution from $G^\pm$ and $J$, which depends on the
 normalisation of the non-principal invariant pairing
 and is not computed here.
\item The shadow depth is $\infty$ (the $G^\pm$ generators
 create nonvanishing higher brackets at all degrees).
\item DS acts as
 $\mathrm{DS}_{f_{\mathrm{sub}}}\colon
 \mathcal H(\widehat{\mathfrak{sl}}_3{}_k)
 \to \mathcal H(W_k(\mathfrak{sl}_3,f_{\mathrm{sub}}))$;
 the ghost central charge is
 $c_{\mathrm{ghost}} = 2\cdot|\Delta^+_{f_{\mathrm{sub}}}|
 = 2\cdot 2 = 4$
 (two positive roots in the nilradical of the
 Jacobson--Morozov $\mathfrak{sl}_2$-triple).
\end{enumerate}
The non-principal case differs from the principal case in two
respects: the ghost central charge depends on the nilpotent
orbit, and the Koszul dual is conjectured to be the
\emph{transpose} reduction
$W_{k^\vee}(\mathfrak{sl}_3,f_{\mathrm{sub}}^t)$
rather than the same nilpotent at dual level. Proving
bar-cobar/Koszul commutation with non-principal DS reduction
is a frontier problem
(Conjecture~\ref{conj:ds-kd-arbitrary-nilpotent}).
\end{computation}

\begin{construction}[Drinfeld--Sokolov reduction of holographic
data]
\label{constr:ds-reduction-holographic}
\index{Drinfeld--Sokolov reduction!holographic data}
Principal DS reduction
$\mathrm{DS}\colon\widehat{\mathfrak{sl}}_N{}_k
\to\mathcal W_N{}_{c(k)}$
with $c(k) = (N-1) - N(N^2-1)(k+N-1)^2/(k+N)$ extends to a morphism
of holographic data
\[
\mathrm{DS}\colon
\mathcal H(\widehat{\mathfrak{sl}}_N{}_k)
\;\longrightarrow\;
\mathcal H(\mathcal W_N{}_{c(k)}).
\]
On each component:
\begin{enumerate}[label=\textup{(\roman*)}]
\item $\cA$: $\widehat{\mathfrak{sl}}_N{}_k \mapsto
 \mathcal W_N{}_{c(k)}$.
\item $\cA^!$: $\widehat{\mathfrak{sl}}_N{}_{-k-2N} \mapsto
 \mathcal W_N{}_{c(-k-2N)}$.
\item $r(z)$: $k\,\Omega_{\mathfrak{sl}_N}/z \mapsto
 r_c^{W_N}(z)$.
\item $\kappa$:
 $\dfrac{(N^2-1)(k+N)}{2N} \mapsto c(k)\,(H_N - 1)$,
 where $H_N = \sum_{j=1}^N j^{-1}$ is the $N$-th harmonic number.
 For $N=2$: $\kappa(\widehat{\mathfrak{sl}}_2) = 3(k+2)/4$,
 $\kappa(\mathcal W_2) = c/2$; for $N=3$:
 $\kappa(\widehat{\mathfrak{sl}}_3) = 4(k+3)/3$,
 $\kappa(\mathcal W_3) = 5c/6$.
\item Shadow depth: $3 \mapsto \infty$ for $N \geq 3$.
\item $\Theta^{(2)}_{\mathrm{pf}}$: $0 \mapsto \neq 0$
 for $N \geq 3$.
\end{enumerate}
\end{construction}

\begin{theorem}[Central charge additivity under DS;
\ClaimStatusConditional]

\label{thm:ds-central-charge-additivity}
\index{Drinfeld--Sokolov reduction!central charge additivity}
\textup{[}Type signature:
$(\mathrm{Open}\cap\mathrm{Cat},\mathrm{factorisation},
\mathrm{level}=1\leftrightarrow 5,
\mathrm{hypothesis\ package}=\mathrm{principal\ }f\in\fg;
k\ne -h^\vee;
\mathrm{Kac\text{--}Roan\text{--}Wakimoto\ BRST\ construction};
\mathrm{Yamada\ generic\text{-}level\ finite\ presentation\ for\ quantum\ }\Walg_k(\fg);
\mathrm{Prochazka\ }\Walg_N\mathrm{\ closure\ \textup{(}finite\text{-}spin\ endpoint\ for\ shadow\ jump\textup{)}};
\mathrm{KZ\ analytic\ SDR\ for\ part~\textup{(iii)}};
\mathrm{Vol\ III\ }(\Sigma_{\mathrm{rank}\,\fg-1},C)\text{-specialisation\ for\ stage\text{-}2\ shadow})$\textup{]}
Let $\mathfrak g$ be a finite-dimensional simple Lie algebra with
dual Coxeter number~$h^\vee$ and positive roots~$\Delta^+$.
For principal DS reduction
$\widehat{\mathfrak g}_k\to\mathcal W_k(\mathfrak g)$:
\begin{enumerate}[label=\textup{(\roman*)}]
\item \emph{Central charge decomposition.}
The DS BRST complex gives a $k$-dependent central charge shift:
\begin{equation}\label{eq:ds-central-charge-additivity}
c(\widehat{\mathfrak g}_k) - c(\mathcal W_k(\mathfrak g))
\;=\;
c_{\mathrm{ghost}}^{bc}
\;+\;
c_{\mathrm{improvement}}(k),
\end{equation}
where $c_{\mathrm{ghost}}^{bc}$ is the $k$-independent ghost
$bc$ contribution and $c_{\mathrm{improvement}}(k)$ is the
stress-tensor improvement shift. For $\mathfrak{sl}_N$:
$c(\widehat{\mathfrak{sl}}_N{}_k) - c(\mathcal{W}_N, k)
= N(N{-}1)\bigl((N{+}1)k + N^2{-}N{-}1\bigr)$.
\item \emph{Modular characteristic additivity.}
\begin{equation}\label{eq:ds-kappa-decomposition}
\kappa(\widehat{\mathfrak g}_k)
\;=\;
\kappa(\mathcal W_k(\mathfrak g))
\;+\;
\kappa(\mathrm{ghost}).
\end{equation}
\item \emph{Genus-$g$ decomposition.}
\begin{equation}\label{eq:ds-genus-g-decomposition}
F_g(\widehat{\mathfrak g}_k)
\;=\;
F_g(\mathcal W_k(\mathfrak g))
\;+\;
F_g(\mathrm{ghost})
\qquad\text{for all }g\ge 1.
\end{equation}
\item \emph{Shadow depth jump.}
For $\operatorname{rank}\mathfrak g\ge 2$, the shadow depth
jumps from~$3$ (affine = Lie/tree archetype) to~$\infty$
($\mathcal W$-algebra = mixed archetype). The planted-forest
shell $\Theta^{(2)}_{\mathrm{pf}}$ is zero on the affine side
and nonzero on the $\mathcal W$-algebra side: DS reduction
creates the planted-forest correction.
\end{enumerate}
\end{theorem}

\begin{proof}
(i)~The DS BRST complex consists of the affine vertex algebra, a
ghost system $\{b_\alpha,c_\alpha\}_{\alpha\in\Delta^+}$,
and an improvement to the stress tensor involving the affine
Cartan current. The ghost conformal weights are
$(1{-}j_\alpha,\,j_\alpha)$ where $j_\alpha$ is the height
of~$\alpha$ in the principal grading.
The ghost $bc$-pairs contribute a $k$-independent central
charge, but the stress-tensor improvement contributes
a $k$-dependent shift, so the total difference
$c(\hat{\mathfrak{g}}_k) - c(\mathcal{W}_k)$ depends on~$k$.
For $\mathfrak{sl}_N$, the Fateev--Lukyanov formula gives
$c(\mathcal{W}_N, k) = (N{-}1) - N(N^2{-}1)(k{+}N{-}1)^2/(k{+}N)$
and direct computation yields
\[
c(\hat{\mathfrak{sl}}_N{}_k) - c(\mathcal{W}_N, k)
\;=\;
N(N{-}1)\bigl((N{+}1)k + N^2{-}N{-}1\bigr),
\]
which is linear in~$k$ and reduces to $N(N{-}1)(N^2{-}N{-}1)$
at~$k = 0$.

Explicit verification for $\mathfrak{sl}_N$:
\begin{center}
\renewcommand{\arraystretch}{1.15}
\begin{tabular}{lccc}
$\mathfrak g$ & $\dim$ & $|\Delta^+|$
 & $c_{\mathrm{aff}}-c_{W}$ \\
\hline
$\mathfrak{sl}_2$ & $3$ & $1$ & $2(3k{+}1)$ \\
$\mathfrak{sl}_3$ & $8$ & $3$ & $6(4k{+}5)$ \\
$\mathfrak{sl}_4$ & $15$ & $6$ & $12(5k{+}11)$ \\
$\mathfrak{sl}_5$ & $24$ & $10$ & $20(6k{+}19)$
\end{tabular}
\end{center}

\noindent
(ii)~From~(i) by the genus universality theorem:
$\kappa$ is determined by the genus-$1$ curvature
(Theorem~\ref{thm:genus-universality}) applied to each factor.

\noindent
(iii)~The genus-$g$ free energy is
$F_g=\kappa\cdot\lambda_g^{\mathrm{FP}}$ at the scalar level
(Theorem~\ref{thm:genus-universality}); additivity of~$\kappa$
gives additivity of~$F_g$. At the non-scalar level, the
graph-sum formula~\eqref{eq:preface-theta} is additive under
tensor products by
Proposition~\ref{prop:independent-sum-factorization}.

\noindent
(iv)~$\ell_4^{(0)}=0$ for $\widehat{\mathfrak g}_k$ (strict
formality, Theorem~\ref{thm:km-strictification}); $\ell_4^{(0)}\neq 0$
for $\mathcal W_k(\mathfrak g)$ when $\operatorname{rank}\ge 2$
(the quartic contact invariant is nonzero).
\end{proof}

\begin{corollary}[Critical dimensions;
\ClaimStatusProvedHere]
\label{cor:critical-dimensions}
\index{critical dimension|textbf}
The Koszul dual bar complex
$\barB(\cA^!)$ is uncurved if and only if
$\kappa(\cA^!)=0$. For the standard families:
\begin{center}
\renewcommand{\arraystretch}{1.2}
\begin{tabular}{llll}
\textbf{Algebra} & $\kappa(\cA)$ & $\kappa(\cA^!)$
 & \textbf{$\kappa^!=0$ at} \\
\hline
$\mathrm{Vir}_c$ & $c/2$ & $(26-c)/2$ & $c=26$ \\
$\mathcal W_N{}_c$ & $c\cdot\varrho_N$
 & $\varrho_N K_N - c\cdot\varrho_N$
 & $c=K_N$ \\
$\widehat{\mathfrak g}_k$
 & $(k{+}h^\vee)\dim\mathfrak g/(2h^\vee)$
 & $-(k{+}h^\vee)\dim\mathfrak g/(2h^\vee)$
 & $k=-h^\vee$
\end{tabular}
\end{center}
\end{corollary}

\begin{proof}
$\nabla^{\mathrm{hol}}=d-\kappa\omega_g-\cdots$; $\kappa=0$
gives flatness at leading order.
\end{proof}


%% ====================================================================
%% FACTORIZATION-ENVELOPE SHADOW FUNCTOR
%% ====================================================================
\subsection{The factorization-envelope shadow functor}
\label{subsec:factorization-envelope-shadow-functor}
\index{factorization envelope!shadow functor|textbf}
\index{envelope-shadow functor|textbf}
\index{Lie conformal algebra!cyclically admissible|textbf}

The preceding subsections organized the holographic programme around
a single datum~$\mathcal{H}(T)$ attached to an HT system. But the
datum is not yet functorial: it is assembled from a specific bulk
algebra~$\cA$ and its Koszul dual. The factorization-envelope
technology (Nishinaka's construction of holomorphic factorization
algebras from Lie conformal algebras~\cite{Nishinaka25}, and Vicedo's
full/non-chiral analogue~\cite{Vicedo25}) supplies the missing
functorial input. Starting from a Lie conformal algebra~$R$, the
factorization envelope produces a vertex algebra isomorphic to the
enveloping vertex algebra~$\Uvert(R)$; the shadow obstruction tower
then becomes a functor on the input data, not a case-by-case invariant
of isolated vertex algebras.

\begin{definition}[Cyclically admissible Lie conformal algebra]
\label{def:v1-cyclically-admissible-lca}
\index{Lie conformal algebra!cyclically admissible!definition}
A \emph{cyclically admissible Lie conformal algebra} on a smooth
projective curve~$X$ is a Lie conformal algebra~$R$ satisfying:
\begin{enumerate}[label=\textup{(\roman*)}]
\item \emph{Grading.}
 $R = \bigoplus_{h \geq 0} R_h$ carries a conformal-weight grading
 with finite-dimensional homogeneous components.
\item \emph{Filtration.}
 The descending filtration
 $F^m R := \bigoplus_{h \geq m} R_h$ is complete:
 $R = \varprojlim_m R / F^m R$.
\item \emph{Bounded pole order.}
 The $\lambda$-bracket
 $[\,\cdot\,{}_\lambda\,\cdot\,]\colon R \otimes R \to R[\lambda]$
 has bounded pole order: $[a_\lambda b] = \sum_{j=0}^{N(a,b)}
 \lambda^j c_j$ for some finite $N(a,b)$.
\item \emph{Invariant residue pairing.}
 There exists a nondegenerate residue pairing
 \[
 \langle -, - \rangle \colon R \otimes R \longrightarrow \omX
 \]
 compatible with translation ($\langle \partial a, b \rangle
 + \langle a, \partial b \rangle = 0$) and satisfying
 $\langle a, b \rangle = -(-1)^{|a||b|} \langle b, a \rangle$.
\end{enumerate}
Write $\LCA_{\mathrm{cyc}}(X)$ for the category of cyclically
admissible Lie conformal algebras on~$X$.
\end{definition}

\begin{definition}[Envelope-shadow functor]
\label{def:envelope-shadow-functor-frontier}
\index{envelope-shadow functor!definition}
For $R \in \LCA_{\mathrm{cyc}}(X)$ and $r \geq 2$, define the
\emph{envelope-shadow functor} by
\begin{equation}\label{eq:envelope-shadow-functor-frontier}
\Thetaenv_{\leq r}(R)
\;:=\;
\Theta^{\leq r}_{\Uvert(R)},
\end{equation}
where $\Uvert(R)$ is the enveloping vertex algebra and
$\Theta^{\leq r}$ denotes the order-$r$ truncation of the
shadow obstruction tower
(Definition~\ref{def:shadow-postnikov-tower}).
Define the \emph{envelope-shadow complexity}
\begin{equation}\label{eq:envelope-shadow-complexity-frontier}
\chienv(R)
\;:=\;
\inf\!\bigl\{\, r \geq 2 \,:\,
 o_{r+1}\bigl(\Uvert(R)\bigr) = 0
 \text{ and all } o_{s+1} = 0 \text{ for } s > r
\,\bigr\}
\;\in\;
\{2, 3, 4, \ldots\} \cup \{\infty\},
\end{equation}
the smallest degree at which the shadow obstruction tower terminates.
\end{definition}

\begin{proposition}[Finite-jet rigidity; \ClaimStatusProvedHere]
\label{prop:finite-jet-rigidity-frontier}
\index{finite-jet rigidity!envelope-shadow}
\textup{[}Type signature:
$(\mathrm{Open}\cap\mathrm{Chain},\mathrm{chain\text{-}level},
\mathrm{level}=1\to 4,
\mathrm{hypothesis\ package}=R\in\LCA_{\mathrm{cyc}}(X);
\mathrm{descending\ PBW\ filtration}\ F^\bullet\Uvert(R))$\textup{]}
Let $R \in \LCA_{\mathrm{cyc}}(X)$ and let $r \geq 2$.
The order-$r$ envelope-shadow
$\Thetaenv_{\leq r}(R)$
depends only on the $r$-jet of the enveloping vertex algebra:
it is determined by the quotient
$\Uvert(R) / F^{r+1}\Uvert(R)$, where $F^\bullet$ is the
descending filtration by PBW degree.
\end{proposition}

\begin{proof}
The shadow obstruction tower at degree~$r$ involves the transferred
cyclic minimal model operations $m_2, \ldots, m_r$ and graph sums
over connected stable graphs of total weight at most~$r$
(Definition~\ref{def:modular-bar-hamiltonian}).
Each such graph involves at most~$r$ vertex labels from the
transferred operations, and each transferred operation~$m_k$ is
determined by the first~$k$ layers of the PBW filtration.
Hence the weight-$r$ piece of the bar complex, and therefore the
degree-$r$ shadow, depends only on $\Uvert(R) / F^{r+1}\Uvert(R)$.
\end{proof}

\begin{theorem}[Level-polynomial theorem; \ClaimStatusProvedHere]
\label{thm:level-polynomial}
\index{level-polynomial theorem|textbf}
\textup{[}Type signature:
$(\mathrm{Open}\cap\mathrm{Chain},\mathrm{chain\text{-}level},
\mathrm{level}=1\to 4,
\mathrm{hypothesis\ package}=\mathrm{cyclically\ admissible\ LCA\ family\ \textup{(}Definition~\ref{def:v1-cyclically-admissible-lca}\textup{)}};
\mathrm{single\ scalar\ central\ extension\ direction})$\textup{]}
Let $\{R_t\}_{t \in \C}$ be a one-parameter family of cyclically
admissible Lie conformal algebras differing only by a scalar
central extension: $R_t = R_0 \oplus \C \cdot K$ with
$[a_\lambda K] = t \cdot \langle a, - \rangle$, where
$\langle -, - \rangle$ is the invariant pairing.
Then for each $r \geq 2$, the envelope-shadow
$\Thetaenv_{\leq r}(R_t)$ is polynomial in~$t$ of degree at
most~$r - 1$. In particular:
\begin{enumerate}[label=\textup{(\alph*)}]
\item $\kappa\bigl(\Uvert(R_t)\bigr)$ is linear in~$t$;
\item the cubic shadow $\mathfrak{C}$ is at most quadratic in~$t$;
\item the quartic resonance class
 $R^{\mathrm{mod}}_{4}$ is at most cubic in~$t$.
\end{enumerate}
\end{theorem}

\begin{proof}
The transferred cyclic operations $m_k$ on $\Uvert(R_t)$ arise
from bar-cobar transfer applied to the relations of the
enveloping vertex algebra. Each operation~$m_k$ involves at
most~$k - 1$ OPE contractions, each of which contributes at
most one power of~$t$ from the central extension. Therefore the
graph-sum formula at degree~$r$ is a polynomial of degree at
most~$r - 1$ in~$t$: the degree-$r$ graph sum involves edges
carrying the complementarity propagator $P = H^{-1}$, which is
independent of~$t$, and vertex labels~$m_k$ each contributing
at most $k - 1$ powers of~$t$. Since all graphs in the sum have
total vertex-weight at most~$r$ and each vertex contributes at
most $\operatorname{val}(v) - 1$ powers, the total $t$-degree is
bounded by $\sum_v (\operatorname{val}(v) - 1) \leq r - 1$ for
connected graphs of weight~$r$.
\end{proof}

\begin{theorem}[Gaussian collapse; \ClaimStatusProvedHere]
\label{thm:gaussian-collapse}
\index{Gaussian collapse!abelian envelope|textbf}
\textup{[}Type signature:
$(\mathrm{Open}\cap\mathrm{Chain},\mathrm{chain\text{-}level},
\mathrm{level}=1\to 5,
\mathrm{hypothesis\ package}=\mathrm{abelian\ cyclically\ admissible\ }R;
\mathrm{class}\ \mathbf{G}\ \mathrm{archetype\ scope;\ does\ not\ propagate\ to\ Lie/Heisenberg/free\text{-}field\ subfamilies\ with\ nontrivial\ brackets})$\textup{]}
If $R$ is abelian (that is, $[a_\lambda b] = 0$ for all
$a, b \in R$), then $\chienv(R) = 2$. Equivalently, the
shadow obstruction tower of an abelian envelope terminates at degree~$2$:
\[
\Thetaenv_{\leq r}(R) = \Thetaenv_{\leq 2}(R)
\qquad \text{for all } r \geq 2.
\]
The only shadow is the scalar modular characteristic
$\kappa\bigl(\Uvert(R)\bigr)$.
\end{theorem}

\begin{proof}
If $R$ is abelian, the enveloping vertex algebra $\Uvert(R)$ is the
symmetric algebra on~$R$, which is the Heisenberg-type factorization
algebra. Its bar complex is quasi-free, and the transferred minimal
model has $m_k = 0$ for all $k \geq 3$. Since the graph-sum
formula at degree~$r \geq 3$ requires at least one vertex with
$m_k$, $k \geq 3$, all such contributions vanish. Therefore
$o_{r+1}(\Uvert(R)) = 0$ for all $r \geq 2$, and
$\chienv(R) = 2$.
\end{proof}

\begin{proposition}[Independent sums factor; \ClaimStatusProvedHere]
\label{prop:independent-sums-factor}
\index{independent sums!shadow factorization}
Let $R = R_1 \oplus R_2$ with all mixed $\lambda$-brackets
vanishing. Then the enveloping vertex algebra factors as
$\Uvert(R) \cong \Uvert(R_1) \,\widehat{\otimes}\, \Uvert(R_2)$
and all shadow data separate:
\begin{align}
\kappa\bigl(\Uvert(R)\bigr)
 &= \kappa\bigl(\Uvert(R_1)\bigr)
 + \kappa\bigl(\Uvert(R_2)\bigr),
\label{eq:kappa-additivity-frontier}\\
T^{\mathrm{br}}_R
 &= T^{\mathrm{br}}_{R_1} \oplus T^{\mathrm{br}}_{R_2},
\label{eq:branch-additivity}\\
\Delta_R(x)
 &= \Delta_{R_1}(x) \cdot \Delta_{R_2}(x),
\label{eq:discriminant-multiplicativity}\\
R^{\mathrm{mod}}_{4}(R)
 &= R^{\mathrm{mod}}_{4}(R_1)
 + R^{\mathrm{mod}}_{4}(R_2).
\label{eq:quartic-additivity}
\end{align}
\end{proposition}

\begin{proof}
The tensor-product decomposition of enveloping vertex algebras for
orthogonal direct sums is standard
(PBW for vertex algebras).
Under this decomposition the bar complex splits as a tensor product
of bar complexes, the transferred cyclic operations have no mixed
contributions, and the graph-sum formula at each degree decomposes
into independent contributions from each summand. The additivity
and multiplicativity of the shadows follow from the corresponding
properties of the Maurer--Cartan equation applied to each
independent factor.
\end{proof}

\begin{conjecture}[Universal modular factorization envelope;
\ClaimStatusConjectured]
\label{conj:universal-modular-factorization-envelope}
\index{modular factorization envelope!universal}
\index{primitive-current functor!modular}
\textup{[}Type signature:
$(\mathrm{Open}\cap\mathrm{Cat},\mathrm{factorisation},
\mathrm{level}=1\leftrightarrow 5,
\mathrm{hypothesis\ package}=\mathrm{Nishinaka\ envelope\ at\ genus\ }0;
\mathrm{modular\ category\ extension\ of\ adjunction};
\mathrm{KSDual\ sublocus\ for\ Verdier\text{-}symmetric\ adjunction};
\mathrm{modular\ data\ realised\ via\ trace\ +\ clutching\ on\ }\overline{\cM}_{g,n},\ \mathrm{not\ as\ closed\text{-}algebra\ property\ of\ }\Uvert(R))$\textup{]}
There exists a functor
\begin{equation}\label{eq:universal-modular-envelope}
U_X^{\mathrm{mod}}
\;\colon\;
\LCA_{\mathrm{cyc}}(X)
\;\longrightarrow\;
\mathrm{ModKoszul}(X)
\end{equation}
from cyclically admissible Lie conformal algebras to modular Koszul
factorization objects, left adjoint to the modular
primitive-current functor
$\operatorname{Prim}_{\mathrm{mod}} \colon
\mathrm{ModKoszul}(X) \to \LCA_{\mathrm{cyc}}(X)$:
\begin{equation}\label{eq:envelope-adjunction}
\Hom_{\mathrm{ModKoszul}(X)}\!\bigl(
 U_X^{\mathrm{mod}}(R),\, \cA
\bigr)
\;\cong\;
\Hom_{\LCA_{\mathrm{cyc}}(X)}\!\bigl(
 R,\, \operatorname{Prim}_{\mathrm{mod}}(\cA)
\bigr).
\end{equation}
At genus~$0$, this recovers the Nishinaka envelope. The output
carries, functorially and simultaneously, the full modular
characteristic class~$\Theta$, the scalar shadow~$\kappa$, the
spectral discriminant~$\Delta$, the branch operator~$T^{\mathrm{br}}$,
and the quartic resonance class~$R^{\mathrm{mod}}_4$.
\end{conjecture}

\begin{evidence}
At genus~$0$, the Nishinaka envelope is already left adjoint to the
forgetful/primitive functor from vertex algebras to Lie conformal
algebras; the conjecture extends this adjunction to the modular
category carrying the clutching and Verdier data. The PBW theorem
for universal enveloping vertex algebras
(Proposition~\ref{prop:pbw-universality}) ensures that the
genus-$0$ input is always chirally Koszul, so the modular Koszul
target category is nonempty. The remaining content is that the
modular completion of the envelope inherits the adjunction.
The envelope-shadow functor
$\Thetaenv_{\leq r}(R) := \Theta^{\leq r}_{\Uvert(R)}$
(Definition~\ref{def:envelope-shadow-functor-frontier}) already
produces the shadow obstruction tower \emph{functorially} from $R$, including
level-polynomial dependence
(Theorem~\ref{thm:level-polynomial}), Gaussian collapse
(Theorem~\ref{thm:gaussian-collapse}), and independent-sum
factorization (Proposition~\ref{prop:independent-sums-factor}).
The conjecture is that this functorial shadow structure
lifts to a genuine adjunction at the modular category level.
\end{evidence}


%% ====================================================================
%% THE QUARTIC SHADOW FROM FILTERED MC
%% ====================================================================
\subsection{The quartic shadow from filtered Maurer--Cartan theory}
\label{subsec:quartic-shadow-filtered-mc}
\index{quartic resonance class!from filtered MC|textbf}
\index{Maurer--Cartan element!quartic class from filtered|textbf}

The shadow obstruction tower constructs nonlinear invariants
$\kappa \to \mathfrak{C} \to \mathfrak{Q} \to \cdots$ of increasing
degree. A natural question is: why should the first \emph{stable}
nonlinear invariant appear at degree~$4$, rather than~$3$? The answer
is a general fact about filtered dg~Lie algebras, independent of the
chiral context, which explains why cubic ambiguities are removable by
gauge while the quartic class is canonical.

\begin{theorem}[Quartic stability from filtered MC;
\ClaimStatusProvedHere]
\label{thm:quartic-stability-filtered-mc}
\index{quartic stability!filtered MC theorem|textbf}
\textup{[}Type signature:
$(\mathrm{Open}\cap\mathrm{Chain},\mathrm{chain\text{-}level},
\mathrm{level}=3\to 4,
\mathrm{hypothesis\ package}=\mathrm{complete\ descending\text{-}filtered\ dg\ Lie\ algebra};
\mathrm{vanishing}\ H^1(F^3/F^4,d_2)=0;
\mathrm{H^0\text{-}action\ assumption\ for\ part~(iii)})$\textup{]}
Let $(\mathfrak{g}, F^\bullet)$ be a complete descending-filtered
dg~Lie algebra with differential~$d$ compatible with the filtration,
and let
\[
\Theta = \Theta_2 + \Theta_3 + \Theta_4 + \cdots
\in F^2 \mathfrak{g}^1
\]
be a Maurer--Cartan element, with $\Theta_k \in
F^k \mathfrak{g} / F^{k+1} \mathfrak{g}$ the weight-$k$ component.
Write $d_2 := d + [\Theta_2, -]$ for the twisted differential.
Assume:
\begin{equation}\label{eq:h1-vanishing-quartic}
H^1\bigl(F^3 \mathfrak{g} / F^4 \mathfrak{g},\, d_2\bigr) = 0.
\end{equation}
Then:
\begin{enumerate}[label=\textup{(\roman*)}]
\item The cubic component $\Theta_3$ is $d_2$-exact:
 there exists $\xi \in F^3 \mathfrak{g}^0$ such that
 $\Theta_3 = d_2 \xi$.
\item After the gauge transformation
 $\Theta \mapsto \Theta' = e^{-\xi} \cdot \Theta$,
 one may arrange $\Theta' = \Theta_2 + \Theta'_4 + O(F^5)$.
\item The class
 \begin{equation}\label{eq:quartic-class-definition}
 [\Theta'_4]
 \;\in\;
 H^2\bigl(F^4 \mathfrak{g} / F^5 \mathfrak{g},\, d_2\bigr)
 \end{equation}
 is well-defined, independent of the choice of~$\xi$ up to
 the natural $H^0(F^3 \mathfrak{g} / F^4 \mathfrak{g}, d_2)$-action.
 When $H^0(F^3 \mathfrak{g} / F^4 \mathfrak{g}, d_2) = 0$,
 the class~$[\Theta'_4]$ is canonical.
\end{enumerate}
\end{theorem}

\begin{proof}
Projecting the MC equation
$d\Theta + \tfrac{1}{2}[\Theta, \Theta] = 0$
to $F^3 / F^4$ gives
\[
d_2 \Theta_3 = 0
\qquad\text{in } F^3 \mathfrak{g}^2 / F^4 \mathfrak{g}^2.
\]
The vanishing hypothesis~\eqref{eq:h1-vanishing-quartic} then gives
$\Theta_3 = d_2 \xi$ for some
$\xi \in F^3 \mathfrak{g}^0$. The infinitesimal gauge
transformation $\Theta \mapsto \Theta - d_2 \xi
- [\xi, \Theta_3] + O(F^5) = \Theta_2 + \Theta'_4 + O(F^5)$
kills the cubic term.

Projecting the MC equation to $F^4 / F^5$ after the gauge gives
\[
d_2 \Theta'_4 + \tfrac{1}{2}[\Theta_2, \Theta_2]_4 = 0,
\]
where $[\Theta_2, \Theta_2]_4$ denotes the $F^4/F^5$-component of
the bracket. Since $[\Theta_2, \Theta_2]$ is $d$-closed and
$d_2^2 = 0$ on $F^4/F^5$ (the MC equation at weight~$2$ gives
$d\Theta_2 + \tfrac{1}{2}[\Theta_2, \Theta_2] = 0$ in $F^2/F^3$,
hence $d_2^2 = 0$), the class $[\Theta'_4]$ is a well-defined
element of~$H^2(F^4/F^5, d_2)$. A different choice
$\xi' = \xi + \zeta$ with $d_2 \zeta = 0$ changes $\Theta'_4$ by
$[\zeta, \Theta_2]$, which is in the image of the $H^0$-action.
\end{proof}

\begin{remark}[Application to the shadow obstruction tower]
\label{rem:quartic-stability-shadow-application}
In the modular convolution dg~Lie algebra
$\gAmod$
(Definition~\ref{def:nms-modular-convolution-lie}), with the degree
filtration $F^r \gAmod := \bigoplus_{n \geq r}
(\gAmod)_n$, the MC element is $\Theta_\cA$ and $\Theta_2 = \kappa$.
The twisted differential $d_2 = d + [\kappa, -]$ governs the
linearized deformation theory around the scalar shadow. When
$H^1(F^3/F^4, d_2) = 0$ (which holds for all
strongly-generated Koszul algebras at generic parameters), the cubic
shadow $\mathfrak{C}$ is gauge-trivial, and the first stable
nonlinear invariant is the quartic resonance class
$R^{\mathrm{mod}}_4(\cA)$
(Definition~\ref{def:nms-modular-quartic-resonance-class}). This
is the abstract reason that the quartic class, not the cubic one,
is the primary nonlinear shadow.
\end{remark}


%% ====================================================================
%% ANALYTIC COMPLETION AND SEWING ENVELOPE
%% ====================================================================
\subsection{Analytic completion and sewing envelope}
\label{subsec:analytic-completion-sewing-envelope}
\index{sewing envelope|textbf}
\index{analytic Koszul pair|textbf}
\index{HS-sewing condition|textbf}

The algebraic chiral core~$\cA$ determines the combinatorial
bar-cobar package. But for genuine partition functions, trace-class
amplitudes, and analytic modular invariants, the bare algebra is
insufficient: one needs a completion that remembers every local-to-global
sewing operation at once. The construction by
Moriwaki~\cite{Moriwaki26a,Moriwaki26b} of conformally flat
factorization homology valued in $\mathrm{IndHilb}$ (with sphere
values recovering CFT partition functions under positivity and
continuity assumptions) supplies the prototype. The Heisenberg
case, where the ind-Hilbert completion of the affine Heisenberg VOA
is identified with a conformally flat 2-disk algebra on the symmetric
algebra of the Bergman space, shows that the passage from formal OPE
data to analytic sewing objects is already possible in rank one.

\begin{definition}[Sewing envelope]
\label{def:sewing-envelope-frontier}
\index{sewing envelope!definition}
Let $\cA_{\mathrm{alg}}$ be an algebraic chiral core with
algebraic amplitudes
$A^{\mathrm{alg}}_\Sigma$ for every finite conformally flat
bordism~$\Sigma$ with parametrized in/out collars. For every
matrix coefficient of every such amplitude, define a seminorm on
$\cA_{\mathrm{alg}}$ by
\begin{equation}\label{eq:sewing-seminorm}
p_{\Sigma,\xi,\eta}(a)
\;:=\;
\bigl|
\langle \eta,\,
 A^{\mathrm{alg}}_\Sigma(\xi_1 \otimes \cdots \otimes a
 \otimes \cdots \otimes \xi_m)
\rangle
\bigr|.
\end{equation}
The \emph{sewing envelope} $\widehat{\cA}^{\mathrm{sew}}$ is the
Hausdorff completion of $\cA_{\mathrm{alg}}$ for the locally convex
topology generated by all seminorms~$p_{\Sigma,\xi,\eta}$.
\end{definition}

\begin{proposition}[Universal property of the sewing envelope;
\ClaimStatusProvedHere]
\label{prop:sewing-envelope-universal}
\index{sewing envelope!universal property}
If $E$ is any complete locally convex realization of
$\cA_{\mathrm{alg}}$ for which every algebraic amplitude extends
continuously, then the identity on $\cA_{\mathrm{alg}}$ extends
uniquely to a continuous map
$\widehat{\cA}^{\mathrm{sew}} \to E$.
\end{proposition}

\begin{proof}
Every such realization makes each seminorm $p_{\Sigma,\xi,\eta}$
continuous. Hence the identity
$\cA_{\mathrm{alg}} \to E$ is continuous for the sewing topology.
By completeness of~$E$, it extends uniquely from the Hausdorff
completion.
\end{proof}

\begin{definition}[HS-sewing condition]
\label{def:hs-sewing-condition}
\index{HS-sewing condition!definition}
Let $H = \bigoplus_{n \geq 0} H_n$ be a graded pre-Hilbert core,
and write the pair-of-pants multiplication as sector maps
$m^c_{a,b} \colon H_a \,\widehat{\otimes}\, H_b \to H_c$.
Say~$\cA$ satisfies the \emph{HS-sewing condition} if there exists
$0 < q < 1$ such that
\begin{equation}\label{eq:hs-sewing-condition}
\sum_{a,b,c \geq 0}
q^{a+b+c}\, \|m^c_{a,b}\|_{\mathrm{HS}}^2
\;<\; \infty.
\end{equation}
Under HS-sewing, the weighted pair-of-pants map on the weighted
completion $H_q := \widehat{\bigoplus}_{n \geq 0}\, q^n H_n$ is
Hilbert--Schmidt; consequently, every closed amplitude obtained from
finitely many pairs of pants, cups, and caps is trace class on the
relevant weighted state space.
\end{definition}

\begin{definition}[Analytic Koszul pair]
\label{def:analytic-koszul-pair-frontier}
\index{analytic Koszul pair!definition}
An \emph{analytic Koszul pair} $(\cA, \cA^!)$ consists of:
\begin{enumerate}[label=\textup{(\roman*)}]
\item a sewing envelope $\widehat{\cA}^{\mathrm{sew}}$;
\item an analytic bar coalgebra $\barB^{\mathrm{an}}(\cA)$, defined
 as the closure of the algebraic bar sectors inside the sewing
 envelope in the graph norm of the closable bar differential;
\item a genus-$0$ equivalence
 \[
 D^{\mathrm{an}}_{\mathrm{compl}}(\cA\text{-}\mathrm{mod})
 \;\simeq\;
 D^{\mathrm{an}}_{\mathrm{conil}}
 \bigl(\barB^{\mathrm{an}}(\cA)\text{-}\mathrm{comod}\bigr);
 \]
\item under finite-type dualizability, an identification of
 conilpotent analytic comodules with complete
 $\cA^!$-modules.
\end{enumerate}
\end{definition}

\begin{remark}[Curvature forces coderived passage]
\label{rem:curvature-coderived-passage}
\index{coderived category!curvature forcing}
At genus $g \geq 1$, if the completed bar object carries a curvature
term $m_0^{(g)} \neq 0$, then the fiberwise differential satisfies
$\dfib^2 = \kappa \cdot \omega_g \neq 0$, and ordinary cohomology
ceases to be functorial under sewing. The correct receptacle is the
coderived/contraderived analytic category, not an ordinary derived
category. This is not cosmetic: it is the mechanism by which
higher-genus local-to-global data remain nontrivial after completion.
The passage from the algebraic bar complex to the analytic bar
coalgebra, and from derived to coderived categories, is the exact
point where formal factorization becomes analytic local-to-global
theory.
\end{remark}

\begin{conjecture}[Analytic realization criterion for unitary VOAs;
\ClaimStatusConjectured]
\label{conj:analytic-realization-criterion}
\index{analytic realization!unitary VOA criterion}
\textup{[}Type signature:
$(\mathrm{Open}\cap\mathrm{Cat},\mathrm{factorisation},
\mathrm{level}=1\to 5,
\mathrm{hypothesis\ package}=\mathrm{conformal\ Osterwalder\text{--}Schrader\ axioms};
\mathrm{polynomial\ spectral\ growth};
\mathrm{HS\text{-}sewing\ after\ exponential\ damping};
\mathrm{KSDual\ sublocus\ for\ Verdier\text{-}consistent\ analytic\ Koszul\ pair};
\mathrm{modularity\ via\ trace\ +\ clutching\ on\ open\ category},\ \mathrm{not\ closed\text{-}algebra\ property})$\textup{]}
Let $V$ be a unitary full VOA such that:
\textup{(a)}~its quasi-primary correlators satisfy conformal
Osterwalder--Schrader axioms
\textup{(}\textit{cf.}~Adamo--Moriwaki--Tanimoto~\textup{\cite{AMT24}}\textup{)};
\textup{(b)}~its graded pieces have polynomial spectral growth;
\textup{(c)}~its intertwining/screening OPE maps satisfy HS-sewing
after exponential weight damping.
Then $V$ admits:
\begin{enumerate}[label=\textup{(\roman*)}]
\item a sewing envelope
 $V_{\mathrm{alg}} \to
 \widehat{V}^{\mathrm{sew}} \in \mathrm{IndHilb}$;
\item a conformally flat 2-disk algebra;
\item a genus-$0$ analytic Koszul duality package;
\item a higher-genus coderived shadow $Q_\bullet(V)$ with
 complementarity involution.
\end{enumerate}
\end{conjecture}

\begin{evidence}
Moriwaki~\cite{Moriwaki26b} constructs a conformally flat 2-disk
algebra on $\operatorname{Sym} A^2(D)$ and identifies it with the
ind-Hilbert completion of the affine Heisenberg VOA. The
genus-$0$ analytic Koszul package for Heisenberg is therefore
already implicit in that identification. The Adamo--Moriwaki--Tanimoto
result~\cite{AMT24} extends analytic control to a broader class of
unitary full VOAs; the conjecture asks for the remaining bridge
from conformal OS axioms to the bar-cobar package and coderived
modular shadows.

Hypothesis~\textup{(c)} is the critical convergence condition;
it is established for the entire standard landscape of universal
algebras by the general HS-sewing theorem
(Theorem~\ref{thm:general-hs-sewing}: polynomial OPE growth +
subexponential sector growth $\Longrightarrow$ convergence at
all genera). The Heisenberg sewing theorem
(Theorem~\ref{thm:heisenberg-sewing-vol1}) proves the full
analytic chain for rank-$1$ Heisenberg, including the Fredholm
determinant formula. The conjecture asks for the extension to
unitary VOAs beyond the standard Lie-theoretic landscape.
\end{evidence}


%% ====================================================================
%% NON-PRINCIPAL W-ALGEBRA DUALITY
%% ====================================================================
\subsection{Non-principal \texorpdfstring{$\Walg$}{W}-algebra duality}
\label{subsec:non-principal-w-algebra-duality}
\index{W-algebra!non-principal duality|textbf}
\index{Drinfeld--Sokolov reduction!non-principal transport|textbf}
\index{hook-type!transport corridor|textbf}
\index{Barbasch--Vogan duality!W-algebra transport|textbf}

The principal $\Walg$-algebra duality
$\Walg^k(\mathfrak{g})^! \simeq
\Walg^{-k-2h^\vee}(\mathfrak{g})$
(Feigin--Frenkel) is the proved core. The non-principal extension,
which should relate $\Walg^k(\mathfrak{g}, f)^!$ to
$\Walg^{k^\vee}(\mathfrak{g}, f^{\mathrm{BV}})$ for the
Barbasch--Vogan dual nilpotent~$f^{\mathrm{BV}}$, is the active
frontier. Quantum Drinfeld--Sokolov reduction exists for arbitrary
nilpotent~$f$ by the Kac--Roan--Wakimoto construction. What remains
difficult is proving that bar-cobar/Koszul duality commutes with
arbitrary-nilpotent reduction, identifying the correct dual orbit
under Barbasch--Vogan (or Spaltenstein) duality, and determining the
non-principal level transform.

The architecture is controlled by the following factorization of the
problem.

\begin{remark}[Transport/orbit factorization]
\label{rem:transport-orbit-factorization}
\index{transport/orbit factorization}
The non-principal duality problem splits into two axes:
\[
\text{non-principal duality}
\;=\;
\text{transport along reduction network}
\;+\;
\text{identification of dual orbit/level}.
\]
The first axis is now proved on a substantial type~$A$ corridor; the
second remains the true frontier.
\end{remark}

After the principal Feigin--Frenkel case, the first genuinely
non-principal proved corridor is the hook-type family in type~$A$.
Full hook-type duality is proved; inverse reductions connect all
hook-type nodes; and the 2023--2025 work of Fehily,
Creutzig--Linshaw--Nakatsuka--Sato,
Creutzig--Fasquel--Linshaw--Nakatsuka, Genra--Juillard, and
Butson--Nair places hook-type algebras as building blocks inside
a wider reduction network.

\begin{proposition}[Transport propagation lemma;
\ClaimStatusProvedHere]
\label{prop:transport-propagation-frontier}
\index{transport propagation!hook-type corridor}
\textup{[}Type signature:
$(\mathrm{Open}\cap\mathrm{Cat},\mathrm{factorisation},
\mathrm{level}=1\leftrightarrow 2,
\mathrm{hypothesis\ package}=\mathrm{generic\ level};
\mathrm{Fehily\ +\ CLNS\ hook\ corridor};
\mathrm{Genra\ stages,\ Butson\text{--}Nair\ inverse\ reduction\ \textup{(}finite\text{-}spin\ endpoint\ scopes\textup{)}};
\mathrm{KSDual\ sublocus\ on\ self\text{-}transpose\ vertices})$\textup{]}
Let $\Gamma_N$ be the graph whose vertices are partitions
$\lambda \vdash N$, and whose edges are the proved reduction or
inverse-reduction functors between the corresponding type-$A$
affine $\Walg$-algebras at generic level. Suppose a candidate
duality functor~$D$ satisfies:
\begin{enumerate}[label=\textup{(\alph*)}]
\item \emph{Hook seed.}
 $D\bigl(\Walg^k(\mathfrak{sl}_N, f_\eta)\bigr)
 \simeq \Walg^{k^\vee_\eta}(\mathfrak{sl}_N, f_{\eta^t})$
 for every hook partition~$\eta$;
\item \emph{Edge compatibility.}
 $D$ intertwines every edge of~$\Gamma_N$ with the transposed
 edge.
\end{enumerate}
Then $D$ is forced to satisfy
\begin{equation}\label{eq:transport-propagation}
D\bigl(\Walg^k(\mathfrak{sl}_N, f_\lambda)\bigr)
\;\simeq\;
\Walg^{k^\vee_\lambda}(\mathfrak{sl}_N, f_{\lambda^t})
\end{equation}
for every partition~$\lambda$ in the transport-closure of the hook
vertices inside~$\Gamma_N$.
\end{proposition}

\begin{proof}
Pick a path in $\Gamma_N$ from a hook vertex to~$\lambda$.
Start from the known hook formula and propagate inductively across
the path: at each edge, edge-compatibility forces the duality to
transport the hook formula one step further. The result follows
by induction on path length, using the functoriality of~$D$ on
morphisms of~$\Gamma_N$.
\end{proof}

\begin{conjecture}[Type-$A$ transport-to-transpose;
\ClaimStatusConjectured]
\label{conj:type-a-transport-to-transpose-frontier}
\index{transport-to-transpose!type A conjecture}
\textup{[}Type signature:
$(\mathrm{Open}\cap\mathrm{Cat},\mathrm{factorisation},
\mathrm{level}=1\leftrightarrow 2,
\mathrm{hypothesis\ package}=\mathrm{generic\ level};
\mathrm{full\ }\operatorname{Par}(N)\mathrm{\ coverage\ of\ hook\ corridor};
\mathrm{Fehily/CLNS/CFLN/Genra/Butson\text{--}Nair\ inverse\text{-}reduction\ network\ \textup{(}finite\text{-}spin\ MA\text{-}12\ endpoints\textup{)}};
\mathrm{KSDual\ sublocus\ on\ self\text{-}transpose\ partitions};
\mathrm{Vol\ III\ }(\Sigma_{N-1},C)\text{-specialisation\ at\ stage\text{-}2\ shadow})$\textup{]}
At generic level, the transport-closure of the hook vertices is all
of $\operatorname{Par}(N)$, and Koszul/bar-cobar duality intertwines
every proved reduction or inverse-reduction edge. Consequently,
\begin{equation}\label{eq:type-a-transport-to-transpose}
\bigl(\Walg^k(\mathfrak{sl}_N, f_\lambda)\bigr)^!
\;\simeq\;
\Walg^{k^\vee_\lambda}(\mathfrak{sl}_N, f_{\lambda^t})
\end{equation}
for all $\lambda \vdash N$.
\end{conjecture}

\begin{evidence}
The hook-type corridor supplies the proved seed family. The
2023--2025 inverse Hamiltonian reduction results in type~$A$
make the edge set of~$\Gamma_N$ dense enough that every partition
is reachable from a hook vertex. The remaining content is the
edge-compatibility hypothesis: that Koszul duality commutes with
each proved reduction functor. This is the correct intermediate
conjecture, sharper than the old ``arbitrary Barbasch--Vogan
duality'' slogan but still large enough for the frontier label.
\end{evidence}

\begin{remark}[The key obstruction]
\label{rem:non-principal-key-obstruction}
\index{non-principal duality!key obstruction}
The obstruction is not the existence of quantum Drinfeld--Sokolov
reduction for arbitrary nilpotent~$f$; that construction already
exists. The remaining difficulty is to prove that bar-cobar/Koszul
duality commutes with arbitrary-nilpotent reduction, to identify the
correct dual nilpotent under Barbasch--Vogan or Spaltenstein duality,
and to determine the non-principal level-shift formula. The hook-type
corridor has a proved seed family and a named
obstruction package, but not yet a global theorem.
\end{remark}

\subsection{The relative holographic deformation complex}%
\label{subsec:relative-holographic-deformation}%
\index{relative holographic deformation complex|textbf}%
\index{holographic deformation!relative fiber}

\begin{definition}[Holographic forgetful morphism]%
\label{def:holographic-forgetful}%
\index{holographic forgetful morphism|textbf}%
Let $\mathfrak{g}_T^{\mathrm{SC}}$ be the Swiss-cheese
convolution $L_\infty$-algebra of an HT theory~$T$
with boundary datum~$\cA$. The \emph{holographic
forgetful morphism} is the $L_\infty$-morphism
\begin{equation}\label{eq:holographic-forgetful}
p\colon \mathfrak{g}_T^{\mathrm{SC}}
\;\longrightarrow\;
\gAmod
\end{equation}
that discards all open and mixed vertices and retains
only the closed modular sector.
\end{definition}

\begin{definition}[Relative holographic deformation complex]%
\label{def:relative-holographic-deformation}%
\index{relative holographic deformation complex!definition|textbf}%
Fix the canonical basepoint
$\alpha_T \in \MC(\mathfrak{g}_T^{\mathrm{SC}})$
with $p(\alpha_T) = \Theta_\cA$.
The \emph{relative holographic deformation complex} is
\begin{equation}\label{eq:relative-fiber}
\mathfrak{h}_T
\;:=\;
\operatorname{fib}(p),
\qquad
\mathfrak{h}_T^\Theta
\;:=\;
\bigl(\operatorname{fib}(p)\bigr)_{\alpha_T}\,,
\end{equation}
where $(\cdot)_{\alpha_T}$ denotes the
$\alpha_T$-twisted $L_\infty$-structure with
brackets
\begin{equation}\label{eq:twisted-brackets}
\ell_m^{\alpha_T}(x_1,\ldots,x_m)
\;:=\;
\sum_{k \geq 0}\frac{1}{k!}\,
\ell_{m+k}(\alpha_T^{\otimes k},x_1,\ldots,x_m).
\end{equation}
\end{definition}

\begin{proposition}[Lifts as relative MC elements;
\ClaimStatusProvedHere]%
\label{prop:lifts-as-relative-mc}%
\index{relative holographic deformation complex!MC elements}%
\textup{[}Type signature:
$(\mathrm{Open}\cap\mathrm{Cat},\mathrm{factorisation},
\mathrm{level}=3\to 4,
\mathrm{hypothesis\ package}=\mathrm{Swiss\text{-}cheese\ convolution\ }L_\infty\mathrm{\text{-}algebra}\ \mathfrak{g}_T^{\mathrm{SC}};
\mathrm{canonical\ basepoint}\ \alpha_T\ \mathrm{with}\ p(\alpha_T)=\Theta_\cA;
\mathrm{KSDual\ sublocus\ for\ Verdier\text{-}consistent\ closed\ basepoint})$\textup{]}
Lifts of the fixed boundary modular datum~$\Theta_\cA$
to full $3d$ bulk/boundary/line data are exactly
Maurer--Cartan elements of the twisted fiber:
\begin{equation}\label{eq:lifts-relative-mc}
\bigl\{\alpha \in \MC(\mathfrak{g}_T^{\mathrm{SC}})
\;:\; p(\alpha) = \Theta_\cA\bigr\}
\;\;\xleftrightarrow{\;1{:}1\;}\;\;
\MC(\mathfrak{h}_T^\Theta).
\end{equation}
\end{proposition}

\begin{proof}
A lift with the same closed part has the form
$\alpha_T + \beta$ with $p(\beta) = 0$, hence
$\beta \in \operatorname{fib}(p)$.
The MC equation for $\alpha_T + \beta$ is
$\sum_{n \geq 1}\frac{1}{n!}\,\ell_n(\alpha_T{+}\beta,
\ldots,\alpha_T{+}\beta) = 0$.
Subtracting the MC equation for~$\alpha_T$ and collecting
terms with at least one~$\beta$ gives
$\sum_{m \geq 1}\frac{1}{m!}\,
\ell_m^{\alpha_T}(\beta,\ldots,\beta) = 0$,
which is precisely the MC equation in
$\mathfrak{h}_T^\Theta$.
Conversely, any MC element~$\beta$ in the twisted fiber
produces a lift $\alpha_T + \beta$.
\end{proof}

\begin{corollary}[\ClaimStatusProvedElsewhere]%
\label{cor:holographic-deformation-cohomology}%
\index{relative holographic deformation complex!cohomology}%
The true open content of holography is controlled by
$H^\bullet(\mathfrak{h}_T^\Theta)$\textup:
\begin{enumerate}[label=\textup{(\roman*)}]
\item first-order relative deformations live in
 $H^1(\mathfrak{h}_T^\Theta)$\textup;
\item obstructions to extending deformations live in
 $H^2(\mathfrak{h}_T^\Theta)$.
\end{enumerate}
\end{corollary}

\begin{definition}[Holographic obstruction tower]%
\label{def:holographic-obstruction-tower}%
\index{holographic obstruction tower|textbf}%
Impose the filtration $F^\bullet$ on
$\mathfrak{h}_T^\Theta$ by open/mixed vertex count
\textup{(}color depth in the Swiss-cheese graph
expansion\textup{)}.
Given a partial lift
$\beta^{(<r)} \in \mathfrak{h}_T^\Theta / F^r$
solving the MC equation modulo~$F^r$, the
\emph{order-$r$ holographic obstruction} is
\begin{equation}\label{eq:holographic-obstruction}
\mathfrak{o}_r(\beta)
\;:=\;
\pi_r\!\Bigl(
\sum_{n \geq 1}\frac{1}{n!}\,
\ell_n^\Theta(\beta^{(<r)},\ldots,\beta^{(<r)})
\Bigr)
\;\in\;
Z^2\!\bigl(\operatorname{gr}^F_r
\mathfrak{h}_T^\Theta\bigr),
\end{equation}
and its cohomology class
$[\mathfrak{o}_r(\beta)]
\in H^2(\operatorname{gr}^F_r \mathfrak{h}_T^\Theta)$
vanishes if and only if the lift extends to the next
stage. The stages have the following interpretation\textup:
\begin{enumerate}[label=\textup{(\roman*)}]
\item stage~$1$ $=$ line-side deformations with fixed
 boundary\textup;
\item mixed stages $=$ bulk/boundary comparison
 data\textup;
\item higher stages $=$ nonlinear/genus compatibility.
\end{enumerate}
\end{definition}

\begin{definition}[Relative quartic obstruction]%
\label{def:relative-quartic-obstruction}%
\index{relative quartic obstruction|textbf}%
\index{quartic resonance class!relative}%
The \emph{relative quartic obstruction} is the
class at color depth~$4$\textup:
\begin{equation}\label{eq:relative-quartic}
\mathfrak{R}_4^{\mathrm{rel}}(T)
\;:=\;
[\mathfrak{o}_4(\beta)]
\;\in\;
H^2\!\bigl(\operatorname{gr}^F_4
\mathfrak{h}_T^\Theta\bigr).
\end{equation}
It measures the failure of binary spectral data
$r(z)$ together with bulk/boundary comparison data
to extend to a full mixed open/closed nonlinear
system\textup: the first invariant detecting when
a boundary modular datum cannot arise from a planar
holographic realization.
\end{definition}

\begin{conjecture}[Relative quartic computation;
\ClaimStatusConjectured]%
\label{conj:relative-quartic-computation}%
\index{relative quartic obstruction!computation}%
\textup{[}Type signature:
$(\mathrm{Open}\cap\mathrm{Cat},\mathrm{factorisation},
\mathrm{level}=3\to 4,
\mathrm{hypothesis\ package}=\mathrm{relative\ holographic\ deformation\ complex\ \textup{(}Definition~\ref{def:relative-holographic-deformation}\textup{)}};
\mathrm{five\text{-}archetype\ classification\ G/L/C/M\ on\ KSDual\ sublocus};
\mathrm{KZ\ analytic\ SDR\ +\ weight\text{-}completed\ ambient\ for\ class\text{-}M\ vanishing\ \textup{(}MA\text{-}13\textup{)}})$\textup{]}
\begin{enumerate}[label=\textup{(\roman*)}]
\item For the Heisenberg family,
 $\mathfrak{R}_4^{\mathrm{rel}}(\cH) = 0$.
\item For the affine family at generic level,
 $\mathfrak{R}_4^{\mathrm{rel}}(V_k(\fg)) = 0$
 or reduces to CYBE/KZ compatibility.
\item For the Virasoro and principal~$\cW_N$ at
 generic~$c$,
 $\mathfrak{R}_4^{\mathrm{rel}} \neq 0$:
 the quartic resonance class
 $\mathfrak{R}_{4,g,n}^{\mathrm{mod}}(\cA)$
 lifts to the first nontrivial relative class.
\end{enumerate}
\end{conjecture}

\begin{conjecture}[Loop--Connes transfer;
\ClaimStatusConjectured]%
\label{conj:loop-connes-transfer}%
\index{Loop--Connes principle|textbf}%
\index{genus-raising!as boundary trace image}%
\textup{[}Type signature:
$(\mathrm{Open}\cap\mathrm{Cat},\mathrm{factorisation},
\mathrm{level}=3\to 4,
\mathrm{hypothesis\ package}=\mathrm{chiral\ Hochschild\ derived\ centre\ on\ KSDual\ sublocus};
\mathrm{open\ BV\ operator\ }\Delta_{\mathrm{open}}\mathrm{\ defined\ on\ cyclic\ chains};
\mathrm{closed\ genus\text{-}raising\ }\Lambda_P\mathrm{\ realised\ via\ trace\ +\ clutching\ on\ open\ category},\ \mathrm{not\ a\ closed\text{-}algebra\ property})$\textup{]}
The derived-center map
$\beta_{\mathrm{der}}^*\colon
\operatorname{HH}^*_{\mathrm{ch}}(\cA^!)
\to Z_{\mathrm{der}}(\cA^!)$
sends the open BV operator~$\Delta_{\mathrm{open}}$
to the closed genus-raising operator:
\begin{equation}\label{eq:loop-connes}
\beta_{\mathrm{der}}^*(\Delta_{\mathrm{open}})
\;=\;
\Lambda_P\,.
\end{equation}
At the shadow level\textup: closed genus creation is the
image of the open cyclic trace under the derived-center
map. In the language of
Theorem~\textup{\ref{thm:nms-genus-loop-model-families}}\textup:
loops in the bulk are traces on the boundary.
\end{conjecture}

\subsection{Reduction table for the eight frontier conjectures}
\label{subsec:frontier-reduction-table}

\begin{remark}[Obstruction anatomy]
\label{rem:frontier-obstruction-anatomy}
\index{frontier conjectures!reduction table}
Each of the eight conjectures in this section is reduced to a
precisely identified minimal obstruction. The algebraic content
of every conjecture is proved; what remains is in each case either a
3d geometric construction, a categorical lift, or a modular-category
extension of a genus-$0$ adjunction. The Vol~III stage-two
specialisation
$\Phi_d^{(\Sigma_{d-1},C)}=\mathrm{Sp}^{\mathrm{ch}}_{\Sigma_{d-1},C}\circ
\Phi_d^{\mathrm{FA}}$
\textup{(}Convention~\textup{\ref{conv:master-type-signature}}\textup{)}
specialises Vol~III data on the curve $C$ to a
single $E_1$-chiral shadow on the boundary; that specialisation, not
inversion, is the cross-volume comparison map between Vol~I bar-cobar
data and Vol~III six-dimensional hCS data
\textup{(}master-concordance Theorem~A surface\textup{)}.
\begin{center}
\renewcommand{\arraystretch}{1.25}
\small
\begin{tabular}{@{}p{3.7cm}p{3.8cm}p{4.5cm}@{}}
\toprule
\textbf{Conjecture} & \textbf{Proved content}
 & \textbf{Minimal obstruction} \\
\midrule
Bulk enhancement
 \textup{(\ref{conj:frontier-bulk-enhancement-universal-mc})}
 & $\Theta_\cA$ proved; items (ii)--(iv) algebraic
 & Construct 3d bulk fact.\ algebra \\[4pt]
Boundary--defect
 \textup{(\ref{conj:boundary-defect-realization})}
 & Genus-$0$ identification $= $ Koszul duality
 & Geometric KK-reduction $=$ algebraic $\cA^!$ at modular level \\[4pt]
Singular-fiber descent
 \textup{(\ref{conj:singular-fiber-descent})}
 & Generic family modular Koszul
 & Coderived specialization at rational points \\[4pt]
HT deformation quant.\
 \textup{(\ref{conj:ht-deformation-quantization})}
 & Formal algebraic functor proved \textup{(Vol II)}
 & 3d Poisson sigma-model realization \\[4pt]
Subregular $\mathfrak{sl}_3$ datum
 \textup{(\ref{comp:subregular-sl3-holographic-datum})}
 & $\cA$, $\kappa$, shadow depth computed
 & $\cA^!$ identification
 $\Leftarrow$ \textup{(\ref{conj:type-a-transport-to-transpose-frontier})} \\[4pt]
Modular envelope
 \textup{(\ref{conj:universal-modular-factorization-envelope})}
 & Genus-$0$ Nishinaka adjunction; shadow functor
 & Modular completion of adjunction \\[4pt]
Analytic realization
 \textup{(\ref{conj:analytic-realization-criterion})}
 & HS-sewing for standard landscape; Heisenberg sewing
 & Extension beyond Lie-theoretic families \\[4pt]
Type-$A$ transport
 \textup{(\ref{conj:type-a-transport-to-transpose-frontier})}
 & Hook-type seed; propagation lemma
 & Edge-compatibility: Koszul commutes with DS \\
\bottomrule
\end{tabular}
\end{center}
\end{remark}

% ======================================================================
%
% STRING FIELD THEORY VERTICES FROM THE BAR COMPLEX
%
% ======================================================================

\subsection{String field theory vertices from the bar complex}
\label{subsec:frontier-sft-vertices}
\index{string field theory!vertices from bar complex|textbf}
\index{bar complex!string field theory vertices}

The bar complex $B(\cA)$ of a chiral algebra~$\cA$ encodes
string field theory vertices at the algebraic level.

\begin{remark}[SFT vertex comparison; \ClaimStatusHeuristic]

\label{rem:frontier-sft-vertices}
\index{string field theory!vertex identification|textbf}
The following comparisons are proposed after a conformal-overlap
comparison between the Ran bar construction and the geometric vertices
of string field theory:
\begin{enumerate}[label=\textup{(\roman*)}]
\item \emph{Open string.}
 The bar differential $d_B$ encodes the BRST
 operator~$Q$ of open string field theory. The bar
 coproduct $\Delta_B \colon B(\cA) \to B(\cA) \otimes B(\cA)$
 encodes the cubic open string vertex~$V_3$
 \textup{(}Witten~\cite{Witten86SFT}\textup{)}.
 The higher $A_\infty$ operations
 $m_n$ on $B(\cA)$ encode the $(n{+}1)$-point open
 string vertices~$V_{n+1}$.

\item \emph{Closed string.}
 The genus-$g$, degree-$n$ shadow amplitudes
 $\mathrm{Sh}_{g,n}(\Theta_\cA)$
 map to the closed string field theory vertices
 $V_{g,n}$ of Zwiebach~\cite{Zwiebach93} under a comparison morphism:
 \begin{equation}\label{eq:frontier-csft-vertex}
 \rho_{\mathrm{SFT}}(V_{g,n})
 \;=\;
 \mathrm{Sh}_{g,n}(\Theta_\cA).
 \end{equation}
 In particular: the tadpole $V_{1,1} = \kappa/24$
 is the genus-$1$ shadow at degree~$1$.

\item \emph{Master equation.}
 The closed SFT master equation
 \[
 \sum_{g,n}
 \frac{\hbar^g}{n!}\,
 V_{g,n}(\Psi^{\otimes n})
 \;=\; 0
 \]
 is the MC equation for
 $\Theta_\cA$
 \textup{(}Theorem~\textup{\ref{thm:mc2-bar-intrinsic}}\textup{)}.
\end{enumerate}

The algebraic bar complex structure is proved; the
identification with the specific geometric SFT vertices of
Witten and Zwiebach is structural, requiring the matching of
the bar construction on $\mathrm{Ran}(X)$ with the conformal-mapping
geometry of the SFT overlap conditions.

Part~(i): the bar construction on a curve $X = \mathbb{P}^1$
extracts collision residues from the OPE. The bar
differential $d_B$ acts by the genus-$0$ collision
operation on $\mathrm{FM}_k(X)$, which is the BRST
operator of the open string. The coproduct~$\Delta_B$
deconcatenates the bar elements, which on the
conformal-mapping side reduces to the midpoint overlap
condition of Witten's cubic vertex. The Neumann
coefficients $N^{rs}_{mn}$ parametrize this overlap:
$N^{11}_{mn} = (-1)^{m+n+1}(2/(m+n))
\binom{2m}{m}4^{-m}\binom{2n}{n}4^{-n}\sqrt{mn}$
(Gross--Jevicki~\cite{Gross-Jevicki87},
Rastelli--Sen--Zwiebach~\cite{RSZ01}).

Part~(ii): the shadow amplitudes
$\mathrm{Sh}_{g,n}(\Theta_\cA)$
are the $n$-point, genus-$g$ projections of the
universal MC element. Their identification with
Zwiebach's vertices follows from the geometric
construction: $V_{g,n}$ is defined as the
integral over $\mathcal{M}_{g,n}$ of the appropriate
form built from OPE data, which is precisely the
shadow amplitude at genus~$g$, degree~$n$.

Part~(iii): the MC equation
$D_\cA\Theta_\cA + \tfrac{1}{2}[\Theta_\cA, \Theta_\cA] = 0$
(Theorem~\ref{thm:mc2-bar-intrinsic})
expands genus-by-genus to the closed SFT equation.

The comparison
$\rho_{\mathrm{SFT}}(V_{g,n}) = \mathrm{Sh}_{g,n}(\Theta_\cA)$
would follow from matching the modular operad structure: the CSFT
vertices are the string amplitudes on $\overline{\cM}_{g,n}$, while
the MC element encodes the algebraic shadow amplitudes through the
bar-intrinsic construction
(Theorem~\ref{thm:mc2-bar-intrinsic}). The matching map
$\rho_{\mathrm{SFT}}$ is the open geometric input.
The convergent shadow partition function
$\sum_{g \geq 1} F_g\,\hbar^{2g}$
exhibits Bernoulli decay $F_{g+1}/F_g \to 1/(2\pi)^2$
(Theorem~\ref{thm:shadow-double-convergence}),
which stands in contrast to the Weil--Petersson
divergent string perturbation series
whose coefficients grow as $(2g)!$
(Shenker, 1990).
If the comparison map preserves the relevant amplitudes, the bar
complex provides a convergent algebraic shadow of the closed string
field theory: the shadow partition function converges where the
Weil--Petersson perturbation series diverges.
\end{remark}

\begin{remark}[Tachyon potential from the bar coproduct]
\label{rem:frontier-tachyon}
\index{tachyon potential!bar coproduct}
The open string tachyon potential
$V(T) = \sum_n V_{0,n}\,T^n/n!$ satisfies
Sen's conjecture $V(T_0)/T_D = -1$ at the stable
vacuum \cite{Sen99}. Level truncation of the bar
coproduct reproduces the convergent sequence:
$68.4\%$ at level $(0,0)$, $95.9\%$ at level $(2,4)$,
$99.93\%$ at level $(2,6)$. The bar complex provides
the off-shell algebraic formalism in which Sen's
conjecture is the statement that the MC element at
genus~$0$ has a unique stable critical point.
\end{remark}

\begin{remark}[Positive geometry of the bar complex: FM boundary and amplituhedron]
\label{rem:frontier-positive-geometry-bar}
\ClaimStatusConjectured
\index{positive geometry!bar complex|textbf}
\index{associahedron!ABHY realization}
\index{Fulton--MacPherson!positive geometry}
The genus-$0$ bar complex of a free (abelian) chiral algebra
admits a geometric realization as the ABHY associahedron
(Arkani-Hamed--Bai--He--Yan, 2017).
Three structures align:
\begin{enumerate}[label=\textup{(\roman*)}]
\item \emph{Combinatorial.}
 The bar differential on
 $B(\cA) = T^c(s^{-1}\bar\cA)$
 is a signed sum over planar binary trees; each summand
 corresponds to a facet of the Stasheff
 associahedron~$K_{n-1}$
 (Theorem~\ref{thm:fm-associahedron}).
 The ABHY construction endows $K_{n-1}$ with a convex
 realization in kinematic space whose canonical form
 $\Omega(K_{n-1})$ computes the tree-level bi-adjoint
 scalar amplitude $m_n(\mathbb{I} | \mathbb{I})$.

\item \emph{Geometric.}
 The Fulton--MacPherson compactification
 $\FM_n(X)$ (the iterated blowup of~$X^n$ along all
 diagonals) has boundary strata indexed by stable
 trees with the same face poset as~$K_{n-1}$.
 The bar propagators $d\log E(z_i, z_j)$
 (weight~$1$; see \S\ref{subsec:frontier-sft-vertices})
 furnish the canonical differential forms on these
 boundary strata: each codimension-$k$ boundary
 component of $\FM_n(X)$ contributes a $k$-fold
 residue of the product $\prod_{(i,j) \in T} d\log E(z_i, z_j)$
 indexed by the internal edges of the tree~$T$.

\item \emph{Color decoration via the R-matrix.}
 Passing from the free theory to an interacting
 chiral algebra~$\cA$ at level~$k$ dresses the
 bar-complex combinatorics with the classical
 r-matrix $r(z) = k\Omega/z$
 (which vanishes at $k = 0$, recovering the free case).
 The ordered bar $B^{\mathrm{ord}}(\cA)$ with its
 R-matrix twist encodes color-dressed amplitudes;
 the symmetric bar $B^{\Sigma}(\cA)$
 is the $\Sigma_n$-coinvariant shadow, which
 retains the scalar invariant
 $\kappa$, namely $\mathrm{av}(r(z))$ in abelian and scalar
 families and $\mathrm{av}(r(z))+\dim(\fg)/2$ for non-abelian
 affine Kac--Moody, but loses the color ordering.
 The R-matrix thus plays the role of the
 color-kinematic dressing that distinguishes
 Yang--Mills amplitudes from bi-adjoint scalars.
\end{enumerate}

At genus~$0$, the ABHY associahedron is the positive
geometry whose boundary structure \emph{is} the bar
differential of the free theory: the facets of~$K_{n-1}$
are the summands of~$d_B$, and the canonical form is the
tree-level amplitude.
Whether this identification extends to a full positive
geometry for the interacting bar complex (with R-matrix
twist encoding the color-dressed canonical form) and to
higher genus (where the relevant polytopes would be
the graph complexes on $\overline{\cM}_{g,n}$) remains
open.
\end{remark}

% ======================================================================
%
% CELESTIAL OPE FROM KOSZUL DUALITY
%
% ======================================================================

\subsection{Celestial OPE from Koszul duality}
\label{subsec:frontier-celestial}
\index{celestial holography!OPE from Koszul duality|textbf}
\index{collinear algebra!bar complex}

Celestial holography recasts four-dimensional scattering
amplitudes as correlators of a two-dimensional celestial
CFT on the celestial sphere~$S^2$. The collinear
singularities of massless particles define a chiral algebra
whose bar complex controls celestial amplitudes.

\begin{proposition}[Celestial OPE from the bar complex;
\ClaimStatusConditional]
\label{prop:frontier-celestial-ope}
\index{celestial OPE|textbf}
\textup{[}Type signature:
$(\mathrm{Open}\cap\mathrm{Cat},\mathrm{factorisation},
\mathrm{level}=1\to 2,
\mathrm{hypothesis\ package}=\mathrm{collinear\text{-}algebra\ comparison};
\mathrm{quadratic\ collinear\ OPE\ for\ part\ (i)};
\mathrm{PRS}\,\mathfrak{w}_{1+\infty}\mathrm{\ wedge\ closure\ for\ parts\ (ii)\text{--}(iii)};
\mathrm{quadratic\ OPE\ alone\ does\ not\ give\ all\text{-}loop\ HT:\ full\ celestial\ recovery\ requires\ Yamada\ finite\ presentation\ +\ KZ\ analytic\ SDR};
\mathrm{Vol\ III\ }(\Sigma_2,C)\text{-specialisation\ at\ }d=3,\ C\ \mathrm{the\ celestial\ sphere})$\textup{]}
Assume a collinear-algebra comparison map from the celestial OPE of
the four-dimensional theory to the chiral bar complex on the celestial
sphere; that the gauge collinear OPE has at most a single simple pole
\textup{(}quadratic OPE structure, with the Koszul dual
$\mathrm{Com}^!=\mathrm{Lie}$ identification dependent on the
single-pole hypothesis\textup{)}; and that the gravitational collinear
algebra realises $\mathfrak{w}_{1+\infty}[\lambda{=}1]$ in the
Pope--Romans--Shen \textup{(PRS)} wedge presentation. Then the
Vol~I bar formalism gives the following boundary shadow statements.
\begin{enumerate}[label=\textup{(\roman*)}]
\item \emph{Gauge theory.}
 For self-dual Yang--Mills with gauge group~$G$, the
 collinear algebra is the holomorphic current algebra
 $J_a(z)$ at level $k = 0$. The bar complex
 $B(J_{\mathrm{col}})$ is chirally Koszul with
 $H^*(B(J_{\mathrm{col}}))$ concentrated in bar degree~$1$.
 The Koszul dual is $\mathrm{Com}^! = \mathrm{Lie}$:
 the collinear gluon algebra is Koszul dual to the
 free Lie algebra on the adjoint representation.
 The collision residue is
 $r_{ab} = f^c_{ab}\,J_c$ The celestial R-matrix is the exponential of this
 constant r-matrix acting on tensor products of
 collinear modules.

\item \emph{Gravity.}
 For self-dual gravity, the collinear algebra is
 $\mathfrak{w}_{1+\infty}[\lambda = 1]$, the wedge subalgebra
 of $W_{1+\infty}$ at the self-dual coupling.
 This contains the BMS algebra (supertranslations $+$
 superrotations) and extends to an infinite tower of
 soft graviton symmetries.

\item \emph{Conformally soft theorems.}
 The image of the soft graviton theorem at sub${}^n$-leading order
 under the collinear comparison map is a bar cohomology class of the
 $\mathfrak{w}_{1+\infty}$ bar complex at
 $\Delta = 1 - n$.
 The hierarchy: $S_0$ (supertranslation, BMS),
 $S_1$ (superrotation, Virasoro extension),
 $S_2$ ($\mathfrak{w}_{1+\infty}$ extension,
 spin-$3$ soft theorem).
\end{enumerate}
\end{proposition}

\begin{proof}
Part~(i): after the collinear comparison has identified the celestial
gauge OPE with the holomorphic current algebra, the relation
$J_a(z)J_b(w) \sim
f^c_{ab}J_c(w)/(z{-}w)$ has a single simple pole.
The bar differential extracts this via $d\log(z{-}w)$,
giving the Chevalley--Eilenberg differential of the
loop algebra. Koszulness follows from the quadratic
OPE (single simple pole implies bar cohomology concentrated
in degree~$1$).

Part~(ii): after the gravitational collinear comparison, the
$\mathfrak{w}_{1+\infty}$ algebra has
generators $w_s(z)$ of spin $s = 1, 2, 3, \ldots$ with
OPE poles of order up to $2s$. The T-line
(Virasoro sub-tower) has shadow data identical to
$\mathrm{Vir}_c$ at the appropriate central charge.

Part~(iii): the conditional bar statement is that the comparison image
of the soft theorem at order~$n$ is represented by a nonzero class in
$H^n(B(\mathfrak{w}_{1+\infty}))$ at conformal dimension
$\Delta = 1 - n$. This follows from the
$\mathfrak{w}_{1+\infty}$ OPE and the resulting spectral sequence on
the bar complex once the comparison map has been supplied.
Computation:
\texttt{compute/lib/celestial\_koszul\_ope.py}.
\end{proof}

% ======================================================================
%
% SEIBERG-WITTEN PREPOTENTIAL FROM THE SHADOW CONNECTION
%
% ======================================================================

\subsection{Seiberg--Witten theory from the shadow connection}
\label{subsec:frontier-seiberg-witten}
\index{Seiberg--Witten!shadow connection|textbf}

\begin{proposition}[Shadow connection and Picard--Fuchs;
\ClaimStatusConditional]
\label{prop:frontier-sw-shadow}
\index{Picard--Fuchs equation!shadow connection}
\textup{[}Type signature:
$(\mathrm{Open}\cap\mathrm{Cat},\mathrm{factorisation},
\mathrm{level}=4\to 5,
\mathrm{hypothesis\ package}=\mathrm{Theorem~\ref{thm:shadow-connection-kz}};
\mathrm{Picard\text{--}Fuchs\ Fuchsian\ ODE\ comparison};
\mathrm{KSDual\ sublocus\ for\ Verdier\text{-}consistent\ monodromy};
\mathrm{Virasoro\ scope\ for\ part~(i);\beyond\ this\ requires\ Prochazka\ }\Walg_N\ \mathrm{closure\ for\ analogous\ statement};
\mathrm{Vol\ III\ }(\Sigma_2,C)\text{-specialisation\ for\ stage\text{-}2\ shadow})$\textup{]}
The shadow connection
$\nabla^{\mathrm{sh}} = d - Q'_L/(2Q_L)\,dt$
\textup{(}Theorem~\textup{\ref{thm:shadow-connection}}\textup{)}
is a logarithmic connection with regular singular points
at the zeros of~$Q_L$ and at infinity.
For the Virasoro algebra, $Q_L$ is quadratic in~$t$
with three regular singular points on~$\mathbb{P}^1$,
yielding a Gauss hypergeometric equation.
The identification with the Seiberg--Witten Picard--Fuchs
equation is:
\begin{center}
\renewcommand{\arraystretch}{1.25}
\begin{tabular}{l@{\quad$\longleftrightarrow$\quad}l}
\textbf{Shadow obstruction tower} & \textbf{Seiberg--Witten} \\
\midrule
$Q_L(t, c)$ & discriminant of the SW curve \\
$\nabla^{\mathrm{sh}}$ & Gauss--Manin connection \\
Koszul sign $-1$ & monodromy around singularities \\
$\kappa = c/2$ & prepotential coefficient \\
zeros of $Q_L$ & singular fibers $u = \pm\Lambda^2$
\end{tabular}
\end{center}
\end{proposition}

\begin{proof}
The shadow connection on a one-dimensional primary slice
has connection form $-Q'_L/(2Q_L)$, whose poles are at the
zeros of $Q_L$. For Virasoro,
$Q_L = (2\kappa + 3\alpha t)^2 + 2\Delta t^2$
is quadratic, so the connection has two finite singular
points and one at infinity: a rigid Fuchsian system.
The monodromy around each zero is $-1$
(Theorem~\ref{thm:shadow-connection}: Koszul sign).

For the $\mathrm{SU}(2)$ $N_f = 0$ Seiberg--Witten curve
$y^2 = (x^2 - u)^2 - \Lambda^4$, the Picard--Fuchs equation
for the period integrals is a second-order Fuchsian ODE with
singular points at $u = \pm\Lambda^2$ and $u = \infty$. The
monodromy matrices $M_\pm$ and $M_\infty$ satisfy
$M_+ M_- M_\infty = 1$. The identification maps the shadow
metric zeros to the singular fibers, and the Koszul monodromy
$-1$ to the $(-1)$-eigenvalue of the monodromy matrices.

The instanton coefficients of the prepotential
$F_{\mathrm{inst}}(a)
= \sum_{k \ge 1} F_k\,\Lambda^{4k}/a^{4k}$
with $F_1 = 1/(2^5\pi^2)$, $F_2 = 5/(2^{14}\pi^4)$
are controlled by the Taylor coefficients of the
shadow connection's flat sections.
Computation:
\texttt{compute/lib/seiberg\_witten\_shadow.py}.
\end{proof}

% ======================================================================
%
% CHERN-SIMONS FROM BAR-COBAR
%
% ======================================================================

\subsection{Chern--Simons partition functions from bar-cobar duality}
\label{subsec:frontier-chern-simons}
\index{Chern--Simons theory!bar-cobar|textbf}

\begin{proposition}[Chern--Simons/WZW comparison from the shadow obstruction tower;
\ClaimStatusConditional]
\label{prop:frontier-cs-shadow}
\index{Chern--Simons theory!shadow tower}
\index{Verlinde formula!shadow tower}
\textup{[}Type signature:
$(\mathrm{Open}\cap\mathrm{Cat},\mathrm{factorisation},
\mathrm{level}=4\to 5,
\mathrm{hypothesis\ package}=\mathrm{positive\ integral\ level\ }k\in\Z_{>0};
\mathrm{finite\ semisimplicity\ of\ MTC\ \textup{(}finite\text{-}spin\ Reshetikhin\text{--}Turaev\ endpoint\textup{)}};
\mathrm{comparison\ map\ from\ CS\ state\ space\ to\ factorisation\ homology};
\mathrm{KSDual\ sublocus\ for\ self\text{-}adjoint\ Verlinde\ identification\ at\ }k=k^\vee;
\mathrm{Vol\ III\ }(\Sigma_2,C)\text{-specialisation\ at\ }d=3;
\mathrm{modularity\ realised\ via\ }\Tr_\cC+\mathrm{clutching\ on\ open\ category})$\textup{]}
Assume the WZW/Chern--Simons correspondence at positive integral
level, the finite semisimplicity of the corresponding modular tensor
category, and a chosen comparison map from the Chern--Simons state
space to the Vol~I factorization homology object. Then the affine
boundary algebra supplies the following typed comparison.
For the affine algebra
$\widehat{\mathfrak{g}}_k$ at level~$k$
with modular characteristic
$\kappa(\widehat{\mathfrak{g}}_k)
= \dim(\mathfrak{g})(k + h^\vee)/(2h^\vee)$:
\begin{enumerate}[label=\textup{(\roman*)}]
\item \emph{$S^3$ partition function.}
 The Chern--Simons partition function on~$S^3$ is
 \[
 Z(S^3, G, k)
 \;=\;
 \sqrt{\frac{2}{k{+}h^\vee}}^{\,\mathrm{rank}(\mathfrak{g})}
 \cdot
 \prod_{\alpha > 0}
 2\sin\!\Bigl(
 \frac{\pi\langle\alpha, \rho\rangle}{k{+}h^\vee}
 \Bigr).
 \]
 For $\mathrm{SU}(2)$: $Z = \sqrt{2/(k{+}2)}\sin(\pi/(k{+}2))$.

\item \emph{Verlinde formula.}
 The dimension of the space of conformal blocks at
 genus~$g$ is
 $\dim V_g(G, k)
 = \sum_\lambda S_{0\lambda}^{2-2g}$,
 where $S_{0\lambda}$ is the modular $S$-matrix.
 The comparison map identifies this with the degree-zero
 projection of the factorization homology object
 $\mathbf{Q}_g(\widehat{\mathfrak{g}}_k)$ in the shadow obstruction
 tower.

\item \emph{Genus-$1$ scalar and one-loop determinant.}
 The genus-$1$ shadow amplitude
 \[
 F_1^{\mathrm{sh}}
 =
 \kappa/24
 =
 \frac{\dim(\mathfrak g)(k+h^\vee)}{48h^\vee}
 \]
 is the Vol~I scalar curvature coefficient. It is not the logarithm
 of the $S^3$ one-loop determinant. For
 $G=\mathrm{SU}(2)$, the perturbative determinant contains the
 logarithmic term $-\tfrac32\log(k+2)$, whereas
 $F_1^{\mathrm{sh}}=3(k+2)/96$ is linear in the shifted level.
 Any equality between these quantities requires an additional
 renormalized determinant comparison, not supplied by Theorem~D.
\end{enumerate}
\end{proposition}

\begin{proof}
The WZW/CS correspondence identifies the Hilbert space of
CS theory on~$\Sigma_g$ with the space of WZW conformal
blocks. The modular characteristic
$\kappa = \dim(\mathfrak{g})(k{+}h^\vee)/(2h^\vee)$
determines only the Vol~I genus-$1$ scalar. The Verlinde
formula computes the dimension of conformal blocks; under the
chosen comparison map this is the degree-$0$ projection of the
shadow complex $\mathbf{Q}_g$. The $S^3$ partition function is a
three-dimensional TQFT invariant controlled by quantum dimensions,
and its perturbative determinant is compared to the boundary scalar
only after the extra determinant map has been constructed.
Computation:
\texttt{compute/lib/chern\_simons\_barcobar.py}.
\end{proof}

\begin{remark}[AGT as factorization homology]

\label{rem:agt-factorization-homology}
\index{AGT correspondence!factorization homology interpretation}
\index{gauge origami!independent sum factorization}
The AGT correspondence~\cite{AGT09} identifies Nekrasov instanton
partition functions with $\mathcal{W}_N$ conformal blocks.
In the factorization formalism, both sides are shadows of the
same object: the factorization homology
$\int_X \cA$ of the boundary chiral algebra over the
curve.
Three structural parallels make this precise.
\begin{enumerate}[label=\textup{(\roman*)}]
\item \emph{Instanton number $=$ bar degree.}
The $k$-instanton sector of the Nekrasov partition function
$Z = \sum_{k \geq 0} q^k Z_k$ corresponds to the degree-$k$
bar complex sector $B^{(k)}(\cA)$.
The bar comultiplication (factorization splitting) dualizes
the CoHA multiplication (extension of sheaves on the
instanton moduli space).
\item \emph{Gauge origami $=$ independent sum factorization.}
The Nekrasov gauge origami construction places branes along
coordinate hyperplanes in~$\mathbb{C}^4$, with the total
partition function factoring as a product over spikes.
This matches
Proposition~\ref{prop:independent-sum-factorization}:
for $L = L_1 \oplus L_2$ with vanishing mixed OPE,
$\kappa$ is additive and $\Delta$ is multiplicative.
\item \emph{Genus expansion.}
The shadow genus expansion
$F_g = \kappa \cdot \lambda_g^{\mathrm{FP}}$
(Theorem~\ref{thm:modular-koszul-duality-main})
reproduces the vacuum sector of the Nekrasov genus expansion
$F = \sum_{g \geq 0} \hbar^{2g-2} F_g(a,q)$
in the unrefined limit.
\end{enumerate}
Computation:
\texttt{compute/lib/gauge\_origami\_comparison\_engine.py}
($94$~tests).
\end{remark}

% ======================================================================
%
% TWISTED HOLOGRAPHY AMPLITUDES
%
% ======================================================================

\subsection{Twisted holography amplitudes from Koszul duality}
\label{subsec:frontier-twisted-holography}
\index{twisted holography!amplitudes from Koszul duality|textbf}

\begin{theorem}[Twisted holography datum; \ClaimStatusConditional]
\label{thm:frontier-twisted-holography}
\index{twisted holography|textbf}
\index{D3 brane!holographic datum}
\textup{[}Type signature:
$(\mathrm{Open}\cap\mathrm{Cat},\mathrm{factorisation},
\mathrm{level}=1\to 5,
\mathrm{hypothesis\ package}=\mathrm{Costello\text{--}Li\ holomorphic\ twist};
\mathrm{Costello\text{--}Gaiotto\ boundary\ chiral\ algebra\ identification};
\mathrm{KSDual\ sublocus\ for\ }r(z)\mathrm{\ identification\ \textup{(}affine\ }\mathfrak{gl}_N\mathrm{\ at\ }k=1\textup{)}};
\mathrm{Vol\ III\ }(\Sigma_2,C)\text{-specialisation\ at\ }d=3\ \mathrm{for\ stage\text{-}2\ shadow};
\mathrm{modularity\ here\ is\ }\Tr_\cC+\mathrm{clutching\ on\ open\ category},\ \mathrm{not\ a\ closed\text{-}algebra\ property\ of\ }\cA)$\textup{]}
The Costello--Li twisted holography programme identifies
the holomorphic twist of type~IIB string theory on
$\mathrm{AdS}_3 \times S^3 \times T^4$ with Koszul
duality data attached to a boundary chiral algebra~$\cA$ and its
Verdier/Koszul branch~$\cA^!$. This is a typed boundary comparison,
not an equality of the physical bulk with~$\cA^!$. The holographic
modular Koszul datum
\[
\mathcal{H}(\mathcal{T})
\;=\;
(\cA,\; \cA^i,\; \cA^!,\; \mathcal{C},\; r(z),\; \Theta_\cA,\;
\nabla^{\mathrm{hol}})
\]
is computed for two brane systems:
\begin{enumerate}[label=\textup{(\roman*)}]
\item \emph{D3 brane \textup{(}twisted $\mathcal{N} = 4$ SYM\textup{)}.}
 The boundary chiral algebra is affine
 $\mathfrak{gl}_N$ at level $k = 1$:
 \[
 \kappa(\mathfrak{gl}_N,\, k{=}1)
 \;=\;
 \frac{(N^2{-}1)(N{+}1)}{2N} + 1.
 \]
 The holographic R-matrix is
 $r(z) = k\Omega/z$ (classical Yang--Baxter; at the D3 locus
 $k=1$ this specializes to the level-$1$ kernel, while the
 level-prefixed form $k\Omega/z$ vanishes at $k=0$; $\Omega$ the Casimir element), and the genus expansion
 $F_g = \kappa\cdot\lambda_g^{\mathrm{FP}}$ \textup{(}UNIFORM-WEIGHT\textup{}) computes
 twisted $\mathcal{N} = 4$ amplitudes at genus~$g$.
 The anomaly matching condition
 $\kappa_{\mathrm{eff}}
 = \kappa(\mathrm{matter}) + \kappa(\mathrm{ghost}) = 0$
 is a consistency check.

\item \emph{Sphere reconstruction.}
 The genus-$0$, $n$-point shadow amplitudes
 $\mathrm{Sh}_{0,n}(\Theta_\cA)$ are the Witten
 diagram amplitudes of the holomorphic twist. The
 commuting differentials of Gaiotto--Zinenko~\cite{GZ26}
 are the scalar shadow of $\mathrm{Sh}_{0,n}$.
\end{enumerate}
\end{theorem}

\begin{proof}
The boundary VOA of the holomorphic twist is identified
by Costello--Gaiotto~\cite{Costello-Gaiotto18}.
The modular characteristic is computed by the additivity
formula
$\kappa(\mathfrak{gl}_N) =
\kappa(\mathfrak{sl}_N) + \kappa(\mathfrak{u}(1))$
(Proposition~\ref{prop:independent-sum-factorization}),
with $\kappa(\mathfrak{sl}_N, k)
= (N^2{-}1)(k{+}N)/(2N)$ and
$\kappa(\mathfrak{u}(1), k) = k$
(the authoritative formulas from
Theorem~\ref{thm:modular-characteristic}).

The R-matrix
$r(z) = \mathrm{Res}^{\mathrm{coll}}_{0,2}(\Theta_\cA)
= k\,\Omega_{\mathrm{tr}}/z$
is the trace-form collision residue. The KZ-normalized representative
$\Omega/((k+h^\vee)z)$ is obtained by Casimir and coupling rescaling
\textup{(}at $k=0$ the trace-form collision residue vanishes; at the
D3 level $k=1$ this gives the level-$1$ specialization of the same
generic family\textup{)}:
the current algebra OPE
$J_a(z)J_b(w) \sim k\delta_{ab}/(z{-}w)^2 +
f^c_{ab}J_c/(z{-}w)$
has at most a double pole; the $d\log$ extraction
 gives a simple-pole r-matrix $r_{ab}(z) = k\Omega_{ab}/z$.

The sphere reconstruction uses the genus-$0$ shadow
amplitudes, which are the tree-level Witten diagrams
computed from the stable graph sum on~$\overline{\mathcal{M}}_{0,n}$.
Computation:
\texttt{compute/lib/twisted\_holography\_amplitudes.py}.
\end{proof}

% ======================================================================
%
% ABJM HOLOGRAPHIC DATUM
%
% ======================================================================

\subsection{The ABJM holographic modular Koszul datum}
\label{subsec:frontier-abjm}
\index{ABJM theory!holographic datum|textbf}
\index{M2 brane!holographic datum}

\begin{theorem}[ABJM holographic datum; \ClaimStatusConditional]
\label{thm:frontier-abjm}
\index{ABJM theory|textbf}
\textup{[}Type signature:
$(\mathrm{Open}\cap\mathrm{Cat},\mathrm{scalar},
\mathrm{level}=1\to 5,
\mathrm{hypothesis\ package}=\mathrm{Costello\text{--}Gaiotto\ boundary\ BRST};
\mathrm{Marino\text{--}Putrov\ +\ Fuji\text{--}Hirano\text{--}Moriyama\ localisation};
\mathrm{large\text{-}}N\mathrm{\ comparison\ map\ for\ Airy\ confluence};
\mathrm{KSDual\ sublocus\ irrelevant;\leading\text{-}lane\ }\kappa<0\ \mathrm{breaks\ self\text{-}adjointness};
\mathrm{modularity\ realised\ via\ }\Tr_\cC+\mathrm{clutching},\ \mathrm{not\ as\ closed\text{-}algebra\ Airy\ property};
\mathrm{Vol\ III}\,(\Sigma_2,C)\mathrm{\ specialisation\ at\ }d=3\ \mathrm{for\ stage\text{-}2\ shadow})$\textup{]}
Assume the Costello--Gaiotto identification of the holomorphic-twist
boundary BRST vertex algebra and the Fuji--Hirano--Moriyama /
Marino--Putrov localization formula for the physical ABJM partition
function. For ABJM theory with gauge group
$U(N)_k \times U(N)_{-k}$ and $4N^2$ symplectic
boson bifundamentals:
\begin{enumerate}[label=\textup{(\roman*)}]
\item \emph{Boundary VOA.}
 The boundary chiral algebra of the holomorphic twist is
 the BRST reduction
 \[
 \cA_{\mathrm{ABJM}}(N, k)
 \;=\;
 H^0_{\mathrm{BRST}}\!\bigl(
 V_k(\mathfrak{gl}_N)
 \otimes V_{-k}(\mathfrak{gl}_N)
 \otimes \mathrm{Sb}^{4N^2}
 \bigr),
 \]
 where $\mathrm{Sb}^{4N^2}$ denotes $4N^2$ symplectic
 bosons.

\item \emph{Pre-BRST and reduced modular characteristics.}
 By additivity, the modular characteristic receives
 contributions from the two Chern--Simons sectors
 and from the $4N^2$ symplectic boson pairs:
 \begin{equation}\label{eq:frontier-abjm-kappa}
 \kappa_{\mathrm{pre}}(\cA_{\mathrm{ABJM}}(N,k))
 \;=\;
 \underbrace{(N^2-1)}_{\text{CS}_{k}\,+\,\text{CS}_{-k}}
 \;+\;
 \underbrace{(-2N^2)}_{4N^2\;\text{Sb}_{1/2}}
 \;=\;
 -(N^2+1).
 \end{equation}
 The BRST-reduced compute model gives
 $\kappa_{\mathrm{red}}=-N^2$; the finite-$N$ equality between a
 reduced physical datum and the pre-BRST value is not asserted without
 computing the ghost cancellation in the same BRST complex.
 The negative sign reflects the non-unitarity of the
 symplectic boson system.
 At large~$N$, $\kappa \sim -N^2$.

\item \emph{Genus-$1$ free energy.}
 The pre-BRST scalar gives
 $F_1^{\mathrm{pre}} = \kappa_{\mathrm{pre}}/24 = -(N^2+1)/24$,
 while the reduced compute datum gives
 $F_1^{\mathrm{red}}=-N^2/24$.
 At large~$N$, the boundary scalar is
 $F_1^{\mathrm{red}}\sim -N^2/24$. This is the Vol~I genus-one
 shadow coefficient. It is not the physical Airy free energy, whose
 leading M2-brane scaling is $N^{3/2}$ after localization and
 grand-potential resummation.

\item \emph{Airy connection.}
 The shadow connection $\nabla^{\mathrm{sh}}$,
 after the rescaling $N \mapsto t$ and the large-$N$ limit,
 degenerates to the Airy connection
 $\nabla^{\mathrm{Airy}} = d/dt - \bigl(\begin{smallmatrix}
 0 & 1 \\ t & 0 \end{smallmatrix}\bigr)\,dt$
 on~$\mathbb{C}^2$.
 The ABJM partition function is
 \begin{equation}\label{eq:frontier-abjm-airy}
 Z(N, k)
 \;=\;
 C_k^{-1/3}\,
 \mathrm{Ai}\bigl(
 C_k^{-1/3}(N - B_k)
 \bigr),
 \end{equation}
 where $C_k = 2/(\pi^2 k)$,
 $B_k = k/24 - 1/(6k) + 1/3$,
 and $\mathrm{Ai}$ is the Airy function
 \textup{(}Fuji--Hirano--Moriyama~\cite{FHM11},
 Marino--Putrov~\cite{MP11}\textup{)}.
 The $N^{3/2}$ scaling of the free energy
 $F(N, k) \sim (\pi\sqrt{2k}/3)N^{3/2}$
 is the hallmark of M2-brane degrees of freedom.

\item \emph{Shadow--Airy correspondence.}
 The passage from the shadow connection
 (regular singularities, logarithmic) to the Airy
 connection (irregular singularity at infinity,
 Stokes phenomenon) is an irregular degeneration:
 a confluence of regular singular points in the
 large-$N$ limit. This is the shadow-tower
 manifestation of the $N^{3/2}$ scaling law.
\end{enumerate}
\end{theorem}

\begin{proof}
Part~(i): the holomorphic twist of ABJM identifies the
boundary VOA by Costello--Gaiotto~\cite{Costello-Gaiotto18}.
The BRST reduction produces the boundary algebra from the
CS and matter sectors.

Part~(ii): the modular characteristic is additive
(Proposition~\ref{prop:independent-sum-factorization}).
For the two CS sectors:
$\kappa(\mathfrak{gl}_N, k) + \kappa(\mathfrak{gl}_N, -k)
= [(N^2{-}1)(k{+}N)/(2N) + k]
+ [(N^2{-}1)(-k{+}N)/(2N) + (-k)] = (N^2{-}1)$.
For the $4N^2$ symplectic bosons
$\mathrm{Sb}_\lambda$ at $\lambda = 1/2$:
$\kappa(\mathrm{Sb}_{1/2})
= 6(1/2)^2 - 6(1/2) + 1 = -1/2$ per pair, so
$4N^2 \cdot (-1/2) = -2N^2$.
Total: $\kappa_{\mathrm{pre}} = (N^2{-}1) + (-2N^2) = -(N^2 + 1)$.
This computation uses the additivity of~$\kappa$
(Proposition~\ref{prop:independent-sum-factorization}) on the
pre-BRST system. The reduced datum
$\kappa_{\mathrm{red}}=-N^2$ is the value implemented by the compute
engine after BRST reduction. The precise finite-$N$ comparison between
the two requires a ghost-cancellation computation inside the same BRST
complex.

Part~(iii): immediate from the chosen pre-BRST or reduced scalar and
$\lambda_1^{\mathrm{FP}} = 1/24$.

Part~(iv): the Airy function formula is the
localization result of Fuji--Hirano--Moriyama and
Marino--Putrov. The identification with the shadow
connection is conditional on the large-$N$ comparison map from the
boundary scalar tower to the localization spectral curve:
the shadow metric $Q(t)$ at finite~$N$ has regular
singular points, which coalesce in the large-$N$ limit
to produce the irregular singularity of the Airy
equation $y'' = ty$.

Part~(v): the confluence of regular singularities
is a standard phenomenon in the theory of connections
(Katz~\cite{Katz90}). The shadow obstruction tower captures the
finite-$N$ boundary BRST scalar structure. The Airy limit extracts
the physical $N^{3/2}$ scaling only after the localization comparison
has been imposed.
Computation:
\texttt{compute/lib/abjm\_holographic\_datum.py}.
\end{proof}

\subsection{The M5 brane holographic modular Koszul datum}
\label{subsec:frontier-m5}
\index{M5 brane!holographic datum|textbf}
\index{twisted M-theory!M5 boundary algebra}

\begin{theorem}[M5 brane holographic datum; \ClaimStatusConditional]
\label{thm:frontier-m5}
\index{M5 brane|textbf}
\textup{[}Type signature:
$(\mathrm{Open}\cap\mathrm{Cat},\mathrm{factorisation},
\mathrm{level}=1\to 5,
\mathrm{hypothesis\ package}=\mathrm{Costello\ M5\ boundary\ identification};
\mathrm{Beem\text{--}Rastelli\text{--}van\ Rees\ protected\ chiral\ algebra};
\mathrm{Linshaw\text{--}Rapcak\ involution\ }\lambda\mapsto 2-\lambda;
\mathrm{stress\text{-}tensor\ lane\ \textup{(}finite\text{-}spin\ truncation\ at\ Virasoro\text{-}sub\text{-}tower\textup{)},\ full}\ W_{1+\infty}\ \mathrm{requires\ Prochazka\text{--}Rapcak\ closure};
\mathrm{KSDual\ sublocus\ at\ }N=1\ \mathrm{\textup{(}free\ tensor\textup{)}\ only};
\mathrm{Vol\ III}\,(\Sigma_5,C)\text{-specialisation\ at\ }d=6,\ C\ \mathrm{the\ M5\ boundary\ curve};
\mathrm{stress\text{-}lane\ scope\ for\ }\kappa;\full\ W_{1+\infty}\ \mathrm{shadow\ requires\ higher\text{-}spin\ extension})$\textup{]}
For $N$ M5 branes in twisted M-theory, the boundary chiral
algebra is $W_{1+\infty}[\lambda = N]$
\textup{(}Costello~\cite{Costello2111},
Costello--Paquette~\cite{CostelloP2201};
for the 11d uplift see
Bittleston--Costello~\cite{BittlestonCostello25}\textup{)}.
The holographic modular Koszul datum is:
\begin{enumerate}[label=\textup{(\roman*)}]
\item \emph{Boundary VOA.}
 $\cA_{\mathrm{M5}}(N) = W_{1+\infty}[\lambda = N]$,
 the $W_\infty$ algebra at parameter $\lambda = N$.
 At $N = 1$: the free tensor multiplet
 (symplectic boson $+$ free fermion);
 at $N \geq 2$: the interacting $(2,0)$ chiral algebra.

\item \emph{Central charge.}
 The 2d central charge of the protected chiral algebra
 is \textup{(}Beem--Rastelli--van Rees~\cite{BRvR14-6d}\textup{)}
 \begin{equation}\label{eq:frontier-m5-central-charge}
 c_{2\mathrm{d}}(\mathrm{M5}_N)
 \;=\;
 4N^3 - 3N - 1.
 \end{equation}
 At $N = 1$: $c = 0$ (free tensor multiplet, anomaly-free).
 At $N = 2$: $c = 25$.
 At $N = 3$: $c = 98$.

\item \emph{Stress-tensor modular characteristic.}
 Using the Virasoro-line extraction
 $\kappa_{\mathrm{stress}} = c_{2\mathrm{d}}/2$:
 \begin{equation}\label{eq:frontier-m5-kappa}
 \kappa_{\mathrm{stress}}(\mathrm{M5}_N)
 \;=\;
 \frac{4N^3 - 3N - 1}{2}.
 \end{equation}
 At large~$N$, $\kappa \sim 2N^3$: the cubic scaling
 is the shadow signature of the $N^3$ scaling of M5 brane
 degrees of freedom \textup{(}Henningson--Skenderis~\cite{HS98}\textup{)}.
 At $N = 1$: $\kappa_{\mathrm{stress}} = 0$ (free tensor multiplet).
 The full $W_{1+\infty}$ modular characteristic may receive
 higher-spin contributions; this theorem records the stress-tensor
 lane only.

\item \emph{Koszul dual.}
 The $W_\infty$ involution $\lambda \mapsto 2 - \lambda$ gives the
 dual parameter
 $\cA_{\mathrm{M5}}(N)^! = W_{1+\infty}[\lambda = 2 - N]$
 on the Linshaw--Rapcak lane. The full stress-tensor scalar is the
 formal continuation of~\eqref{eq:frontier-m5-central-charge}:
 \begin{equation}\label{eq:frontier-m5-kappa-dual}
 \kappa_{\mathrm{stress}}^{\mathrm{full}}(\mathrm{M5}_N^!)
 \;=\;
 \frac{4(2-N)^3-3(2-N)-1}{2}
 =
 \frac{-4N^3+24N^2-45N+25}{2}.
 \end{equation}
 The clean anomaly-leading lane instead uses
 $\kappa_{\mathrm{lead}}(\mathrm{M5}_N)=N^3$ and
 $\kappa_{\mathrm{lead}}(\mathrm{M5}_N^!)=(2-N)^3$.

\item \emph{Complementarity sum.}
 On the full stress-tensor lane,
 \begin{equation}\label{eq:frontier-m5-complementarity}
 \kappa_{\mathrm{stress}}(\mathrm{M5}_N) +
 \kappa_{\mathrm{stress}}^{\mathrm{full}}(\mathrm{M5}_N^!)
 \;=\;
 \frac{4N^3 - 3N - 1}{2}
 +
 \frac{4(2-N)^3-3(2-N)-1}{2}
 \;=\;
 12(N-1)^2,
 \end{equation}
 which vanishes exactly at the free tensor point $N=1$.
 At the leading-order level (matching the geometric $N^3$ anomaly
 normalization $\kappa_{\mathrm{lead}} = N^3$,
 $\kappa^!_{\mathrm{lead}} = (2-N)^3$):
 \begin{equation}\label{eq:frontier-m5-complementarity-leading}
 \kappa_{\mathrm{lead}} + \kappa^!_{\mathrm{lead}}
 \;=\;
 N^3 + (2-N)^3
 \;=\;
 6N^2 - 12N + 8
 \;=\;
 2\bigl(3(N-1)^2 + 1\bigr),
 \end{equation}
 which is always positive (minimum $2$ at $N = 1$).

\item \emph{Shadow class.}
 At $N = 1$: class C (free tensor, shadow depth $r_{\max} = 4$).
 At $N \geq 2$: class M (interacting $(2,0)$, shadow
 depth $r_{\max} = \infty$).

\item \emph{M2 vs.\ M5 dichotomy.}
 The M2 brane (ABJM, Theorem~\ref{thm:frontier-abjm})
 has $\kappa \sim -N^2$ with Airy $N^{3/2}$ scaling,
 while the M5 brane has $\kappa \sim +N^3$.
 The sign and scaling distinguish the two:
 \[
 \begin{array}{@{}l@{\;:\;\;}c@{\;\;}c@{\;\;}c@{}}
 & \kappa & \text{scaling} & \text{sign} \\
 \midrule
 \text{M2 (ABJM)} & -(N^2 + 1) & N^{3/2} & - \\
 \text{M5} & (4N^3 - 3N - 1)/2 & N^3 & +
 \end{array}
 \]
 The negative sign for M2 reflects the symplectic boson
 dominance; the positive sign for M5 reflects the
 gravitational anomaly of the $(2,0)$ theory.
\end{enumerate}
\end{theorem}

\begin{proof}
Part~(i): the identification of the M5 boundary VOA with
$W_{1+\infty}[\lambda = N]$ is the central result of
Costello~\cite{Costello2111} and
Costello--Paquette~\cite{CostelloP2201};
see also Costello--Paquette--Sharma~\cite{CPS2208,CPS2306}
for detailed constructions in twisted M-theory.

Part~(ii): the Beem--Rastelli--van Rees formula
$c_{2\mathrm{d}} = 4N^3 - 3N - 1$ for the protected chiral
algebra of the 6d $(2,0)$ $A_{N-1}$ theory is proved in
\cite{BRvR14-6d}. At $N = 1$, the formula gives $c = 0$
\textup{(}free tensor multiplet\textup{)}; at $N = 2$, it gives
$c = 25$.

Part~(iii): from $\kappa = c_{2\mathrm{d}}/2$ (the
Virasoro-line formula, Theorem~\textup{\ref{thm:genus-universality}}).

Part~(iv): the involution $\lambda \mapsto 2 - \lambda$ on
$W_{1+\infty}$ is the Linshaw--Rapcak duality. The full
stress-tensor scalar follows by applying the same central-charge
formula at $\lambda=2-N$; the leading anomaly scalar is recorded
separately and is not substituted into the full stress-tensor formula.

Part~(v): direct computation.

Part~(vi): at $N = 1$ the algebra is free (contact class);
at $N \geq 2$ the non-quadratic OPE of $W_{1+\infty}$
forces infinite shadow depth (class~M).

Part~(vii): comparison of
\eqref{eq:frontier-abjm-kappa} and \eqref{eq:frontier-m5-kappa}.
Computation:
\texttt{compute/lib/m5\_brane\_shadow\_engine.py}.
\end{proof}

\begin{proposition}[Leading phantom M5 Koszul dual; \ClaimStatusConditional]
\label{prop:phantom-m5-koszul-dual}
\index{M5 brane!Koszul dual|textbf}%
\index{phantom theory!M5 Koszul dual|textbf}%
\index{Koszul duality!phantom M5|textbf}%
\index{phantom algebra}%
\index{Koszul duality!M5 brane}%
\textup{[}Type signature:
$(\mathrm{Open}\cap\mathrm{Cat},\mathrm{factorisation},
\mathrm{level}=1\to 5,
\mathrm{hypothesis\ package}=\mathrm{Theorem~\ref{thm:frontier-m5}};
\mathrm{Linshaw\text{--}Rapcak\ involution};
\mathrm{KSDual\ sublocus\ only\ at\ }N=1;
\mathrm{leading\text{-}anomaly\ lane\ scope;\ full\ }W_{1+\infty}\ \mathrm{shadow\ requires\ Prochazka\text{--}Rapcak\ finite\text{-}spin\ closure})$\textup{]}
On the leading anomaly lane, the Koszul dual
$\mathrm{M5}_N^! = W_{1+\infty}[\lambda = 2 - N]$ is a
\emph{phantom} for $N \geq 3$: it satisfies
\[
\kappa_{\mathrm{lead}}(\mathrm{M5}_N^!) = (2-N)^3 < 0,
\]
so the leading scalar genus expansion
$F_g^{\mathrm{lead}}(\mathrm{M5}_N^!)
=\kappa_{\mathrm{lead}}^!\lambda_g^{\mathrm{FP}}$
\textup{(}UNIFORM-WEIGHT\textup{}) has the wrong sign. On the full
stress-tensor lane one must instead use
\[
\kappa_{\mathrm{stress}}^{\mathrm{full}}(\mathrm{M5}_N^!)
=\frac{4(2-N)^3-3(2-N)-1}{2},
\]
which is also negative for $N\geq3$. Neither scalar is an equality
with a physical dual theory without an additional bulk comparison.
\begin{enumerate}[label=\textup{(\roman*)}]
\item \emph{Irreducible anomaly.}
 The complementarity sum at leading order is
 \begin{equation}\label{eq:m5-complementarity-polynomial}
 \kappa_{\mathrm{lead}}(\mathrm{M5}_N)
 + \kappa_{\mathrm{lead}}(\mathrm{M5}_N^!)
 \;=\;
 N^3 + (2-N)^3
 \;=\;
 6N^2 - 12N + 8
 \;=\;
 2\bigl(3(N{-}1)^2 + 1\bigr),
 \end{equation}
 which is strictly positive for all~$N$ \textup{(}discriminant
 $144 - 192 = -48 < 0$; minimum value~$2$ at $N = 1$\textup{)}.
 In particular, the complementarity sum never vanishes: there is
 no anomaly cancellation between M5 and its Koszul dual at any~$N$.
 This is the leading-lane W-algebra-type complementarity of
 Theorem~\textup{\ref{thm:complementarity}}.

\item \emph{$N = 2$: leading uncurved point.}
 At $N = 2$: $\lambda^! = 0$ and
 $\kappa_{\mathrm{lead}}(\mathrm{M5}_2^!) = 0$, so the leading
 anomaly lane is uncurved. The full stress-lane continuation gives
 $\kappa_{\mathrm{stress}}^{\mathrm{full}}(\mathrm{M5}_2^!)=-1/2$,
 so no full stress-tensor vanishing is asserted. The dual algebra is
 $W_{1+\infty}[\lambda = 0] \cong \mathfrak{u}(\infty)_1$, the
 free boson tower (vacuum module of the affine Yangian). This is a
 leading-lane vertex-algebra realization, not a physical dual of the
 six-dimensional theory.

\item \emph{Exclusion of physical candidates.}
 For $N \geq 3$, the dual $\mathrm{M5}_N^!$ does not coincide with
 any known physical theory:
 \begin{itemize}
 \item \emph{Not a $6$d $(1,0)$ theory}: the $(1,0)$ anomaly
 coefficient $a_{(1,0)}$ is always positive
 \textup{(}Ohmori--Shimizu--Tachikawa--Yonekura~\cite{OSTY14}\textup{)},
 but $\kappa^! < 0$. The sign mismatch is absolute.
 \item \emph{Not the S-dual}: Koszul duality
 $\lambda \mapsto 2 - \lambda$ is a homological operation on the
 bar complex, not a duality of the physical theory. In every
 compactification dimension ($6$d, $5$d on~$S^1$, $4$d on~$T^2$,
 $3$d on~$T^3$), the Koszul involution differs from the physical
 duality (S-duality, T-duality, $3$d mirror symmetry).
 \end{itemize}
\end{enumerate}
\end{proposition}

\begin{proof}
The formula $\kappa_{\mathrm{lead}}(\mathrm{M5}_N^!) = (2 - N)^3$
is the leading anomaly scalar from
Theorem~\ref{thm:frontier-m5}. For~(i): the identity
$N^3 + (2-N)^3 = 6N^2 - 12N + 8$ is checked by direct expansion,
and the quadratic $6x^2 - 12x + 8$ has negative discriminant
$\Delta = (-12)^2 - 4 \cdot 6 \cdot 8 = -48$, so it has no real
roots. The full stress-lane scalar is negative for $N\geq3$ by the
formula in~\eqref{eq:frontier-m5-kappa-dual}. For~(ii):
$\kappa_{\mathrm{lead}}(\mathrm{M5}_2^!) = (2-2)^3 = 0$ and the
identification $W_{1+\infty}[\lambda = 0] \cong
\mathfrak{u}(\infty)_1$ is standard. For~(iii): the $(1,0)$
exclusion follows from positivity of $a_{(1,0)}$
\cite{OSTY14,BT15}; the S-duality exclusion follows from the
observation that the Koszul involution
$\lambda \mapsto 2 - \lambda$ differs from every physical duality
in the dimensional reduction tower ($6$d lacks a Lagrangian, $5$d
gives T-duality $R \mapsto 1/R$, $4$d gives Montonen--Olive
$\tau \mapsto -1/\tau$, $3$d gives mirror symmetry).

Computation:
\texttt{compute/lib/phantom\_m5\_koszul\_dual\_engine.py} (67~tests).
\end{proof}

\begin{remark}[Class-$\mathcal S$ $A_1$ on $\Sigma_{0,24}$ as the
$\mathcal N{=}2$ parent of the K3-coupled M5 datum]
\label{rem:frontier-class-S-K3-parent}
\ClaimStatusConjectured
\index{class-S theory!K3 parent}
\index{Chacaltana--Distler!A1 on Sigma024}
\index{Steiner system!S(5,8,24) rigidity}
\index{N=2 superconformal!4d parent}
The M5 boundary datum~$\mathcal{A}_{\mathrm{M5}}(2)
= W_{1+\infty}[\lambda{=}2]$ admits a 4d $\mathcal N{=}2$
parent in the form of a class-$\mathcal S$ theory of type~$A_1$
on the genus-zero Riemann surface $\Sigma_{0,24}$ with $24$ regular
punctures (one per $I_1$ Kodaira fibre on the K3-elliptic side of
Remark~\ref{rem:frontier-bi-based-ran-K3E}). The Chacaltana--Distler
construction packs this configuration into the following invariants
\textup{(}\cite{Arakawa17} for the chiral algebra of class~$\mathcal S$;
the genus-zero $A_1$ formulas
follow from the Riemann--Roch
counts of the BPZ degenerate-representation chiral algebra in the
two-dimensional reduction\textup{)}:
\begin{enumerate}[label=\textup{(\roman*)}]
\item \emph{Conformal anomalies.}
 $c_{4\mathrm d}(A_1, \Sigma_{0,24}, \mathrm{all\ max}) = 107/6$
 in the Shapere--Tachikawa normalisation
 $c_{4\mathrm d} = (2n_v + n_h)/12$
 \textup{(}Shapere--Tachikawa 2008 \S 3, arXiv:0804.1957\textup{)}
 applied to the Chacaltana--Distler class-$\mathcal S$ count
 $(n_v, n_h) = (63, 88)$
 \textup{(}Chacaltana--Distler 2010 eq.~(5.14)\textup{)};
 the protected chiral subalgebra of Beem et
 al.~\cite{Beem-etal-2014} has
 $c_{2\mathrm d}(A_1, \Sigma_{0,24}, \mathrm{all\ max}) = -214$
 \textup{(}standard $c_{2\mathrm d} = -12\,c_{4\mathrm d}$ ratio
 gives $-12\cdot 107/6 = -214$\textup{)}.
 These are the class-$\mathcal S$ parent values, not the stress
 central charge of the M5 boundary VOA in
 Theorem~\ref{thm:frontier-m5}.
\item \emph{Coulomb branch dimension.}
 $\dim_{\mathbb C}\operatorname{Coul}(A_1, \Sigma_{0,24}) = 21$,
 matching the dimension of the moduli space of $24$ unordered
 marked points on~$\mathbb{P}^1$ modulo $\operatorname{PSL}_2$
 \textup{(}$24 - 3 = 21$\textup{)}; each marked point contributes
 one complex Coulomb modulus.
\item \emph{Flavour symmetry.}
 $G_{\mathrm{flav}} = \mathfrak{su}(2)^{24}$, one $\mathfrak{su}(2)$
 per puncture; the maximal torus is $24$-dimensional, matching the
 number of $I_1$ singular fibres of the K3 elliptic fibration.
\item \emph{Steiner-system rigidity.}
 The $24$ punctures of $\Sigma_{0,24}$ admit the structure of a
 Steiner system $S(5,8,24)$
 \textup{(}the unique~$5$-design on $24$ points with block-size $8$,
 underlying the Mathieu group $M_{24}$\textup{)}. Its $759$ blocks
 give a candidate $M_{24}$-equivariant indexing set for free-field
 reductions of the chiral subalgebra. Identifying those blocks with
 Coulomb-branch flat directions or with Mukai-lattice subdata that
 survive the Borcherds lift is an additional K3 comparison problem,
 not a consequence of the Vol~I bar-cobar theorem.
\end{enumerate}
The class-$\mathcal S$ datum is therefore not identified directly with
the M5 boundary VOA
$\mathcal{A}_{\mathrm{M5}}(2)=W_{1+\infty}[\lambda{=}2]$ of
Theorem~\ref{thm:frontier-m5}: the stress central charges differ
($-214$ on the class-$\mathcal S$ parent, $25$ on the
$W_{1+\infty}[\lambda=2]$ stress lane). Any K3-coupled comparison must
pass through an interface, BRST, or protected-index map which changes
the stress normalisation; it is not an isomorphism of chiral algebras
at the level of central charge. The $4\mathrm d$ Coulomb dimension
$21$ and the $2\mathrm d$ central charge $-214$ remain the
Shapere--Tachikawa / Chacaltana--Distler invariants of the proposed
parent.
If the Steiner-system comparison is constructed, it fixes the
configuration as \emph{combinatorially unique}: among genus-zero
$A_1$ theories with $24$ punctures, the $S(5,8,24)$-symmetric
arrangement is the candidate coupling to the K3 boundary M5 datum
compatible with the $M_{24}$-equivariant Borcherds lift.
This is the proposed four-dimensional $\mathcal N{=}2$ parent of the
M5 boundary datum at $N=2$ on the K3-coupled branch.
\end{remark}

\begin{remark}[$\Delta_5$ as 1-loop output of twisted 11D~SUGRA on $K3 \times T^2$ with 24 M5 on $I_1$ Kodaira fibres]
\label{rem:frontier-Delta5-onel-loop-11d-K3T2}
\ClaimStatusConjectured
\index{eleven-dimensional supergravity!K3 times T2}
\index{1-loop anomaly!Delta_5 weight}
\index{M5 brane!I_1 Kodaira fibre wrap}
\index{shadow weight!Delta_5 = 2 + 3}
The fifth-shadow obstruction class on the M5 boundary datum on the
K3-coupled branch has automorphic avatar~$\Delta_5$ and is
conjecturally compared with the one-loop output of twisted
eleven-dimensional supergravity on the eleven-manifold
\begin{equation}\label{eq:frontier-Delta5-eleven-manifold}
\mathcal Y_{11}
\;=\;
\mathbb R^{4|2} \times K3 \times T^2,
\end{equation}
\textup{(}twisted along the
$\mathbb R^{4|2}$~direction in the sense of
Costello--Li--Saberi~\cite{CostelloP2201,CPS2208}\textup{)}
with $24$ M5~branes wrapping the $24$ smooth fibres of the elliptic
fibration $K3 \to \mathbb{P}^1$ that degenerate at the $24$~$I_1$
Kodaira loci of the
discriminant locus of the elliptic fibration
\textup{(}see Bittleston--Costello--Zeng~\cite{BittlestonCostelloZeng24}
for the twisted~$11$D supergravity formalism;
Bittleston--Costello~\cite{BittlestonCostello25} for the
eleven-dimensional uplift of the M5 boundary algebra\textup{)}.
The shadow-weight decomposition
\begin{equation}\label{eq:frontier-Delta5-weight-decomposition}
\boxed{\;
\operatorname{wt}(\Delta_5) \;=\; 5 \;=\; 2 \;+\; 3,\;}
\end{equation}
visible at the $1$-loop level of the twisted SUGRA computation,
splits into a contribution of weight~$2$ from the
quadratic-fluctuation determinant of the bulk supergravity multiplet
on $K3 \times T^2$
(the genus-$1$ Reidemeister-torsion contribution of the elliptic
factor, weight~$2$) and a contribution of weight~$3$ from the
$24$-fold M5 boundary insertion at the $I_1$ Kodaira nodes
(the cubic-coupling vertex of the
M5~boundary chiral algebra at $\lambda = N$, weight~$3$).
The total weight $5 = 2 + 3$ is the candidate fifth-shadow lane of the
K3-coupled M5 datum, with the $2 + 3$ decomposition expected to mirror
the bulk/boundary split of the twisted
eleven-dimensional theory: the bulk sector contributes the
$\Delta_2$ anomaly of the supergravity determinant; the $24$
M5~boundaries contribute the $\Delta_3$ piece of the $r$-matrix
shadow at depth~$3$. Costello--Paquette--Sharma's one-loop
twist-cancellation theorem and the single-fibre computation of
Bittleston--Costello supply the comparison surface. Restoring the K3
boundary geometry, assembling the $24$ fibres, and identifying the
result with~\eqref{eq:frontier-Delta5-weight-decomposition} are the
remaining proof obligations. The resulting $11$D-SUGRA scalar target is
the K3-enlarged Theorem~C bucket of
Remark~\ref{rem:frontier-K3-bucket-enlargement}.
\end{remark}

\begin{remark}[The holographic datum table]
\label{rem:frontier-holographic-datum-table}
\index{holographic datum!standard examples}
The following table separates standard boundary computations from
frontier physical imports. The $\kappa$ entries are Vol~I boundary
scalars on the stated lane; physical holography requires the typed
comparison maps described above.
\begin{center}
\renewcommand{\arraystretch}{1.25}
\small
\begin{tabular}{@{}l@{\;\;}c@{\;\;}c@{\;\;}c@{\;\;}l@{}}
\textbf{System} & $\cA$ & $\kappa(\cA)$
 & \textbf{Class} & \textbf{Key feature} \\
\midrule
Pure gravity
 & $\mathrm{Vir}_c$
 & $c/2$ & M
 & Infinite shadow depth \\
CS + gravity
 & $\widehat{\mathfrak{g}}_k$
 & $\dim\mathfrak{g}(k{+}h^\vee)/(2h^\vee)$
 & L
 & Finite depth, Verlinde \\
Twisted $\mathcal{N}{=}4$
 & $\mathfrak{gl}_N$ at $k{=}1$
 & $(N^2{-}1)(N{+}1)/(2N)+1$
 & L
 & CYBE R-matrix \\
ABJM (M2)
 & $\cA_{\mathrm{ABJM}}(N,k)$
 & $-(N^2{+}1)$ \textup{(pre-BRST)}
 & M
 & Airy, $N^{3/2}$ scaling \\
M5
 & $W_{1+\infty}[\lambda{=}N]$
 & $(4N^3{-}3N{-}1)/2$
 & M
 & $N^3$ scaling \\
Celestial gauge
 & $J_{\mathrm{col}}$
 & $0$ (self-dual)
 & G
 & $\mathrm{Com}^! = \mathrm{Lie}$ \\
Celestial gravity
 & $\mathfrak{w}_{1+\infty}$
 & $c/2$ (T-line)
 & M
 & Soft graviton tower
\end{tabular}
\end{center}
Each boundary datum is determined by the boundary VOA~$\cA$ and the
universal MC element~$\Theta_\cA$ on the strict modular Koszul locus.
The physical rows add comparison functors, completions, and large-$N$
or K3 imports.
\end{remark}

\begin{computation}[Burns space holographic modular Koszul datum;
\ClaimStatusProvedHere]

\label{comp:burns-space-holographic-datum}
\index{Burns space!holographic datum|textbf}
\index{holographic modular Koszul datum!Burns space}
\index{self-dual gravity!Burns space}
Self-dual gravity on Burns space~\cite{CPS2208,CPS2306} has boundary
VOA $\cA_{\mathrm{Burns}} = \beta\gamma^{\otimes 4}_{\lambda=1}$
with $\mathrm{SO}(8)$ flavor symmetry.
The holographic modular Koszul datum is:
\begin{enumerate}[label=\textup{(\roman*)}]
\item $\cA = \beta\gamma^{\otimes 4}_{\lambda=1}$
 \textup{(}four $\beta\gamma$ pairs at conformal
 weight~$\lambda = 1$\textup{)}.
 Central charge
 $c(\cA_{\mathrm{Burns}}) = 4 \cdot 2 = 8$.
\item $\cA^! = bc^{\otimes 4}_{\lambda=1}$
 \textup{(}four $bc$ ghost pairs\textup{)},
 with $\kappa(\cA^!) = -4$.
 Complementarity:
 $\kappa(\cA) + \kappa(\cA^!) = 0$
 \textup{(}free-field anti-symmetry\textup{)}.
\item $\kappa(\cA_{\mathrm{Burns}}) = 4$
 \textup{(}Proposition~\textup{\ref{prop:independent-sum-factorization}}:
 each $\beta\gamma$ pair contributes
 $\kappa(\beta\gamma_1) = 1$; four copies give
 $\kappa = 4$\textup{)}.
\item Shadow class: $\mathsf{C}$ per $\beta\gamma$ pair
 \textup{(}contact, $r_{\max} = 4$\textup{)};
 class~$\mathsf{M}$ on the $T$-line
 \textup{(}the Virasoro subalgebra at $c = 8$ has infinite
 shadow depth\textup{)}.
\item Genus tower \textup{(}Theorem~D,
 uniform-weight lane: all generators at
 $\lambda = 1$\textup{)}:
 \begin{align}
 F_1(\cA_{\mathrm{Burns}})
 &= \kappa \cdot \lambda_1^{\mathrm{FP}}
 = 4 \cdot \tfrac{1}{24}
 = \tfrac{1}{6},
 \label{eq:burns-F1} \\
 F_2(\cA_{\mathrm{Burns}})
 &= \kappa \cdot \lambda_2^{\mathrm{FP}}
 = 4 \cdot \tfrac{7}{5760}
 = \tfrac{7}{1440},
 \label{eq:burns-F2} \\
 F_3(\cA_{\mathrm{Burns}})
 &= \kappa \cdot \lambda_3^{\mathrm{FP}}
 = 4 \cdot \tfrac{31}{967680}
 = \tfrac{31}{241920}.
 \label{eq:burns-F3}
 \end{align}
 The values $F_2$ and~$F_3$ are predictions beyond
 Costello--Paquette--Sharma~\cite{CPS2306}, who compute
 tree-level and one-loop amplitudes explicitly.
\item $T$-line planted-forest correction at genus~$2$:
 $\delta_{\mathrm{pf}}^{(2,0)} = S_3(10S_3 - \kappa)/48
 = 2(20 - 4)/48 = 2/3$
 \textup{(}where $S_3 = 2$ is the universal Virasoro
 cubic shadow\textup{)}.
 This modifies the tautological decomposition on the
 $T$-line but does not affect the scalar genus expansion
 $F_g = \kappa \cdot \lambda_g^{\mathrm{FP}}$ on the
 uniform-weight lane
 \textup{(}Theorem~D\textup{)}.
\end{enumerate}
Computation:
\texttt{compute/lib/theorem\_burns\_f2\_engine.py}
\textup{(}six independent paths for~$F_2$\textup{)}.
\end{computation}

\begin{remark}[From $\kappa$ to the Witten genus: height interpolation
via the shadow tower]
\label{rem:kappa-hodge-to-witten-genus}
\ClaimStatusConjectured
\index{kappa@$\kappa$!chromatic height interpretation}
\index{Witten genus!shadow tower interpolation}
\index{chromatic height!shadow tower}
\index{TMF!shadow tower connection}
The characteristic $\kappa(\cA)$ of Theorem~D is a Hodge-theoretic
invariant: it governs the class
$\mathrm{obs}_g = \kappa(\cA) \cdot \lambda_g$
\textup{(}UNIFORM-WEIGHT\textup{}) on
$\overline{\cM}_g$ and determines the
genus-$1$ curvature $d_{\mathrm{fib}}^2 = \kappa \cdot \omega_1$
of the fiberwise bar differential. In the chromatic hierarchy of
stable homotopy theory
(Ravenel~\cite{Ravenel86}), this places $\kappa$ at height~$0$:
it is a rational invariant, extracted from the ordinary Hodge bundle
$\bE \to \overline{\cM}_g$ with no higher-chromatic content.

The shadow tower $\Theta_\cA \in \MC(\fg^{\mathrm{mod}}_\cA)$
encodes genus-by-genus obstruction data beyond $\kappa$.
At genus~$1$, the curvature $\kappa \cdot \omega_1$ governs the
$E_2^*$-quasi-modular generating function; at genus~$2$,
the three-shell structure introduces corrections
$\delta F_2^{\mathrm{cross}}$ (all-weight) that depend on the full
cubic shadow~$S_3$. These higher-genus shells carry information
invisible to the rational projection. The chromatic-shadow
correspondence
(Remark~\ref{rem:chromatic-shadow-correspondence})
identifies the genus-$n$ shell of the shadow tower with the
$E(n)$-localization in stable homotopy. Under this identification:
\begin{enumerate}[label=\textup{(\roman*)}]
\item Height $0$ (rational):
 $\kappa(\cA)$, the degree-$2$ scalar shadow: in abelian and scalar
 families $\kappa(\cA)=\mathrm{av}(r(z))$, while for non-abelian
 affine Kac--Moody
 $\kappa(V_k(\fg))=\mathrm{av}(r(z))+\dim(\fg)/2$; it controls
 $F_g$ on the uniform-weight lane.
\item Height $1$ ($K$-theory):
 the genus-$1$ shadow, encoding the quasi-modular generating
 function $\sum_g F_g q^g$ as a weak Jacobi form, lifts to an
 element of $K$-theory of $\cM_{1,1}$ via the Chern character.
\item Height $2$ (TMF):
 the full shadow tower, assembled over the universal elliptic curve
 $\cE \to \cM_{1,1}$, should define a class in topological modular
 forms. The elliptic bar propagator $d\log E(z,w)$ on $E_\tau$
 (weight~$1$) varies in the formal group $\hat{E}_\tau$ of the fiber;
 Conjecture~\ref{conj:chromatic-bar-cobar} predicts that the
 resulting factorization coalgebra carries a TMF-module structure
 (Ando--Hopkins--Strickland~\cite{AHS01}).
\end{enumerate}
The Witten genus of a string manifold~$M$
is a TMF-valued invariant that refines the
$\hat{A}$-genus (height~$0$) through the elliptic genus
(height~$1$) to TMF (height~$2$).
The shadow tower provides a parallel refinement on the chiral
algebra side: $\kappa$ is the $\hat{A}$-genus analogue (a single
rational number), the genus-$1$ quasi-modular tower is the
elliptic genus analogue, and the full MC element~$\Theta_\cA$ is
the conjectural TMF-valued invariant.
If both identifications hold, the shadow tower
of~$\cA$ and the Witten genus of the target manifold of
the associated sigma model should agree after applying the
string orientation $\mathrm{MString} \to \mathrm{tmf}$
to the factorization coalgebra of~$\cA$.
This is a refinement of Conjecture~\ref{conj:chromatic-bar-cobar}(c);
the height-$2$ comparison remains open.
\end{remark}

\begin{remark}[The shadow tower as Taylor expansion of a fully extended TQFT]

\label{rem:shadow-tower-tqft-taylor}
\ClaimStatusHeuristic
\index{shadow tower!TQFT interpretation}
\index{TQFT!shadow tower as Taylor expansion}
\index{Koszul conductor!cobordism invariant}
\index{G/L/C/M classification!cobordism type}
The MC element $\Theta_\cA \in \MC(\fg^{\mathrm{mod}}_\cA)$
encodes the genus-by-genus obstruction data of~$\cA$ as a formal
power series in the genus variable. This series admits a conjectural
interpretation as the Taylor expansion of a fully extended
topological field theory $Z_\cA$ in the sense of
Lurie~\cite{Lurie09TFT}, where the cobordism category is
determined by the shadow depth class of~$\cA$.

The $G/L/C/M$ classification
(Definition~\ref{def:shadow-depth-classification})
controls the cobordism type:
\begin{enumerate}[label=\textup{(\roman*)}]
\item Class~$G$ (shadow depth $r = 2$; Heisenberg):
 the shadow tower terminates at degree~$2$.
 The associated field theory would be finite;
 all partition functions are determined by the single
 invariant $\kappa(\cA)$.
\item Class~$L$ (shadow depth $r = 3$; affine Kac--Moody):
 the tower terminates at degree~$3$.
 The corresponding point-value would require affine cobordism data,
 with the cubic shadow~$S_3$ governing the first
 nontrivial sewing relation.
\item Class~$C$ (shadow depth $r = 4$; $\beta\gamma$ systems):
 the tower terminates at degree~$4$.
 The corresponding point-value would require quadratic cobordism data,
 with the quartic obstruction $o_4 = 0$ but
 $\mathfrak{Q} \neq 0$.
\item Class~$M$ (shadow depth $r = \infty$; Virasoro, $W_N$):
 the tower does not terminate.
 The conjectural field theory is infinite-depth: infinitely many
 cobordism relations are needed, and the partition function
 at each genus carries independent information.
\end{enumerate}

The scalar complementarity invariant
$\varkappa(\cA) := \kappa(\cA) + \kappa(\cA^!)$
is an invariant of this cobordism structure.
Under the duality $\cA \leftrightarrow \cA^!$,
the scalar invariant is preserved:
$\varkappa(\cA^!) = \kappa(\cA^!) + \kappa((\cA^!)^!) = \kappa(\cA^!) + \kappa(\cA) = \varkappa(\cA)$.
The value of~$\varkappa$ is family-specific:
$\varkappa = 0$ for affine Kac--Moody, Heisenberg, and free-field algebras;
$\varkappa = 13$ for the Virasoro algebra;
$\varkappa = 250/3$ for $W_3$;
$\varkappa = 98/3$ for Bershadsky--Polyakov
(Proposition~\ref{prop:bp-complementarity-constant}).
The scalar $\varkappa$ is distinct from the Trinity ghost-charge
conductor $K(\cA) := c(\cA) + c(\cA^!) = -c_{\mathrm{ghost}}(\BRST)$,
which takes per-family values $K(\Vir) = 26$, $K(W_3) = 100$,
$K(\mathrm{BP}) = 196$, $K(V_k(\fg)) = 2\dim\fg$, $K(\cW_N) = 4N^3 - 2N - 2$;
the two are related by $\varkappa = \varrho_{\cA} \cdot K$ with $\varrho$
family-specific ($\varrho_{\cW_N} = H_N - 1$,
$\varrho_{\mathrm{KM}} = \varrho_{\mathrm{free}} = 0$,
$\varrho_{\mathrm{BP}} = 1/6$).

At the self-dual value $\kappa(\cA) = K(\cA)/2$,
the scalar shadow tower acquires a symmetry
$\kappa(\cA) = \kappa(\cA^!)$
that can kill odd-parity obstruction classes only under an additional
hypothesis: the Verdier involution must lift to a parity action on the
completed point-value, and the parity splitting must be compatible
with sewing. Under that hypothesis the resulting field theory is
Euler-type:
the partition function $Z_\cA(\Sigma_g)$ is multiplicative
under cutting and gluing along separating curves,
because the odd part of the cross-channel correction
$\delta F_g^{\mathrm{cross}}$ is projected out. Scalar self-duality
alone does not eliminate the full all-weight tower.
For the Virasoro algebra this self-dual point is
$c = 13$, where $\kappa(\mathrm{Vir}_{13}) = 13/2$.

No construction of~$Z_\cA$ as a fully extended functor has been
carried out; the interpretation is heuristic.
The finite-depth cases ($G$, $L$, $C$) are the most accessible,
since the finitely many shadow relations directly supply the
sewing data of a field theory on a finite set of generators.

The cobordism hypothesis~\cite{Lurie09TFT} would sharpen this picture
only after a fully dualizable point-value carrying the completed shadow
tower has been constructed.
A fully extended $n$-dimensional topological field theory is
determined by its value on the point: the $(\infty,n)$-category
assigned to a $0$-manifold recovers, by iterated adjunction, the
entire functor on all cobordisms.
Under such a construction, the shadow tower of~$\cA$ would supply the
Taylor coefficients of this point-value: $\kappa(\cA)$ is the linear term (genus-$0$
two-point data), $S_3$ the quadratic correction (genus-$0$
three-point sewing), $S_4$ and the quartic obstruction~$\fQ$ the
cubic and quartic corrections, and $\Theta_\cA$ the completed
formal series.
The $G/L/C/M$ classification records how many of these algebraic
shadow coefficients are independent: class~$G$ means the linear
term determines everything (a $1$-truncated point-value); class~$L$
requires quadratic data; class~$C$ requires cubic and quartic;
class~$M$ requires the full formal series, and the point-value is
a genuinely infinite object.
From this vantage the scalar complementarity $\varkappa(\cA)$ is a
candidate obstruction to the point-value being self-conjugate under the
Verdier involution $\cA \leftrightarrow \cA^!$. Its vanishing
($\varkappa = 0$ for affine Kac--Moody and free-field families) is the
scalar condition for orientable self-duality; $\varkappa \neq 0$
(Virasoro at $\varkappa = 13$, $W_3$ at $\varkappa = 250/3$,
Bershadsky--Polyakov at $\varkappa = 98/3$) measures scalar failure of
that condition. It is not yet a theorem about orientation reversal in a
constructed cobordism functor. If such a symmetric monoidal functor is
constructed, then the cobordism hypothesis says it is determined by
the algebraic structure of the point-value; the open obligation is to
construct that point-value from the completed shadow tower.
\end{remark}

\begin{remark}[K3 boundary data and the Theorem C bucket $\{0,\;8,\;13,\;250/3,\;98/3\}$]
\label{rem:frontier-K3-bucket-enlargement}
\ClaimStatusConjectured
\index{Theorem C!bucket enlargement K3}
\index{K3!complementarity bucket}
\index{Mukai lattice!chiral Koszul characteristic}
\index{Borcherds lift!bucket}
The scalar complementarity invariant
$\varkappa(\cA) = \kappa(\cA) + \kappa(\cA^!)$
inventoried above (with values $\{0,\;13,\;250/3,\;98/3\}$ on the
Heisenberg / affine Kac--Moody / Virasoro / $\mathcal{W}_3$ /
Bershadsky--Polyakov core of Part~\ref{part:standard-landscape}) acquires
a fifth natural value, $\varkappa = 8$, once one couples the Vol~I
bar-cobar machinery to the $K3$ boundary data that anchors Volume~III.
On a polarized $K3$ surface~$S$ with Mukai lattice
$\widetilde{H}(S,\mathbb Z)
= H^0 \oplus H^2 \oplus H^4$
of signature $(4,20)$, the positive-norm sub-lattice is
$4$-dimensional; its four positive directions supply
$c_+(\mathrm{Mukai}(K3)) = 4$, so the combinatorial ceiling
$K^{\kappa_{\mathrm{ch}}} = 2\,c_+(\mathrm{Mukai}(K3)) = 8$
is a cohomological invariant of the target, independent of any
chiral-algebra presentation
\textup{(}Mukai's pairing,
\cite[Lemma~10.6]{Mukai81}\textup{)}. Three independent faces
converge on the same~$8$:
\begin{enumerate}[label=\textup{(\roman*)}]
\item \emph{Mukai-lattice face.}
 $c_+(\widetilde{H}(K3)) = 4$ and the Koszul conductor reads
 $K^{\kappa_{\mathrm{ch}}} = 2 c_+ = 8$ on the B-family Mukai
 scalar lane. This is not the Euler number
 $\chi(K3)=24$ and no factorization $24=8+16$ is used here.
\item \emph{Heegner/Chern face.}
 Borcherds' reciprocity
 \textup{(\cite{Borcherds1995}, see also the chiral packaging in
 Chapter~\ref{ch:chiral-climax-platonic})} identifies the
 modular Chern class carried by a Heegner divisor on the period
 domain~$\Omega_{K3}$ with the regularized theta-integral of the
 bar coalgebra in question; the discriminant-$8$ combination is
 the $O(2,\rho)$ trace of the positive-definite quadratic form on the
 $K3$ lattice, the same $8$ that appears on the Mukai face.
\item \emph{Lift face.}
 The Borcherds lift of the chiral partition function of a
 free-field, boundary-supported chiral algebra at
 $c_{2\mathrm d} = 8$ (e.g.~$E_8$ at level one, whose lattice is
 positive-definite, rank~$8$, even and unimodular) is the
 celebrated~$\Phi_{12}$-type automorphic product on the $K3$
 moduli space; $\kappa_\ch = 8$ is the unique scalar value at
 which the chiral partition function is a full symmetric
 automorphic form without pole.
\end{enumerate}
The three faces unify, conditionally, under Bruinier's Heegner Chern-class
reciprocity \textup{(}regularized theta-integrals $=$ divisor-class
pullbacks on modular towers\textup{)} to the single B-family Mukai
scalar $\varkappa = 8$. This scalar must be kept separate from the
Vol~III $K3\times E$ constants:
$\kappa_{\mathrm{BKM}}(\Delta_5)=5$,
$2\kappa_{\mathrm{BKM}}(\Delta_5)=10$, the Heisenberg-level scalar
$3$, the compact total-space scalar $0$, and the fibre Euler scalar
$24$. The conditional Theorem~C bucket on the extended B-family is
therefore
\begin{equation}\label{eq:frontier-K3-bucket}
\varkappa(\cA)
\;\in\;
\bigl\{0,\;8,\;13,\;250/3,\;98/3\bigr\}
\qquad
\textup{(}\text{on the extended B-family including K3 boundary data}\textup{)},
\end{equation}
with $\varkappa = 0$ the Heisenberg/affine/free-field anti-symmetry,
$\varkappa = 8$ the $K3$ Mukai-lattice value,
$\varkappa = 13$ the Virasoro self-dual point,
$\varkappa = 250/3$ the $\mathcal{W}_3$ principal value, and
$\varkappa = 98/3$ the Bershadsky--Polyakov value. The fact that
$8 = c_+\cdot 2$ uses the same positive-definite projection that
gives $c_{\mathrm{Monster}} = 24$ on the Leech lattice
\textup{(\cite{Borcherds92})}; $K3$ is the self-dual geometric object
that sits between Heisenberg ($\varkappa = 0$) and Virasoro
($\varkappa = 13$) on this scalar lane. The automorphic product
statements are proved on the Borcherds side
\textup{(}\cite{Borcherds1995}, see~\cite[Ch.~4]{Borcherds92}\textup{)};
their pushforward to a Vol~I bar-cobar theorem is the conjectural
comparison map, not an already proved scalar identity.
\end{remark}

\begin{remark}[Universal $\Phi$-trace identity bridging Vol~I conductor and Vol~III BKM~$\kappa$]
\label{rem:frontier-universal-phi-trace}
\ClaimStatusConjectured
\index{Phi-trace identity@$\Phi$-trace identity}
\index{Koszul conductor!Borcherds lift identification}
\index{BKM algebra!modular characteristic via $\Phi$}
\index{universal trace identity!chiral to automorphic}
The Koszul conductor
$K(\cA) = c(\cA) + c(\cA^!) = -c_{\mathrm{ghost}}(\BRST(\cA))$
of Remark~\ref{rem:shadow-tower-tqft-taylor} is the BRST
supertrace of the bar--cobar ghost system; on the Vol~III side a
Lorentzian lattice BKM algebra~$\mathfrak g_\Lambda$ carries a fixed
Borcherds automorphic product $\Psi_\Lambda$, and its BKM
characteristic is indexed by that product:
$\kappa_{\mathrm{BKM}}(\Psi_\Lambda) = c_\Lambda(0)/2$, where
$c_\Lambda(0)$ is the constant coefficient of the weight-$(1-\rho/2)$
vector-valued modular form that defines the Borcherds
product~\cite{Borcherds1995,Borcherds92}. Let
$\Phi\colon
\mathsf{CY}^{\mathrm{sm}}
\to
\mathsf{ChirAlg}^{\mathrm{mod}}$
denote the Vol~III Calabi--Yau-to-chiral functor that attaches to a
smooth Calabi--Yau~$Y$ its boundary chiral algebra, extended by
nearby-cycles to the Kodaira-degenerate locus. On a
Borcherds-lift-exact comparison surface, where the chiral conductor
and the automorphic product are attached to the same lattice object,
the \emph{universal $\Phi$-trace identity} predicts
\begin{equation}\label{eq:frontier-phi-trace}
\boxed{\;
K(\Phi(Y))
\;=\;
-c_{\mathrm{ghost}}(\BRST(\Phi(Y)))
\;=\;
\operatorname{tr}_{\cZ^{\mathrm{der}}_{\mathrm{ch}}(\Phi(Y))}
 \!\bigl(K_{\Phi(Y)}\bigr)
\;\xrightarrow{\;\Phi_*\;}\;
2\,\kappa_{\mathrm{BKM}}(\Psi_\Lambda)
\;=\;
c_\Lambda(0),
\;}
\end{equation}
where the left-hand equality is the definition of the chiral
conductor as a derived-centre supertrace when the conductor operator
is defined on the completed chiral Hochschild complex
\textup{(}Theorem~\ref{thm:sphere-reconstruction}, via
$\cZ^{\mathrm{der}}_{\mathrm{ch}}(\cA)
\simeq \mathcal C^\bullet_{\mathrm{ch}}(\cA,\cA)$\textup{)}, the
middle arrow is the desired Borcherds pushforward~$\Phi_*$, and the
right-hand equality is the BKM normalization of the automorphic
product. The derived centre here is chiral Hochschild. It is not a
categorical THH object or a physical bulk unless the comparison maps
of the completed holographic datum have been constructed. Concretely,
on the paramodular
$K3\times E$ row $\Psi_\Lambda=\Delta_5$ and the automorphic side is
$2\,\kappa_{\mathrm{BKM}}(\Delta_5)=c_{\Delta_5}(0)=10$; the
Mukai-doubling conductor $K^{\kappa_{\mathrm{ch}}}=8$ of
Remark~\ref{rem:frontier-K3-bucket-enlargement} is a separate
$\mathcal B$-family scalar, not a Borcherds weight of $\Delta_5$.
For $\Lambda = \text{II}_{1,1}$ the identity reduces to the Heisenberg
anti-symmetry $K = 0$; for $\Lambda = $ Leech it specializes to the
Monstrous $c_\Lambda(0) = 24$
\textup{(}Borcherds~\cite{Borcherds92}\textup{)}.
The diagram intertwining chiral and automorphic supertraces is
\begin{center}
\small
\renewcommand{\arraystretch}{1.2}
\begin{tabular}{c@{\;\;\xrightarrow{\;\Phi\;}\;\;}c}
$\bigl(\cA,\;\cA^!,\;\BRST(\cA)\bigr)$
 & $\bigl(\mathfrak g_\Lambda,\;\mathfrak g_\Lambda^-,\;
 \operatorname{Jacobi-form lift}\bigr)$ \\[2pt]
$K(\cA) = c + c^!$
 & $c_\Lambda(0)$ \\[2pt]
$\cZ^{\mathrm{der}}_{\mathrm{ch}}(\cA)$
 & $\operatorname{gr}^W_\bullet \mathfrak g_\Lambda$
 (weight filtration) \\[2pt]
chiral Hochschild
 & automorphic zero-Fourier \\
supertrace $\operatorname{tr} K_\cA$
 & coefficient $c_\Lambda(0)/2$
\end{tabular}
\end{center}
The conditional identity~\eqref{eq:frontier-phi-trace} sharpens the nonuple
holographic datum~$\Pi^{\mathrm{oc}}_X(\cA)$ on the K3-coupled branch:
on the Borcherds-lift-exact surface, the complementarity scalar
$\kappa(\cA)+\kappa(\cA^!)$ is compared with the zero-Fourier
coefficient of the automorphic lift, and the Vol~I bar-cobar
supertrace calculus pushes forward across~$\Phi$ to the Vol~III BKM
side when the map~$\Phi_*$ exists. The automorphic equality is proved
on the Borcherds-lift-exact locus (Lorentzian Heegner lattices of
signature $(2,n)$). The chiral pushforward equality remains a proof
obligation unless the Vol~I to Vol~III comparison functor is
constructed.
\end{remark}

\section{Synthesis}
\label{sec:fmhp-supplement}

\begin{remark}[Holographic datum at the K3 base]
\label{rem:fmhp-1}
The platonic holographic programme
$\mathcal H(T) = (A, A^i, A^!, C, r(z), \Theta_A,
\nabla^{\mathrm{hol}})$ has the following conditional K3 realization
at the bialgebra $\mathbf H_{\Delta_5}$:
\[
\begin{array}{rcl}
A & = & \mathbf H_{\Delta_5}, \\
A^i & = & H^*(B_{\mathrm{ch}}(\mathbf H_{\Delta_5})), \\
A^! & = & \mathbb D_{\Ran}A^i
 \simeq \mathbf H_{\Delta_5}\ \text{only after the self-dual
 strictification at }\hbar^2=-1/8, \\
C & = & \cZ^{\mathrm{der}}_{\mathrm{ch}}(\mathbf H_{\Delta_5})
 \ \text{or, on the line side, the chosen module-category trace}, \\
\Omega_X(B_{\mathrm{ch}}(\mathbf H_{\Delta_5})) & \simeq
 & \mathbf H_{\Delta_5}
\ \text{by bar--cobar inversion, not by Verdier duality}, \\
r(z) & = & R_{\mathrm{Sieg,dyn}}(z;\tau,w)\bigm|_{\hbar\to 0}, \\
\Theta_A & = & \text{universal obstruction in the Hall ordered dGLA}, \\
\nabla^{\mathrm{hol}} & = & \text{KZB connection on $\bar{\mathcal A}_2$ pulled back via Kuga--Satake}.
\end{array}
\]
The Vol~I recovery statements reconstruct the ordered bar complex, the
typed Verdier/Koszul branch, the collision residue, and the universal
MC element on the strict modular Koszul boundary locus. The assertion
that these four recoveries describe the K3 base point is conditional on
the Vol~III K3 chiral-bialgebra construction, the Lusztig
specialization $\hbar^2=-1/8$, the Kuga--Satake comparison, and the
completion needed for the line-category trace. Cross-ref Vol.~III,
Chapter~\ref{V3-ch:k3-chiral-bialgebra-platonic}.
\end{remark}

\begin{remark}[Global holographic triangle on the B-family]
\label{rem:fmhp-2}
The global holographic triangle -- boundary $\leftrightarrow$ bulk
$\leftrightarrow$ line -- closes on the Lorentzian-lattice-parametric
B-family only after the parity-halving and cusp/Euler comparison maps
are constructed on the same object. The automorphic input supplies the
Bruinier 2002 Prop.~5.1 Heegner
Chern-class reciprocity: the $\mathbb Z/8$-class $[L^{\Delta_5}]\in\mathrm{CH}^1(H_1)$
as a candidate boundary modular holonomy (Mukai doubling), bulk
monodromy (Humbert $H_1$), and line fusion determinant
(Lusztig $\ell = 8$). At the Lusztig specialisation $\hbar^2 = -1/8$
the boundary chiral Hochschild sector is proved by Theorem~H to
concentrate in degrees $\{0,1,2\}$ on the Vol~I locus. This is not a
physical bulk theorem. The global boundary/bulk/line triangle remains
conjectural beyond the typed comparison maps. The $\kappa$~indexing on
the two sides of
the triangle is compared by
Proposition~\ref{prop:kappa-bkm-cross-volume-reconciliation} at
$\kappa_{\mathrm{crown}}^{(3\text{-faces})}
 = 2\kappa_{\mathrm{BKM}}^{(\mathrm{Borch})} + \chi(\cO_X) + 2$,
evaluating to $12 = 2\cdot 5 + 0 + 2$ on $K3 \times E$. This is the
Volume~I crown scalar, not the Borcherds weight
$\kappa_{\mathrm{BKM}}(\Delta_5)=5$ and not the
$K3\times E$ four-scalar spectrum $\{0,3,5,24\}$. The $+2$ is the
bar--cobar double-parity correction and is the quantitative signature
expected, at the Lusztig specialisation, to make the two indexings
bulk-boundary projections of the same Drinfeld double
$U_\cA = \cA \bowtie \cA^{!}_\infty$. The class~$\mathsf{M}$ residue
left open by the boundary-linear argument closes cohomologically only
on the DS--Hochschild comparison lane, and at chain level only in the
completed/pro sense, via the Vol.~II DS--Hochschild bridge
(Corollary~\ref{V2-thm:class-M-chain-bulk}); the residue beyond the
DS~lane remains conjectural and is exhibited by the admissible bar
class $\eta_4$ of $L_{-1/2}(\mathfrak{sl}_2)$. See Vol.~III,
Chapter~\ref{V3-ch:k3-chiral-bialgebra-platonic},
\S\ref{sec:four-recovery-theorems-at-supplement}.
\end{remark}
